\documentclass[10pt,xcolor=dvipsnames,aspectratio=169,spanish]{beamer}

\usetheme[progressbar=frametitle]{metropolis}
\usepackage{appendixnumberbeamer}

\usepackage{booktabs}
\usepackage[scale=2]{ccicons}
\usepackage{adjustbox}
\usepackage{fontspec}
\usepackage{ wasysym }

\usepackage{pgfplots}
\usepgfplotslibrary{dateplot}

\usepackage{xspace}
\newcommand{\themename}{\textbf{\textsc{metropolis}}\xspace}

\usepackage{pdfpages}
\pgfdeclareimage[width=\paperwidth]{mybackground}{img/portada-ugr-169.pdf}

\usepackage{babel}          % Automatic translations
\usepackage{amssymb}        % Mathematical symbols
\usepackage{mathtools}      % Mathematical symbols
\usepackage{soul}
\makeatletter
\let\HL\hl
\renewcommand\hl{%
  \let\set@color\beamerorig@set@color
  \let\reset@color\beamerorig@reset@color
  \HL}
\makeatother

\usepackage{framed}

\usepackage{listings} % Para la inclusión de trozos de código
\lstset{
     literate=%
         {á}{{\'a}}1
         {í}{{\'i}}1
         {é}{{\'e}}1
         {ý}{{\'y}}1
         {ú}{{\'u}}1
         {ó}{{\'o}}1
         {ě}{{\v{e}}}1
         {š}{{\v{s}}}1
         {č}{{\v{c}}}1
         {ř}{{\v{r}}}1
         {ž}{{\v{z}}}1
         {ď}{{\v{d}}}1
         {ť}{{\v{t}}}1
         {ň}{{\v{n}}}1                
         {ů}{{\r{u}}}1
         {Á}{{\'A}}1
         {Í}{{\'I}}1
         {É}{{\'E}}1
         {Ý}{{\'Y}}1
         {Ú}{{\'U}}1
         {Ó}{{\'O}}1
         {Ě}{{\v{E}}}1
         {Š}{{\v{S}}}1
         {Č}{{\v{C}}}1
         {Ř}{{\v{R}}}1
         {Ž}{{\v{Z}}}1
         {Ď}{{\v{D}}}1
         {Ť}{{\v{T}}}1
         {Ň}{{\v{N}}}1                
         {Ů}{{\r{U}}}1    
}
\usepackage{lstautogobble}  % Fix relative indenting

\definecolor{bluekeywords}{rgb}{0.13, 0.13, 1}
\definecolor{greencomments}{rgb}{0, 0.5, 0}
\definecolor{redstrings}{rgb}{0.9, 0, 0}
\definecolor{graynumbers}{rgb}{0.5, 0.5, 0.5}
\def\req{\rotatebox{90}{$\scriptstyle=$}}


\lstset{
    autogobble,
    columns=fullflexible,
    showspaces=false,
    showtabs=false,
    breaklines=true,
    showstringspaces=false,
    breakatwhitespace=true,
    escapeinside={(*@}{@*)},
    commentstyle=\color{greencomments},
    keywordstyle=\color{bluekeywords},
    stringstyle=\color{redstrings},
    numberstyle=\color{graynumbers},
    %basicstyle=\ttfamily\footnotesize,
    frame=l,
    framesep=30pt,
    xleftmargin=30pt,
    tabsize=4,
    captionpos=b,
    basicstyle=\ttfamily,
    numbers=left,
    numberstyle=\footnotesize\color{graynumbers},
    numbersep=10pt,
    %escapechar=|,
    numberblanklines=false,
    columns=flexible
}



\setbeamercolor{section title}{fg=Maroon,bg=Maroon}
\setbeamercolor*{structure}{bg=Maroon!20,fg=Maroon}
\setbeamercolor*{palette primary}{use=structure,fg=white,bg=structure.fg}
\setbeamercolor{progress bar}{fg=gray, bg=gray}

\setbeamerfont{itemize/enumerate body}{}
\setbeamerfont{itemize/enumerate subbody}{size=\normalsize}
\setbeamerfont{itemize/enumerate subsubbody}{size=\normalsize}

%gets rid of bottom navigation bars
\setbeamertemplate{footline}[frame number]{}

%gets rid of bottom navigation symbols
\setbeamertemplate{navigation symbols}{}

%gets rid of footer
%will override 'frame number' instruction above
%comment out to revert to previous/default definitions
\setbeamertemplate{footline}{}
\setbeamercovered{transparent}

\newcommand{\hle}[1]{\hl{\emph{#1}}}
\newcommand\Vtextvisiblespace[1][.3em]{%
\mbox{\kern.06em\vrule height.3ex}%
\vbox{\hrule width#1}%
\hbox{\vrule height.3ex}}
\newcommand{\palabra}[1]{\texttt{\textquotesingle{}{#1}\textquotesingle{}}}

%%%%%%%%%%%%%%%%%%%%%%%%%%%%%%%%%%%%%%%%%%%%%%%%%%%%%%%%%%%%%%%%%%%%%%

\title{Una perspectiva algorítmica a la incompletitud de las Matemáticas}
\subtitle{Defensa TFG}
\date{\today}
\author{Miguel Ángel Fernández Gutiérrez}

%%%%%%%%%%%%%%%%%%%%%%%%%%%%%%%%%%%%%%%%%%%%%%%%%%%%%%%%%%%%%%%%%%%%%%

\setbeamertemplate{title page}{

        \begin{picture}(0,0)

            \put(-10mm,-45mm){%
                \pgfuseimage{mybackground}
            }

            \put(20mm,-5mm){%
                \begin{minipage}[b][45mm][t]{226mm}
                \hspace{-22mm}\includegraphics[width=60mm]{img/logo-ugr.pdf}
                
                    \usebeamerfont{title}{
{\large\textsc{trabajo de fin de grado}}

    \begin{flushleft}
      \LARGE
      Una perspectiva algorítmica a la \\incompletitud de las Matemáticas\\[18pt]
      \large
      Miguel Ángel Fernández Gutiérrez\\
      Tutor: Serafín Moral Callejón\\
      \normalsize
      \textrmlf{3 de julio, 2023}
    \end{flushleft}
                    
                    }
                \end{minipage}
            }

            \end{picture}

    }


\begin{document}

\maketitle

\begin{frame}{Agenda}
  \setbeamertemplate{section in toc}[sections numbered]
  \tableofcontents%[hideallsubsections]
\end{frame}

\section[Introducción]{Introducción}

\begin{frame}[fragile]{Introducción}
\begin{adjustbox}{max height={.75\textheight},center}
% !TeX root = ../libro.tex
% !TeX encoding = utf8
%
%*******************************************************
% Introducción
%*******************************************************

% \manualmark
% \markboth{\textsc{Introducción}}{\textsc{Introducción}} 
\chapter{Resumen}

%De acuerdo con la comisión de grado, el TFG debe incluir una introducción en la que se describan claramente los objetivos previstos inicialmente en la propuesta de TFG, indicando si han sido o no alcanzados, los antecedentes importantes para el desarrollo, los resultados obtenidos, en su caso y las principales fuentes consultadas.

%En este capítulo, resumiremos el proceso que hemos seguido hasta llegar a la demostración, desde una perspectiva algorítmica, del Primer Teorema de Incompletitud. En la \cref{fig:resumen}, se puede observar un esquema resumen de este trabajo, el cual es extendido en la \cref{fig:resumen-extendido}. En este resumen final se omiten las \cref{part:introduccion,,part:conclusion}.

Comenzamos este Trabajo de Fin de Grado con el \cref{ch:motivacion}, que destaca la motivación que subyace este trabajo, así como los objetivos que se pretenden cumplir. Inmediatamente después, en el \cref{ch:historia}, hacemos una breve contextualización histórica de esta obra.

El desarrollo teórico comienza con el \cref{ch:programas-python-maquinas-turing}, en cuya \cref{sec:programas-python} introducimos una clase específica de programas en Python, los \emph{programas SISO}, y observamos que esta clase es tan expresiva como la clase de todos los programas en Python. A continuación, en la \cref{sec:maquinas-turing}, profundizamos en las \emph{máquinas de Turing}, herramienta matemática principal de nuestro trabajo. En la \cref{sec:equivalencia}, introducimos nuevos modelos de computación equivalentes, hasta llegar al resultado principal del capítulo: que las máquinas de Turing pueden simular programas en Python, y viceversa; en otras palabras, que los programas en Python y las máquinas de Turing son equivalentes. Esto nos permitirá usar programas en Python convenientemente, sin realizar descripciones de máquinas de Turing. Concluimos introduciendo la tesis de Church-Turing en la \cref{sec:church-turing}, y viendo que la equivalencia probada puede ser un argumento a su favor.

En la \cref{sec:problemas-decidibles} definimos los \emph{problemas computacionales}, y observamos que los programas en Python (y, equivalentemente, las máquinas de Turing), pueden resolverlos. Llamamos \emph{decidibles} a aquellos problemas que pueden resolverse mediante un programa que nunca para. En la \cref{sec:universalidad} observamos que un programa puede simularse a sí mismo, lo cual nos conduce en la \cref{sec:problemas-no-decidibles} a la existencia de problemas no decidibles. En la \cref{sec:problemas-semidecidibles} relajamos la condición de parada para obtener problemas \emph{semidecidibles}. A continuación, introducimos en la \cref{sec:reducciones} la técnica de \emph{reducción}, que nos permite obtener nuevos problemas no decidibles, incluyendo el más importante del trabajo en la \cref{sec:problema-parada}: el \emph{problema de la parada}.

Para demostrar el teorema principal de este trabajo, es necesario formalizar los conceptos de \emph{demostración} y \emph{verdad}. Hacemos esto introduciendo los \emph{sistemas formales} en la \cref{sec:sistemas-formales} y los \emph{sistemas lógicos} en la \cref{sec:sistemas-logicos}. Los sistemas formales son \emph{sintácticos}: formalizan qué es un \emph{teorema}, una \emph{fórmula} demostrable a partir de un conjunto de \emph{axiomas} y \emph{reglas de inferencia}. Por otro lado, los sistemas lógicos incorporan a los sistemas formales un componente semántico: la \emph{verdad}. En las \cref{sec:solidez-completitud-decidibilidad,,sec:consistencia} presentamos propiedades esenciales de los sistemas mencionados, que relacionan los conceptos semánticos y sintácticos. Destacamos en la \cref{sec:aritmetica-peano} la introducción de la \emph{aritmética de Peano}: un sistema lógico de especial importancia.

Como culmen, en el \cref{ch:teorema-incompletitud} demostramos el Primer Teorema de Incompletitud: primero desde la aritmética de Peano en la \cref{sec:primera-aproximacion}, para luego generalizarlo mediante dos versiones. La primera, semántica, es demostrada en la \cref{sec:version-semantica}, y nos permite (bajo unas hipótesis computacionales) encontrar afirmaciones ciertas e indemostrables en sistemas lógicos. A continuación, la versión sintáctica de la \cref{sec:version-sintactica}, más cercana a la de Gödel, nos conduce a la existencia de sentencias indemostrables en sistemas formales. Finalizamos con la \cref{sec:consecuencias}, estudiando importantes consecuencias matemáticas (\cref{subsec:segundo-problema-hilbert}) y filosóficas (\cref{subsec:mente-humana-tesis-church-turing}).

Finalmente, terminamos con algunas conclusiones en el \cref{ch:conclusiones}, y dedicando unas palabras finales para el trabajo posterior en el \cref{ch:trabajo-futuro}.


%\vfill
% ====================
%\begin{figure}[H]
%\centering
%\vspace{8pt}
%\hspace*{-1cm}
%\centerline{% !TeX root = ../libro.tex
% !TeX encoding = utf8
%
%*******************************************************
% Introducción
%*******************************************************

% \manualmark
% \markboth{\textsc{Introducción}}{\textsc{Introducción}} 
\chapter{Resumen}

%De acuerdo con la comisión de grado, el TFG debe incluir una introducción en la que se describan claramente los objetivos previstos inicialmente en la propuesta de TFG, indicando si han sido o no alcanzados, los antecedentes importantes para el desarrollo, los resultados obtenidos, en su caso y las principales fuentes consultadas.

%En este capítulo, resumiremos el proceso que hemos seguido hasta llegar a la demostración, desde una perspectiva algorítmica, del Primer Teorema de Incompletitud. En la \cref{fig:resumen}, se puede observar un esquema resumen de este trabajo, el cual es extendido en la \cref{fig:resumen-extendido}. En este resumen final se omiten las \cref{part:introduccion,,part:conclusion}.

Comenzamos este Trabajo de Fin de Grado con el \cref{ch:motivacion}, que destaca la motivación que subyace este trabajo, así como los objetivos que se pretenden cumplir. Inmediatamente después, en el \cref{ch:historia}, hacemos una breve contextualización histórica de esta obra.

El desarrollo teórico comienza con el \cref{ch:programas-python-maquinas-turing}, en cuya \cref{sec:programas-python} introducimos una clase específica de programas en Python, los \emph{programas SISO}, y observamos que esta clase es tan expresiva como la clase de todos los programas en Python. A continuación, en la \cref{sec:maquinas-turing}, profundizamos en las \emph{máquinas de Turing}, herramienta matemática principal de nuestro trabajo. En la \cref{sec:equivalencia}, introducimos nuevos modelos de computación equivalentes, hasta llegar al resultado principal del capítulo: que las máquinas de Turing pueden simular programas en Python, y viceversa; en otras palabras, que los programas en Python y las máquinas de Turing son equivalentes. Esto nos permitirá usar programas en Python convenientemente, sin realizar descripciones de máquinas de Turing. Concluimos introduciendo la tesis de Church-Turing en la \cref{sec:church-turing}, y viendo que la equivalencia probada puede ser un argumento a su favor.

En la \cref{sec:problemas-decidibles} definimos los \emph{problemas computacionales}, y observamos que los programas en Python (y, equivalentemente, las máquinas de Turing), pueden resolverlos. Llamamos \emph{decidibles} a aquellos problemas que pueden resolverse mediante un programa que nunca para. En la \cref{sec:universalidad} observamos que un programa puede simularse a sí mismo, lo cual nos conduce en la \cref{sec:problemas-no-decidibles} a la existencia de problemas no decidibles. En la \cref{sec:problemas-semidecidibles} relajamos la condición de parada para obtener problemas \emph{semidecidibles}. A continuación, introducimos en la \cref{sec:reducciones} la técnica de \emph{reducción}, que nos permite obtener nuevos problemas no decidibles, incluyendo el más importante del trabajo en la \cref{sec:problema-parada}: el \emph{problema de la parada}.

Para demostrar el teorema principal de este trabajo, es necesario formalizar los conceptos de \emph{demostración} y \emph{verdad}. Hacemos esto introduciendo los \emph{sistemas formales} en la \cref{sec:sistemas-formales} y los \emph{sistemas lógicos} en la \cref{sec:sistemas-logicos}. Los sistemas formales son \emph{sintácticos}: formalizan qué es un \emph{teorema}, una \emph{fórmula} demostrable a partir de un conjunto de \emph{axiomas} y \emph{reglas de inferencia}. Por otro lado, los sistemas lógicos incorporan a los sistemas formales un componente semántico: la \emph{verdad}. En las \cref{sec:solidez-completitud-decidibilidad,,sec:consistencia} presentamos propiedades esenciales de los sistemas mencionados, que relacionan los conceptos semánticos y sintácticos. Destacamos en la \cref{sec:aritmetica-peano} la introducción de la \emph{aritmética de Peano}: un sistema lógico de especial importancia.

Como culmen, en el \cref{ch:teorema-incompletitud} demostramos el Primer Teorema de Incompletitud: primero desde la aritmética de Peano en la \cref{sec:primera-aproximacion}, para luego generalizarlo mediante dos versiones. La primera, semántica, es demostrada en la \cref{sec:version-semantica}, y nos permite (bajo unas hipótesis computacionales) encontrar afirmaciones ciertas e indemostrables en sistemas lógicos. A continuación, la versión sintáctica de la \cref{sec:version-sintactica}, más cercana a la de Gödel, nos conduce a la existencia de sentencias indemostrables en sistemas formales. Finalizamos con la \cref{sec:consecuencias}, estudiando importantes consecuencias matemáticas (\cref{subsec:segundo-problema-hilbert}) y filosóficas (\cref{subsec:mente-humana-tesis-church-turing}).

Finalmente, terminamos con algunas conclusiones en el \cref{ch:conclusiones}, y dedicando unas palabras finales para el trabajo posterior en el \cref{ch:trabajo-futuro}.


%\vfill
% ====================
%\begin{figure}[H]
%\centering
%\vspace{8pt}
%\hspace*{-1cm}
%\centerline{% !TeX root = ../libro.tex
% !TeX encoding = utf8
%
%*******************************************************
% Introducción
%*******************************************************

% \manualmark
% \markboth{\textsc{Introducción}}{\textsc{Introducción}} 
\chapter{Resumen}

%De acuerdo con la comisión de grado, el TFG debe incluir una introducción en la que se describan claramente los objetivos previstos inicialmente en la propuesta de TFG, indicando si han sido o no alcanzados, los antecedentes importantes para el desarrollo, los resultados obtenidos, en su caso y las principales fuentes consultadas.

%En este capítulo, resumiremos el proceso que hemos seguido hasta llegar a la demostración, desde una perspectiva algorítmica, del Primer Teorema de Incompletitud. En la \cref{fig:resumen}, se puede observar un esquema resumen de este trabajo, el cual es extendido en la \cref{fig:resumen-extendido}. En este resumen final se omiten las \cref{part:introduccion,,part:conclusion}.

Comenzamos este Trabajo de Fin de Grado con el \cref{ch:motivacion}, que destaca la motivación que subyace este trabajo, así como los objetivos que se pretenden cumplir. Inmediatamente después, en el \cref{ch:historia}, hacemos una breve contextualización histórica de esta obra.

El desarrollo teórico comienza con el \cref{ch:programas-python-maquinas-turing}, en cuya \cref{sec:programas-python} introducimos una clase específica de programas en Python, los \emph{programas SISO}, y observamos que esta clase es tan expresiva como la clase de todos los programas en Python. A continuación, en la \cref{sec:maquinas-turing}, profundizamos en las \emph{máquinas de Turing}, herramienta matemática principal de nuestro trabajo. En la \cref{sec:equivalencia}, introducimos nuevos modelos de computación equivalentes, hasta llegar al resultado principal del capítulo: que las máquinas de Turing pueden simular programas en Python, y viceversa; en otras palabras, que los programas en Python y las máquinas de Turing son equivalentes. Esto nos permitirá usar programas en Python convenientemente, sin realizar descripciones de máquinas de Turing. Concluimos introduciendo la tesis de Church-Turing en la \cref{sec:church-turing}, y viendo que la equivalencia probada puede ser un argumento a su favor.

En la \cref{sec:problemas-decidibles} definimos los \emph{problemas computacionales}, y observamos que los programas en Python (y, equivalentemente, las máquinas de Turing), pueden resolverlos. Llamamos \emph{decidibles} a aquellos problemas que pueden resolverse mediante un programa que nunca para. En la \cref{sec:universalidad} observamos que un programa puede simularse a sí mismo, lo cual nos conduce en la \cref{sec:problemas-no-decidibles} a la existencia de problemas no decidibles. En la \cref{sec:problemas-semidecidibles} relajamos la condición de parada para obtener problemas \emph{semidecidibles}. A continuación, introducimos en la \cref{sec:reducciones} la técnica de \emph{reducción}, que nos permite obtener nuevos problemas no decidibles, incluyendo el más importante del trabajo en la \cref{sec:problema-parada}: el \emph{problema de la parada}.

Para demostrar el teorema principal de este trabajo, es necesario formalizar los conceptos de \emph{demostración} y \emph{verdad}. Hacemos esto introduciendo los \emph{sistemas formales} en la \cref{sec:sistemas-formales} y los \emph{sistemas lógicos} en la \cref{sec:sistemas-logicos}. Los sistemas formales son \emph{sintácticos}: formalizan qué es un \emph{teorema}, una \emph{fórmula} demostrable a partir de un conjunto de \emph{axiomas} y \emph{reglas de inferencia}. Por otro lado, los sistemas lógicos incorporan a los sistemas formales un componente semántico: la \emph{verdad}. En las \cref{sec:solidez-completitud-decidibilidad,,sec:consistencia} presentamos propiedades esenciales de los sistemas mencionados, que relacionan los conceptos semánticos y sintácticos. Destacamos en la \cref{sec:aritmetica-peano} la introducción de la \emph{aritmética de Peano}: un sistema lógico de especial importancia.

Como culmen, en el \cref{ch:teorema-incompletitud} demostramos el Primer Teorema de Incompletitud: primero desde la aritmética de Peano en la \cref{sec:primera-aproximacion}, para luego generalizarlo mediante dos versiones. La primera, semántica, es demostrada en la \cref{sec:version-semantica}, y nos permite (bajo unas hipótesis computacionales) encontrar afirmaciones ciertas e indemostrables en sistemas lógicos. A continuación, la versión sintáctica de la \cref{sec:version-sintactica}, más cercana a la de Gödel, nos conduce a la existencia de sentencias indemostrables en sistemas formales. Finalizamos con la \cref{sec:consecuencias}, estudiando importantes consecuencias matemáticas (\cref{subsec:segundo-problema-hilbert}) y filosóficas (\cref{subsec:mente-humana-tesis-church-turing}).

Finalmente, terminamos con algunas conclusiones en el \cref{ch:conclusiones}, y dedicando unas palabras finales para el trabajo posterior en el \cref{ch:trabajo-futuro}.


%\vfill
% ====================
%\begin{figure}[H]
%\centering
%\vspace{8pt}
%\hspace*{-1cm}
%\centerline{\input{fig/resumen}}
%\caption{Esquema resumen}
%\label{fig:resumen}
%\end{figure}
% ====================
%\vfill

\begin{landscape}
% ====================
\begin{figure}[H]
%\captionsetup{margin=200cm}

\centering
%\vspace{-1.5cm}
%\hspace{-1cm}
%\input{fig/resumen-extendido}
\raisebox{-15.5cm}[0pt][0pt]{%
\resizebox{\linewidth}{!}{\input{fig/resumen-extendido}}
}
%\caption{Esquema resumen (extendido)}

%\label{fig:resumen-extendido}
\end{figure}
% ====================
\end{landscape}

\endinput
}
%\caption{Esquema resumen}
%\label{fig:resumen}
%\end{figure}
% ====================
%\vfill

\begin{landscape}
% ====================
\begin{figure}[H]
%\captionsetup{margin=200cm}

\centering
%\vspace{-1.5cm}
%\hspace{-1cm}
%

\tikzset{every picture/.style={line width=0.75pt}} %set default line width to 0.75pt        

\begin{tikzpicture}[x=0.75pt,y=0.75pt,yscale=-1,xscale=1]
%uncomment if require: \path (0,694); %set diagram left start at 0, and has height of 694

%Straight Lines [id:da2056540562354816] 
\draw  [dash pattern={on 4.5pt off 4.5pt}]  (672,64.47) -- (810,64.47) ;
%Straight Lines [id:da23044041717189923] 
\draw  [dash pattern={on 4.5pt off 4.5pt}]  (412.67,64.47) -- (610,64.47) ;
\draw [shift={(410.67,64.47)}, rotate = 0] [color={rgb, 255:red, 0; green, 0; blue, 0 }  ][line width=0.75]    (10.93,-3.29) .. controls (6.95,-1.4) and (3.31,-0.3) .. (0,0) .. controls (3.31,0.3) and (6.95,1.4) .. (10.93,3.29)   ;
%Straight Lines [id:da8398435684999113] 
\draw  [dash pattern={on 4.5pt off 4.5pt}]  (810,64.47) -- (810,530) ;
%Rounded Rect [id:dp027018508739026448] 
\draw  [draw opacity=0][fill={rgb, 255:red, 230; green, 230; blue, 230 }  ,fill opacity=0.5 ] (20,439.33) .. controls (20,445.22) and (24.78,450) .. (30.67,450) -- (409.33,450) .. controls (415.22,450) and (420,445.22) .. (420,439.33) -- (420,140.67) .. controls (420,134.78) and (415.22,130) .. (409.33,130) -- (30.67,130) .. controls (24.78,130) and (20,134.78) .. (20,140.67) -- cycle ;
%Rounded Rect [id:dp9415679498487528] 
\draw  [dash pattern={on 4.5pt off 4.5pt}] (29.83,292.77) .. controls (29.83,285.88) and (35.41,280.3) .. (42.3,280.3) -- (398.54,280.3) .. controls (405.42,280.3) and (411,285.88) .. (411,292.77) -- (411,337.17) .. controls (411,344.05) and (405.42,349.64) .. (398.54,349.64) -- (42.3,349.64) .. controls (35.41,349.64) and (29.83,344.05) .. (29.83,337.17) -- cycle ;
%Rounded Rect [id:dp05100022295548956] 
\draw  [draw opacity=0][fill={rgb, 255:red, 230; green, 230; blue, 230 }  ,fill opacity=0.5 ] (130.5,575) .. controls (130.5,583.28) and (137.22,590) .. (145.5,590) -- (615,590) .. controls (623.28,590) and (630,583.28) .. (630,575) -- (630,395) .. controls (630,386.72) and (623.28,380) .. (615,380) -- (145.5,380) .. controls (137.22,380) and (130.5,386.72) .. (130.5,395) -- cycle ;
%Rounded Rect [id:dp6496809750851564] 
\draw  [draw opacity=0][fill={rgb, 255:red, 230; green, 230; blue, 230 }  ,fill opacity=0.5 ] (450.61,440.79) .. controls (450.61,445.82) and (454.69,449.9) .. (459.72,449.9) -- (780.89,449.9) .. controls (785.92,449.9) and (790,445.82) .. (790,440.79) -- (790,139.11) .. controls (790,134.08) and (785.92,130) .. (780.89,130) -- (459.72,130) .. controls (454.69,130) and (450.61,134.08) .. (450.61,139.11) -- cycle ;
%Rounded Rect [id:dp549636265291094] 
\draw  [fill={rgb, 255:red, 230; green, 230; blue, 230 }  ,fill opacity=0.5 ][dash pattern={on 4.5pt off 4.5pt}] (29.83,42.77) .. controls (29.83,35.88) and (35.41,30.3) .. (42.3,30.3) -- (398.54,30.3) .. controls (405.42,30.3) and (411,35.88) .. (411,42.77) -- (411,87.17) .. controls (411,94.05) and (405.42,99.64) .. (398.54,99.64) -- (42.3,99.64) .. controls (35.41,99.64) and (29.83,94.05) .. (29.83,87.17) -- cycle ;
%Rounded Rect [id:dp1378581119146769] 
\draw  [fill={rgb, 255:red, 255; green, 255; blue, 255 }  ,fill opacity=1 ] (40.4,49.82) .. controls (40.4,44.29) and (44.89,39.8) .. (50.42,39.8) -- (169.58,39.8) .. controls (175.11,39.8) and (179.6,44.29) .. (179.6,49.82) -- (179.6,79.88) .. controls (179.6,85.41) and (175.11,89.9) .. (169.58,89.9) -- (50.42,89.9) .. controls (44.89,89.9) and (40.4,85.41) .. (40.4,79.88) -- cycle ;
%Straight Lines [id:da8729153203483555] 
\draw    (188.42,63.77) -- (250.99,63.77)(188.42,66.77) -- (250.99,66.77) ;
\draw [shift={(257.99,65.27)}, rotate = 180] [color={rgb, 255:red, 0; green, 0; blue, 0 }  ][line width=0.75]    (10.93,-4.9) .. controls (6.95,-2.3) and (3.31,-0.67) .. (0,0) .. controls (3.31,0.67) and (6.95,2.3) .. (10.93,4.9)   ;
\draw [shift={(181.42,65.27)}, rotate = 0] [color={rgb, 255:red, 0; green, 0; blue, 0 }  ][line width=0.75]    (10.93,-4.9) .. controls (6.95,-2.3) and (3.31,-0.67) .. (0,0) .. controls (3.31,0.67) and (6.95,2.3) .. (10.93,4.9)   ;
%Straight Lines [id:da845243590120559] 
\draw    (220.17,99.97) -- (220.17,137.64) ;
\draw [shift={(220.17,139.64)}, rotate = 270] [color={rgb, 255:red, 0; green, 0; blue, 0 }  ][line width=0.75]    (10.93,-3.29) .. controls (6.95,-1.4) and (3.31,-0.3) .. (0,0) .. controls (3.31,0.3) and (6.95,1.4) .. (10.93,3.29)   ;
%Straight Lines [id:da7497001546502899] 
\draw    (380,420.33) -- (380,478) ;
\draw [shift={(380,480)}, rotate = 270] [color={rgb, 255:red, 0; green, 0; blue, 0 }  ][line width=0.75]    (10.93,-3.29) .. controls (6.95,-1.4) and (3.31,-0.3) .. (0,0) .. controls (3.31,0.3) and (6.95,1.4) .. (10.93,3.29)   ;
%Straight Lines [id:da5689893955334542] 
\draw    (300,420) -- (460,420) ;
%Rounded Rect [id:dp18582062333838767] 
\draw  [fill={rgb, 255:red, 255; green, 255; blue, 255 }  ,fill opacity=1 ] (260.4,49.82) .. controls (260.4,44.29) and (264.89,39.8) .. (270.42,39.8) -- (389.58,39.8) .. controls (395.11,39.8) and (399.6,44.29) .. (399.6,49.82) -- (399.6,79.88) .. controls (399.6,85.41) and (395.11,89.9) .. (389.58,89.9) -- (270.42,89.9) .. controls (264.89,89.9) and (260.4,85.41) .. (260.4,79.88) -- cycle ;
%Rounded Rect [id:dp9523169763067778] 
\draw  [fill={rgb, 255:red, 255; green, 255; blue, 255 }  ,fill opacity=1 ] (470,299.92) .. controls (470,294.39) and (474.49,289.9) .. (480.02,289.9) -- (599.18,289.9) .. controls (604.71,289.9) and (609.2,294.39) .. (609.2,299.92) -- (609.2,329.98) .. controls (609.2,335.51) and (604.71,340) .. (599.18,340) -- (480.02,340) .. controls (474.49,340) and (470,335.51) .. (470,329.98) -- cycle ;
%Rounded Rect [id:dp48544030815291705] 
\draw  [fill={rgb, 255:red, 255; green, 255; blue, 255 }  ,fill opacity=1 ] (220,490) .. controls (220,484.48) and (224.48,480) .. (230,480) -- (530,480) .. controls (535.52,480) and (540,484.48) .. (540,490) -- (540,570) .. controls (540,575.52) and (535.52,580) .. (530,580) -- (230,580) .. controls (224.48,580) and (220,575.52) .. (220,570) -- cycle ;
%Rounded Rect [id:dp08222627508455749] 
\draw  [fill={rgb, 255:red, 255; green, 255; blue, 255 }  ,fill opacity=1 ] (540.4,49.82) .. controls (540.4,44.29) and (544.89,39.8) .. (550.42,39.8) -- (690.65,39.8) .. controls (696.18,39.8) and (700.67,44.29) .. (700.67,49.82) -- (700.67,79.88) .. controls (700.67,85.41) and (696.18,89.9) .. (690.65,89.9) -- (550.42,89.9) .. controls (544.89,89.9) and (540.4,85.41) .. (540.4,79.88) -- cycle ;
%Straight Lines [id:da5868346364154593] 
\draw  [dash pattern={on 4.5pt off 4.5pt}]  (810,530) -- (542,530) ;
\draw [shift={(540,530)}, rotate = 360] [color={rgb, 255:red, 0; green, 0; blue, 0 }  ][line width=0.75]    (10.93,-3.29) .. controls (6.95,-1.4) and (3.31,-0.3) .. (0,0) .. controls (3.31,0.3) and (6.95,1.4) .. (10.93,3.29)   ;
%Rounded Rect [id:dp7392784471609453] 
\draw  [fill={rgb, 255:red, 255; green, 255; blue, 255 }  ,fill opacity=1 ] (30.33,150.47) .. controls (30.33,144.58) and (35.11,139.8) .. (41,139.8) -- (400.34,139.8) .. controls (406.23,139.8) and (411,144.58) .. (411,150.47) -- (411,229.34) .. controls (411,235.23) and (406.23,240.01) .. (400.34,240.01) -- (41,240.01) .. controls (35.11,240.01) and (30.33,235.23) .. (30.33,229.34) -- cycle ;
%Rounded Rect [id:dp90901637181607] 
\draw  [fill={rgb, 255:red, 255; green, 255; blue, 255 }  ,fill opacity=1 ] (41,189.82) .. controls (41,184.29) and (45.49,179.8) .. (51.02,179.8) -- (145.02,179.8) .. controls (150.55,179.8) and (155.04,184.29) .. (155.04,189.82) -- (155.04,219.88) .. controls (155.04,225.41) and (150.55,229.9) .. (145.02,229.9) -- (51.02,229.9) .. controls (45.49,229.9) and (41,225.41) .. (41,219.88) -- cycle ;
%Rounded Rect [id:dp06115018179625542] 
\draw  [fill={rgb, 255:red, 255; green, 255; blue, 255 }  ,fill opacity=1 ] (163.65,189.82) .. controls (163.65,184.29) and (168.13,179.8) .. (173.67,179.8) -- (267.67,179.8) .. controls (273.2,179.8) and (277.69,184.29) .. (277.69,189.82) -- (277.69,219.88) .. controls (277.69,225.41) and (273.2,229.9) .. (267.67,229.9) -- (173.67,229.9) .. controls (168.13,229.9) and (163.65,225.41) .. (163.65,219.88) -- cycle ;
%Rounded Rect [id:dp6195275548153432] 
\draw  [fill={rgb, 255:red, 255; green, 255; blue, 255 }  ,fill opacity=1 ] (286.3,189.82) .. controls (286.3,184.29) and (290.78,179.8) .. (296.32,179.8) -- (390.31,179.8) .. controls (395.85,179.8) and (400.33,184.29) .. (400.33,189.82) -- (400.33,219.88) .. controls (400.33,225.41) and (395.85,229.9) .. (390.31,229.9) -- (296.32,229.9) .. controls (290.78,229.9) and (286.3,225.41) .. (286.3,219.88) -- cycle ;
%Rounded Rect [id:dp41283415100159493] 
\draw  [fill={rgb, 255:red, 255; green, 255; blue, 255 }  ,fill opacity=1 ] (40.8,299.92) .. controls (40.8,294.39) and (45.29,289.9) .. (50.82,289.9) -- (169.98,289.9) .. controls (175.51,289.9) and (180,294.39) .. (180,299.92) -- (180,329.98) .. controls (180,335.51) and (175.51,340) .. (169.98,340) -- (50.82,340) .. controls (45.29,340) and (40.8,335.51) .. (40.8,329.98) -- cycle ;
%Rounded Rect [id:dp40154854661776396] 
\draw  [fill={rgb, 255:red, 255; green, 255; blue, 255 }  ,fill opacity=1 ] (260.4,299.82) .. controls (260.4,294.29) and (264.89,289.8) .. (270.42,289.8) -- (389.58,289.8) .. controls (395.11,289.8) and (399.6,294.29) .. (399.6,299.82) -- (399.6,329.88) .. controls (399.6,335.41) and (395.11,339.9) .. (389.58,339.9) -- (270.42,339.9) .. controls (264.89,339.9) and (260.4,335.41) .. (260.4,329.88) -- cycle ;
%Straight Lines [id:da6602630444249074] 
\draw    (180.17,315.8) -- (260,315.8) ;
%Straight Lines [id:da30905866347109745] 
\draw    (219.86,315.97) -- (219.86,386.81) ;
\draw [shift={(219.86,388.81)}, rotate = 270] [color={rgb, 255:red, 0; green, 0; blue, 0 }  ][line width=0.75]    (10.93,-3.29) .. controls (6.95,-1.4) and (3.31,-0.3) .. (0,0) .. controls (3.31,0.3) and (6.95,1.4) .. (10.93,3.29)   ;
%Rounded Rect [id:dp44627062039563214] 
\draw  [fill={rgb, 255:red, 255; green, 255; blue, 255 }  ,fill opacity=1 ] (139.67,399.82) .. controls (139.67,394.29) and (144.15,389.8) .. (149.69,389.8) -- (290.31,389.8) .. controls (295.85,389.8) and (300.33,394.29) .. (300.33,399.82) -- (300.33,429.88) .. controls (300.33,435.41) and (295.85,439.9) .. (290.31,439.9) -- (149.69,439.9) .. controls (144.15,439.9) and (139.67,435.41) .. (139.67,429.88) -- cycle ;
%Straight Lines [id:da8375870493480149] 
\draw    (220,240) -- (220,277.67) ;
\draw [shift={(220,279.67)}, rotate = 270] [color={rgb, 255:red, 0; green, 0; blue, 0 }  ][line width=0.75]    (10.93,-3.29) .. controls (6.95,-1.4) and (3.31,-0.3) .. (0,0) .. controls (3.31,0.3) and (6.95,1.4) .. (10.93,3.29)   ;
%Rounded Rect [id:dp37285624114093485] 
\draw  [fill={rgb, 255:red, 255; green, 255; blue, 255 }  ,fill opacity=1 ] (460,400.02) .. controls (460,394.49) and (464.49,390) .. (470.02,390) -- (610.65,390) .. controls (616.18,390) and (620.67,394.49) .. (620.67,400.02) -- (620.67,430.08) .. controls (620.67,435.61) and (616.18,440.1) .. (610.65,440.1) -- (470.02,440.1) .. controls (464.49,440.1) and (460,435.61) .. (460,430.08) -- cycle ;
%Rounded Rect [id:dp6883616335272991] 
\draw  [fill={rgb, 255:red, 255; green, 255; blue, 255 }  ,fill opacity=1 ] (629.2,299.92) .. controls (629.2,294.39) and (633.69,289.9) .. (639.22,289.9) -- (758.38,289.9) .. controls (763.91,289.9) and (768.4,294.39) .. (768.4,299.92) -- (768.4,329.98) .. controls (768.4,335.51) and (763.91,340) .. (758.38,340) -- (639.22,340) .. controls (633.69,340) and (629.2,335.51) .. (629.2,329.98) -- cycle ;
%Straight Lines [id:da19870419404011752] 
\draw    (540,340) -- (540,388) ;
\draw [shift={(540,390)}, rotate = 270] [color={rgb, 255:red, 0; green, 0; blue, 0 }  ][line width=0.75]    (10.93,-3.29) .. controls (6.95,-1.4) and (3.31,-0.3) .. (0,0) .. controls (3.31,0.3) and (6.95,1.4) .. (10.93,3.29)   ;
%Rounded Rect [id:dp9692784599420237] 
\draw  [fill={rgb, 255:red, 255; green, 255; blue, 255 }  ,fill opacity=1 ] (230,529.92) .. controls (230,524.39) and (234.49,519.9) .. (240.02,519.9) -- (359.18,519.9) .. controls (364.71,519.9) and (369.2,524.39) .. (369.2,529.92) -- (369.2,559.98) .. controls (369.2,565.51) and (364.71,570) .. (359.18,570) -- (240.02,570) .. controls (234.49,570) and (230,565.51) .. (230,559.98) -- cycle ;
%Rounded Rect [id:dp749684165723072] 
\draw  [fill={rgb, 255:red, 255; green, 255; blue, 255 }  ,fill opacity=1 ] (390,530.02) .. controls (390,524.49) and (394.49,520) .. (400.02,520) -- (519.18,520) .. controls (524.71,520) and (529.2,524.49) .. (529.2,530.02) -- (529.2,560.08) .. controls (529.2,565.61) and (524.71,570.1) .. (519.18,570.1) -- (400.02,570.1) .. controls (394.49,570.1) and (390,565.61) .. (390,560.08) -- cycle ;
%Rounded Rect [id:dp28001422897762485] 
\draw  [fill={rgb, 255:red, 255; green, 255; blue, 255 }  ,fill opacity=1 ][dash pattern={on 4.5pt off 4.5pt}] (460,148.92) .. controls (460,143.99) and (463.99,140) .. (468.92,140) -- (771.08,140) .. controls (776.01,140) and (780,143.99) .. (780,148.92) -- (780,241.55) .. controls (780,246.48) and (776.01,250.47) .. (771.08,250.47) -- (468.92,250.47) .. controls (463.99,250.47) and (460,246.48) .. (460,241.55) -- cycle ;
%Shape: Brace [id:dp9411387904179627] 
\draw   (774.67,286.8) .. controls (774.67,282.13) and (772.34,279.8) .. (767.67,279.8) -- (629.96,279.8) .. controls (623.29,279.8) and (619.96,277.47) .. (619.96,272.8) .. controls (619.96,277.47) and (616.63,279.8) .. (609.96,279.8)(612.96,279.8) -- (472.25,279.8) .. controls (467.58,279.8) and (465.25,282.13) .. (465.25,286.8) ;
%Straight Lines [id:da9902399507274253] 
\draw    (620,273) -- (620,253.14) ;
\draw [shift={(620,251.14)}, rotate = 90] [color={rgb, 255:red, 0; green, 0; blue, 0 }  ][line width=0.75]    (10.93,-3.29) .. controls (6.95,-1.4) and (3.31,-0.3) .. (0,0) .. controls (3.31,0.3) and (6.95,1.4) .. (10.93,3.29)   ;

% Text Node
\draw (110,64.85) node  [color={rgb, 255:red, 0; green, 0; blue, 0 }  ,opacity=1 ] [align=left] {\begin{minipage}[lt]{92.93pt}\setlength\topsep{0pt}
\begin{center}
programa en Python\\
\small{(\cref{sec:programas-python})}
\end{center}

\end{minipage}};
% Text Node
\draw (219.89,51.97) node  [font=\small,color={rgb, 255:red, 0; green, 0; blue, 0 }  ,opacity=1 ] [align=left] {\begin{minipage}[lt]{54.86pt}\setlength\topsep{0pt}
\begin{center}
\Cref{sec:equivalencia}
\end{center}

\end{minipage}};
% Text Node
\draw (69.5,15.5) node  [font=\small,color={rgb, 255:red, 0; green, 0; blue, 0 }  ,opacity=1 ] [align=left] {\begin{minipage}[lt]{41.55pt}\setlength\topsep{0pt}
\cref{ch:programas-python-maquinas-turing}
\end{minipage}};
% Text Node
\draw (330,64.85) node  [color={rgb, 255:red, 0; green, 0; blue, 0 }  ,opacity=1 ] [align=left] {\begin{minipage}[lt]{92.93pt}\setlength\topsep{0pt}
\begin{center}
máquina de Turing\\
\small{(\cref{sec:maquinas-turing})}
\end{center}

\end{minipage}};
% Text Node
\draw (539.6,314.95) node  [color={rgb, 255:red, 0; green, 0; blue, 0 }  ,opacity=1 ] [align=left] {\begin{minipage}[lt]{92.93pt}\setlength\topsep{0pt}
\begin{center}
sistema lógico\\
\small{(\cref{sec:sistemas-logicos})}
\end{center}

\end{minipage}};
% Text Node
\draw (380.5,500.11) node  [color={rgb, 255:red, 0; green, 0; blue, 0 }  ,opacity=1 ] [align=left] {\begin{minipage}[lt]{159.55pt}\setlength\topsep{0pt}
\begin{center}
\textbf{Primer Teorema de Incompletitud}
\end{center}

\end{minipage}};
% Text Node
\draw (620.53,64.85) node  [color={rgb, 255:red, 0; green, 0; blue, 0 }  ,opacity=1 ] [align=left] {\begin{minipage}[lt]{106.99pt}\setlength\topsep{0pt}
\begin{center}
tesis de Church-Turing\\
\small{(\cref{sec:church-turing})}
\end{center}

\end{minipage}};
% Text Node
\draw (480.28,52.8) node  [font=\small,color={rgb, 255:red, 0; green, 0; blue, 0 }  ,opacity=1 ] [align=left] {\begin{minipage}[lt]{81.9pt}\setlength\topsep{0pt}
\begin{center}
\cref{sec:church-turing}
\end{center}

\end{minipage}};
% Text Node
\draw (589.5,518.3) node  [font=\small,color={rgb, 255:red, 0; green, 0; blue, 0 }  ,opacity=1 ] [align=left] {\begin{minipage}[lt]{54.86pt}\setlength\topsep{0pt}
\begin{center}
\Cref{subsec:mente-humana-tesis-church-turing}
\end{center}

\end{minipage}};
% Text Node
\draw (220.2,159.62) node  [color={rgb, 255:red, 0; green, 0; blue, 0 }  ,opacity=1 ] [align=left] {\begin{minipage}[lt]{230.11pt}\setlength\topsep{0pt}
\begin{center}
problema computacional
\end{center}

\end{minipage}};
% Text Node
\draw (98.02,204.85) node  [color={rgb, 255:red, 0; green, 0; blue, 0 }  ,opacity=1 ] [align=left] {\begin{minipage}[lt]{76.13pt}\setlength\topsep{0pt}
\begin{center}
decidible\\
\small{(\cref{sec:problemas-decidibles})}
\end{center}

\end{minipage}};
% Text Node
\draw (220.67,204.85) node  [color={rgb, 255:red, 0; green, 0; blue, 0 }  ,opacity=1 ] [align=left] {\begin{minipage}[lt]{76.13pt}\setlength\topsep{0pt}
\begin{center}
no decidible\\
\small{(\cref{sec:problemas-no-decidibles})}
\end{center}

\end{minipage}};
% Text Node
\draw (343.31,204.85) node  [color={rgb, 255:red, 0; green, 0; blue, 0 }  ,opacity=1 ] [align=left] {\begin{minipage}[lt]{76.13pt}\setlength\topsep{0pt}
\begin{center}
semidecidible\\
\small{(\cref{sec:problemas-semidecidibles})}
\end{center}

\end{minipage}};
% Text Node
\draw (110.4,314.95) node  [color={rgb, 255:red, 0; green, 0; blue, 0 }  ,opacity=1 ] [align=left] {\begin{minipage}[lt]{92.93pt}\setlength\topsep{0pt}
\begin{center}
universalidad\\
\small{(\cref{sec:universalidad})}
\end{center}

\end{minipage}};
% Text Node
\draw (330,314.85) node  [color={rgb, 255:red, 0; green, 0; blue, 0 }  ,opacity=1 ] [align=left] {\begin{minipage}[lt]{92.93pt}\setlength\topsep{0pt}
\begin{center}
reducción\\
\small{(\cref{sec:reducciones})}
\end{center}

\end{minipage}};
% Text Node
\draw (220,414.85) node  [color={rgb, 255:red, 0; green, 0; blue, 0 }  ,opacity=1 ] [align=left] {\begin{minipage}[lt]{107.26pt}\setlength\topsep{0pt}
\begin{center}
problema de la parada\\
\small{(\cref{sec:problema-parada})}
\end{center}

\end{minipage}};
% Text Node
\draw (540.33,415.05) node  [color={rgb, 255:red, 0; green, 0; blue, 0 }  ,opacity=1 ] [align=left] {\begin{minipage}[lt]{107.26pt}\setlength\topsep{0pt}
\begin{center}
aritmética de Peano\\
\small{(\cref{sec:aritmetica-peano})}
\end{center}

\end{minipage}};
% Text Node
\draw (698.8,314.95) node  [color={rgb, 255:red, 0; green, 0; blue, 0 }  ,opacity=1 ] [align=left] {\begin{minipage}[lt]{92.93pt}\setlength\topsep{0pt}
\begin{center}
sistema formal\\
\small{(\cref{sec:sistemas-formales})}
\end{center}

\end{minipage}};
% Text Node
\draw (299.6,544.95) node  [color={rgb, 255:red, 0; green, 0; blue, 0 }  ,opacity=1 ] [align=left] {\begin{minipage}[lt]{92.93pt}\setlength\topsep{0pt}
\begin{center}
versión semántica\\
\small{(\cref{sec:version-semantica})}
\end{center}

\end{minipage}};
% Text Node
\draw (459.6,545.05) node  [color={rgb, 255:red, 0; green, 0; blue, 0 }  ,opacity=1 ] [align=left] {\begin{minipage}[lt]{92.93pt}\setlength\topsep{0pt}
\begin{center}
versión sintáctica\\
\small{(\cref{sec:version-sintactica})}
\end{center}

\end{minipage}};
% Text Node
\draw (60.5,464.5) node  [font=\small,color={rgb, 255:red, 0; green, 0; blue, 0 }  ,opacity=1 ] [align=left] {\begin{minipage}[lt]{41.55pt}\setlength\topsep{0pt}
\cref{ch:problema-parada}
\end{minipage}};
% Text Node
\draw (749.5,465.5) node  [font=\small,color={rgb, 255:red, 0; green, 0; blue, 0 }  ,opacity=1 ] [align=left] {\begin{minipage}[lt]{41.55pt}\setlength\topsep{0pt}
\begin{flushright}
\cref{ch:sistemas-logicos}
\end{flushright}

\end{minipage}};
% Text Node
\draw (170.5,604.5) node  [font=\small,color={rgb, 255:red, 0; green, 0; blue, 0 }  ,opacity=1 ] [align=left] {\begin{minipage}[lt]{41.55pt}\setlength\topsep{0pt}
\cref{ch:teorema-incompletitud}
\end{minipage}};
% Text Node
\draw (643.76,195.55) node  [color={rgb, 255:red, 0; green, 0; blue, 0 }  ,opacity=1 ] [align=left] {\begin{minipage}[lt]{217.92pt}\setlength\topsep{0pt}
-- sólido {\small{(\cref{sec:solidez-completitud-decidibilidad})}}\\
-- decidible {\small{(\cref{sec:solidez-completitud-decidibilidad})}}\\
-- (semánticamente) completo {\small{(\cref{sec:solidez-completitud-decidibilidad})}}\\
-- sintácticamente completo {\small{(\cref{sec:consistencia})}}\\
-- sem./sint. consistente {\small{(\cref{sec:consistencia})}}
\end{minipage}};
% Text Node
\draw (414.84,645.25) node  [color={rgb, 255:red, 0; green, 0; blue, 0 }  ,opacity=1 ] [align=left] {\begin{minipage}[lt]{534.7pt}\setlength\topsep{0pt}
\begin{center}
Esquema resumen de este Trabajo de Fin de Grado
\end{center}
\end{minipage}};
% Text Node
\draw (380.15,407.8) node  [font=\small,color={rgb, 255:red, 0; green, 0; blue, 0 }  ,opacity=1 ] [align=left] {\begin{minipage}[lt]{108.93pt}\setlength\topsep{0pt}
\begin{center}
\cref{sec:primera-aproximacion}
\end{center}

\end{minipage}};

\end{tikzpicture}
\raisebox{-15.5cm}[0pt][0pt]{%
\resizebox{\linewidth}{!}{

\tikzset{every picture/.style={line width=0.75pt}} %set default line width to 0.75pt        

\begin{tikzpicture}[x=0.75pt,y=0.75pt,yscale=-1,xscale=1]
%uncomment if require: \path (0,694); %set diagram left start at 0, and has height of 694

%Straight Lines [id:da2056540562354816] 
\draw  [dash pattern={on 4.5pt off 4.5pt}]  (672,64.47) -- (810,64.47) ;
%Straight Lines [id:da23044041717189923] 
\draw  [dash pattern={on 4.5pt off 4.5pt}]  (412.67,64.47) -- (610,64.47) ;
\draw [shift={(410.67,64.47)}, rotate = 0] [color={rgb, 255:red, 0; green, 0; blue, 0 }  ][line width=0.75]    (10.93,-3.29) .. controls (6.95,-1.4) and (3.31,-0.3) .. (0,0) .. controls (3.31,0.3) and (6.95,1.4) .. (10.93,3.29)   ;
%Straight Lines [id:da8398435684999113] 
\draw  [dash pattern={on 4.5pt off 4.5pt}]  (810,64.47) -- (810,530) ;
%Rounded Rect [id:dp027018508739026448] 
\draw  [draw opacity=0][fill={rgb, 255:red, 230; green, 230; blue, 230 }  ,fill opacity=0.5 ] (20,439.33) .. controls (20,445.22) and (24.78,450) .. (30.67,450) -- (409.33,450) .. controls (415.22,450) and (420,445.22) .. (420,439.33) -- (420,140.67) .. controls (420,134.78) and (415.22,130) .. (409.33,130) -- (30.67,130) .. controls (24.78,130) and (20,134.78) .. (20,140.67) -- cycle ;
%Rounded Rect [id:dp9415679498487528] 
\draw  [dash pattern={on 4.5pt off 4.5pt}] (29.83,292.77) .. controls (29.83,285.88) and (35.41,280.3) .. (42.3,280.3) -- (398.54,280.3) .. controls (405.42,280.3) and (411,285.88) .. (411,292.77) -- (411,337.17) .. controls (411,344.05) and (405.42,349.64) .. (398.54,349.64) -- (42.3,349.64) .. controls (35.41,349.64) and (29.83,344.05) .. (29.83,337.17) -- cycle ;
%Rounded Rect [id:dp05100022295548956] 
\draw  [draw opacity=0][fill={rgb, 255:red, 230; green, 230; blue, 230 }  ,fill opacity=0.5 ] (130.5,575) .. controls (130.5,583.28) and (137.22,590) .. (145.5,590) -- (615,590) .. controls (623.28,590) and (630,583.28) .. (630,575) -- (630,395) .. controls (630,386.72) and (623.28,380) .. (615,380) -- (145.5,380) .. controls (137.22,380) and (130.5,386.72) .. (130.5,395) -- cycle ;
%Rounded Rect [id:dp6496809750851564] 
\draw  [draw opacity=0][fill={rgb, 255:red, 230; green, 230; blue, 230 }  ,fill opacity=0.5 ] (450.61,440.79) .. controls (450.61,445.82) and (454.69,449.9) .. (459.72,449.9) -- (780.89,449.9) .. controls (785.92,449.9) and (790,445.82) .. (790,440.79) -- (790,139.11) .. controls (790,134.08) and (785.92,130) .. (780.89,130) -- (459.72,130) .. controls (454.69,130) and (450.61,134.08) .. (450.61,139.11) -- cycle ;
%Rounded Rect [id:dp549636265291094] 
\draw  [fill={rgb, 255:red, 230; green, 230; blue, 230 }  ,fill opacity=0.5 ][dash pattern={on 4.5pt off 4.5pt}] (29.83,42.77) .. controls (29.83,35.88) and (35.41,30.3) .. (42.3,30.3) -- (398.54,30.3) .. controls (405.42,30.3) and (411,35.88) .. (411,42.77) -- (411,87.17) .. controls (411,94.05) and (405.42,99.64) .. (398.54,99.64) -- (42.3,99.64) .. controls (35.41,99.64) and (29.83,94.05) .. (29.83,87.17) -- cycle ;
%Rounded Rect [id:dp1378581119146769] 
\draw  [fill={rgb, 255:red, 255; green, 255; blue, 255 }  ,fill opacity=1 ] (40.4,49.82) .. controls (40.4,44.29) and (44.89,39.8) .. (50.42,39.8) -- (169.58,39.8) .. controls (175.11,39.8) and (179.6,44.29) .. (179.6,49.82) -- (179.6,79.88) .. controls (179.6,85.41) and (175.11,89.9) .. (169.58,89.9) -- (50.42,89.9) .. controls (44.89,89.9) and (40.4,85.41) .. (40.4,79.88) -- cycle ;
%Straight Lines [id:da8729153203483555] 
\draw    (188.42,63.77) -- (250.99,63.77)(188.42,66.77) -- (250.99,66.77) ;
\draw [shift={(257.99,65.27)}, rotate = 180] [color={rgb, 255:red, 0; green, 0; blue, 0 }  ][line width=0.75]    (10.93,-4.9) .. controls (6.95,-2.3) and (3.31,-0.67) .. (0,0) .. controls (3.31,0.67) and (6.95,2.3) .. (10.93,4.9)   ;
\draw [shift={(181.42,65.27)}, rotate = 0] [color={rgb, 255:red, 0; green, 0; blue, 0 }  ][line width=0.75]    (10.93,-4.9) .. controls (6.95,-2.3) and (3.31,-0.67) .. (0,0) .. controls (3.31,0.67) and (6.95,2.3) .. (10.93,4.9)   ;
%Straight Lines [id:da845243590120559] 
\draw    (220.17,99.97) -- (220.17,137.64) ;
\draw [shift={(220.17,139.64)}, rotate = 270] [color={rgb, 255:red, 0; green, 0; blue, 0 }  ][line width=0.75]    (10.93,-3.29) .. controls (6.95,-1.4) and (3.31,-0.3) .. (0,0) .. controls (3.31,0.3) and (6.95,1.4) .. (10.93,3.29)   ;
%Straight Lines [id:da7497001546502899] 
\draw    (380,420.33) -- (380,478) ;
\draw [shift={(380,480)}, rotate = 270] [color={rgb, 255:red, 0; green, 0; blue, 0 }  ][line width=0.75]    (10.93,-3.29) .. controls (6.95,-1.4) and (3.31,-0.3) .. (0,0) .. controls (3.31,0.3) and (6.95,1.4) .. (10.93,3.29)   ;
%Straight Lines [id:da5689893955334542] 
\draw    (300,420) -- (460,420) ;
%Rounded Rect [id:dp18582062333838767] 
\draw  [fill={rgb, 255:red, 255; green, 255; blue, 255 }  ,fill opacity=1 ] (260.4,49.82) .. controls (260.4,44.29) and (264.89,39.8) .. (270.42,39.8) -- (389.58,39.8) .. controls (395.11,39.8) and (399.6,44.29) .. (399.6,49.82) -- (399.6,79.88) .. controls (399.6,85.41) and (395.11,89.9) .. (389.58,89.9) -- (270.42,89.9) .. controls (264.89,89.9) and (260.4,85.41) .. (260.4,79.88) -- cycle ;
%Rounded Rect [id:dp9523169763067778] 
\draw  [fill={rgb, 255:red, 255; green, 255; blue, 255 }  ,fill opacity=1 ] (470,299.92) .. controls (470,294.39) and (474.49,289.9) .. (480.02,289.9) -- (599.18,289.9) .. controls (604.71,289.9) and (609.2,294.39) .. (609.2,299.92) -- (609.2,329.98) .. controls (609.2,335.51) and (604.71,340) .. (599.18,340) -- (480.02,340) .. controls (474.49,340) and (470,335.51) .. (470,329.98) -- cycle ;
%Rounded Rect [id:dp48544030815291705] 
\draw  [fill={rgb, 255:red, 255; green, 255; blue, 255 }  ,fill opacity=1 ] (220,490) .. controls (220,484.48) and (224.48,480) .. (230,480) -- (530,480) .. controls (535.52,480) and (540,484.48) .. (540,490) -- (540,570) .. controls (540,575.52) and (535.52,580) .. (530,580) -- (230,580) .. controls (224.48,580) and (220,575.52) .. (220,570) -- cycle ;
%Rounded Rect [id:dp08222627508455749] 
\draw  [fill={rgb, 255:red, 255; green, 255; blue, 255 }  ,fill opacity=1 ] (540.4,49.82) .. controls (540.4,44.29) and (544.89,39.8) .. (550.42,39.8) -- (690.65,39.8) .. controls (696.18,39.8) and (700.67,44.29) .. (700.67,49.82) -- (700.67,79.88) .. controls (700.67,85.41) and (696.18,89.9) .. (690.65,89.9) -- (550.42,89.9) .. controls (544.89,89.9) and (540.4,85.41) .. (540.4,79.88) -- cycle ;
%Straight Lines [id:da5868346364154593] 
\draw  [dash pattern={on 4.5pt off 4.5pt}]  (810,530) -- (542,530) ;
\draw [shift={(540,530)}, rotate = 360] [color={rgb, 255:red, 0; green, 0; blue, 0 }  ][line width=0.75]    (10.93,-3.29) .. controls (6.95,-1.4) and (3.31,-0.3) .. (0,0) .. controls (3.31,0.3) and (6.95,1.4) .. (10.93,3.29)   ;
%Rounded Rect [id:dp7392784471609453] 
\draw  [fill={rgb, 255:red, 255; green, 255; blue, 255 }  ,fill opacity=1 ] (30.33,150.47) .. controls (30.33,144.58) and (35.11,139.8) .. (41,139.8) -- (400.34,139.8) .. controls (406.23,139.8) and (411,144.58) .. (411,150.47) -- (411,229.34) .. controls (411,235.23) and (406.23,240.01) .. (400.34,240.01) -- (41,240.01) .. controls (35.11,240.01) and (30.33,235.23) .. (30.33,229.34) -- cycle ;
%Rounded Rect [id:dp90901637181607] 
\draw  [fill={rgb, 255:red, 255; green, 255; blue, 255 }  ,fill opacity=1 ] (41,189.82) .. controls (41,184.29) and (45.49,179.8) .. (51.02,179.8) -- (145.02,179.8) .. controls (150.55,179.8) and (155.04,184.29) .. (155.04,189.82) -- (155.04,219.88) .. controls (155.04,225.41) and (150.55,229.9) .. (145.02,229.9) -- (51.02,229.9) .. controls (45.49,229.9) and (41,225.41) .. (41,219.88) -- cycle ;
%Rounded Rect [id:dp06115018179625542] 
\draw  [fill={rgb, 255:red, 255; green, 255; blue, 255 }  ,fill opacity=1 ] (163.65,189.82) .. controls (163.65,184.29) and (168.13,179.8) .. (173.67,179.8) -- (267.67,179.8) .. controls (273.2,179.8) and (277.69,184.29) .. (277.69,189.82) -- (277.69,219.88) .. controls (277.69,225.41) and (273.2,229.9) .. (267.67,229.9) -- (173.67,229.9) .. controls (168.13,229.9) and (163.65,225.41) .. (163.65,219.88) -- cycle ;
%Rounded Rect [id:dp6195275548153432] 
\draw  [fill={rgb, 255:red, 255; green, 255; blue, 255 }  ,fill opacity=1 ] (286.3,189.82) .. controls (286.3,184.29) and (290.78,179.8) .. (296.32,179.8) -- (390.31,179.8) .. controls (395.85,179.8) and (400.33,184.29) .. (400.33,189.82) -- (400.33,219.88) .. controls (400.33,225.41) and (395.85,229.9) .. (390.31,229.9) -- (296.32,229.9) .. controls (290.78,229.9) and (286.3,225.41) .. (286.3,219.88) -- cycle ;
%Rounded Rect [id:dp41283415100159493] 
\draw  [fill={rgb, 255:red, 255; green, 255; blue, 255 }  ,fill opacity=1 ] (40.8,299.92) .. controls (40.8,294.39) and (45.29,289.9) .. (50.82,289.9) -- (169.98,289.9) .. controls (175.51,289.9) and (180,294.39) .. (180,299.92) -- (180,329.98) .. controls (180,335.51) and (175.51,340) .. (169.98,340) -- (50.82,340) .. controls (45.29,340) and (40.8,335.51) .. (40.8,329.98) -- cycle ;
%Rounded Rect [id:dp40154854661776396] 
\draw  [fill={rgb, 255:red, 255; green, 255; blue, 255 }  ,fill opacity=1 ] (260.4,299.82) .. controls (260.4,294.29) and (264.89,289.8) .. (270.42,289.8) -- (389.58,289.8) .. controls (395.11,289.8) and (399.6,294.29) .. (399.6,299.82) -- (399.6,329.88) .. controls (399.6,335.41) and (395.11,339.9) .. (389.58,339.9) -- (270.42,339.9) .. controls (264.89,339.9) and (260.4,335.41) .. (260.4,329.88) -- cycle ;
%Straight Lines [id:da6602630444249074] 
\draw    (180.17,315.8) -- (260,315.8) ;
%Straight Lines [id:da30905866347109745] 
\draw    (219.86,315.97) -- (219.86,386.81) ;
\draw [shift={(219.86,388.81)}, rotate = 270] [color={rgb, 255:red, 0; green, 0; blue, 0 }  ][line width=0.75]    (10.93,-3.29) .. controls (6.95,-1.4) and (3.31,-0.3) .. (0,0) .. controls (3.31,0.3) and (6.95,1.4) .. (10.93,3.29)   ;
%Rounded Rect [id:dp44627062039563214] 
\draw  [fill={rgb, 255:red, 255; green, 255; blue, 255 }  ,fill opacity=1 ] (139.67,399.82) .. controls (139.67,394.29) and (144.15,389.8) .. (149.69,389.8) -- (290.31,389.8) .. controls (295.85,389.8) and (300.33,394.29) .. (300.33,399.82) -- (300.33,429.88) .. controls (300.33,435.41) and (295.85,439.9) .. (290.31,439.9) -- (149.69,439.9) .. controls (144.15,439.9) and (139.67,435.41) .. (139.67,429.88) -- cycle ;
%Straight Lines [id:da8375870493480149] 
\draw    (220,240) -- (220,277.67) ;
\draw [shift={(220,279.67)}, rotate = 270] [color={rgb, 255:red, 0; green, 0; blue, 0 }  ][line width=0.75]    (10.93,-3.29) .. controls (6.95,-1.4) and (3.31,-0.3) .. (0,0) .. controls (3.31,0.3) and (6.95,1.4) .. (10.93,3.29)   ;
%Rounded Rect [id:dp37285624114093485] 
\draw  [fill={rgb, 255:red, 255; green, 255; blue, 255 }  ,fill opacity=1 ] (460,400.02) .. controls (460,394.49) and (464.49,390) .. (470.02,390) -- (610.65,390) .. controls (616.18,390) and (620.67,394.49) .. (620.67,400.02) -- (620.67,430.08) .. controls (620.67,435.61) and (616.18,440.1) .. (610.65,440.1) -- (470.02,440.1) .. controls (464.49,440.1) and (460,435.61) .. (460,430.08) -- cycle ;
%Rounded Rect [id:dp6883616335272991] 
\draw  [fill={rgb, 255:red, 255; green, 255; blue, 255 }  ,fill opacity=1 ] (629.2,299.92) .. controls (629.2,294.39) and (633.69,289.9) .. (639.22,289.9) -- (758.38,289.9) .. controls (763.91,289.9) and (768.4,294.39) .. (768.4,299.92) -- (768.4,329.98) .. controls (768.4,335.51) and (763.91,340) .. (758.38,340) -- (639.22,340) .. controls (633.69,340) and (629.2,335.51) .. (629.2,329.98) -- cycle ;
%Straight Lines [id:da19870419404011752] 
\draw    (540,340) -- (540,388) ;
\draw [shift={(540,390)}, rotate = 270] [color={rgb, 255:red, 0; green, 0; blue, 0 }  ][line width=0.75]    (10.93,-3.29) .. controls (6.95,-1.4) and (3.31,-0.3) .. (0,0) .. controls (3.31,0.3) and (6.95,1.4) .. (10.93,3.29)   ;
%Rounded Rect [id:dp9692784599420237] 
\draw  [fill={rgb, 255:red, 255; green, 255; blue, 255 }  ,fill opacity=1 ] (230,529.92) .. controls (230,524.39) and (234.49,519.9) .. (240.02,519.9) -- (359.18,519.9) .. controls (364.71,519.9) and (369.2,524.39) .. (369.2,529.92) -- (369.2,559.98) .. controls (369.2,565.51) and (364.71,570) .. (359.18,570) -- (240.02,570) .. controls (234.49,570) and (230,565.51) .. (230,559.98) -- cycle ;
%Rounded Rect [id:dp749684165723072] 
\draw  [fill={rgb, 255:red, 255; green, 255; blue, 255 }  ,fill opacity=1 ] (390,530.02) .. controls (390,524.49) and (394.49,520) .. (400.02,520) -- (519.18,520) .. controls (524.71,520) and (529.2,524.49) .. (529.2,530.02) -- (529.2,560.08) .. controls (529.2,565.61) and (524.71,570.1) .. (519.18,570.1) -- (400.02,570.1) .. controls (394.49,570.1) and (390,565.61) .. (390,560.08) -- cycle ;
%Rounded Rect [id:dp28001422897762485] 
\draw  [fill={rgb, 255:red, 255; green, 255; blue, 255 }  ,fill opacity=1 ][dash pattern={on 4.5pt off 4.5pt}] (460,148.92) .. controls (460,143.99) and (463.99,140) .. (468.92,140) -- (771.08,140) .. controls (776.01,140) and (780,143.99) .. (780,148.92) -- (780,241.55) .. controls (780,246.48) and (776.01,250.47) .. (771.08,250.47) -- (468.92,250.47) .. controls (463.99,250.47) and (460,246.48) .. (460,241.55) -- cycle ;
%Shape: Brace [id:dp9411387904179627] 
\draw   (774.67,286.8) .. controls (774.67,282.13) and (772.34,279.8) .. (767.67,279.8) -- (629.96,279.8) .. controls (623.29,279.8) and (619.96,277.47) .. (619.96,272.8) .. controls (619.96,277.47) and (616.63,279.8) .. (609.96,279.8)(612.96,279.8) -- (472.25,279.8) .. controls (467.58,279.8) and (465.25,282.13) .. (465.25,286.8) ;
%Straight Lines [id:da9902399507274253] 
\draw    (620,273) -- (620,253.14) ;
\draw [shift={(620,251.14)}, rotate = 90] [color={rgb, 255:red, 0; green, 0; blue, 0 }  ][line width=0.75]    (10.93,-3.29) .. controls (6.95,-1.4) and (3.31,-0.3) .. (0,0) .. controls (3.31,0.3) and (6.95,1.4) .. (10.93,3.29)   ;

% Text Node
\draw (110,64.85) node  [color={rgb, 255:red, 0; green, 0; blue, 0 }  ,opacity=1 ] [align=left] {\begin{minipage}[lt]{92.93pt}\setlength\topsep{0pt}
\begin{center}
programa en Python\\
\small{(\cref{sec:programas-python})}
\end{center}

\end{minipage}};
% Text Node
\draw (219.89,51.97) node  [font=\small,color={rgb, 255:red, 0; green, 0; blue, 0 }  ,opacity=1 ] [align=left] {\begin{minipage}[lt]{54.86pt}\setlength\topsep{0pt}
\begin{center}
\Cref{sec:equivalencia}
\end{center}

\end{minipage}};
% Text Node
\draw (69.5,15.5) node  [font=\small,color={rgb, 255:red, 0; green, 0; blue, 0 }  ,opacity=1 ] [align=left] {\begin{minipage}[lt]{41.55pt}\setlength\topsep{0pt}
\cref{ch:programas-python-maquinas-turing}
\end{minipage}};
% Text Node
\draw (330,64.85) node  [color={rgb, 255:red, 0; green, 0; blue, 0 }  ,opacity=1 ] [align=left] {\begin{minipage}[lt]{92.93pt}\setlength\topsep{0pt}
\begin{center}
máquina de Turing\\
\small{(\cref{sec:maquinas-turing})}
\end{center}

\end{minipage}};
% Text Node
\draw (539.6,314.95) node  [color={rgb, 255:red, 0; green, 0; blue, 0 }  ,opacity=1 ] [align=left] {\begin{minipage}[lt]{92.93pt}\setlength\topsep{0pt}
\begin{center}
sistema lógico\\
\small{(\cref{sec:sistemas-logicos})}
\end{center}

\end{minipage}};
% Text Node
\draw (380.5,500.11) node  [color={rgb, 255:red, 0; green, 0; blue, 0 }  ,opacity=1 ] [align=left] {\begin{minipage}[lt]{159.55pt}\setlength\topsep{0pt}
\begin{center}
\textbf{Primer Teorema de Incompletitud}
\end{center}

\end{minipage}};
% Text Node
\draw (620.53,64.85) node  [color={rgb, 255:red, 0; green, 0; blue, 0 }  ,opacity=1 ] [align=left] {\begin{minipage}[lt]{106.99pt}\setlength\topsep{0pt}
\begin{center}
tesis de Church-Turing\\
\small{(\cref{sec:church-turing})}
\end{center}

\end{minipage}};
% Text Node
\draw (480.28,52.8) node  [font=\small,color={rgb, 255:red, 0; green, 0; blue, 0 }  ,opacity=1 ] [align=left] {\begin{minipage}[lt]{81.9pt}\setlength\topsep{0pt}
\begin{center}
\cref{sec:church-turing}
\end{center}

\end{minipage}};
% Text Node
\draw (589.5,518.3) node  [font=\small,color={rgb, 255:red, 0; green, 0; blue, 0 }  ,opacity=1 ] [align=left] {\begin{minipage}[lt]{54.86pt}\setlength\topsep{0pt}
\begin{center}
\Cref{subsec:mente-humana-tesis-church-turing}
\end{center}

\end{minipage}};
% Text Node
\draw (220.2,159.62) node  [color={rgb, 255:red, 0; green, 0; blue, 0 }  ,opacity=1 ] [align=left] {\begin{minipage}[lt]{230.11pt}\setlength\topsep{0pt}
\begin{center}
problema computacional
\end{center}

\end{minipage}};
% Text Node
\draw (98.02,204.85) node  [color={rgb, 255:red, 0; green, 0; blue, 0 }  ,opacity=1 ] [align=left] {\begin{minipage}[lt]{76.13pt}\setlength\topsep{0pt}
\begin{center}
decidible\\
\small{(\cref{sec:problemas-decidibles})}
\end{center}

\end{minipage}};
% Text Node
\draw (220.67,204.85) node  [color={rgb, 255:red, 0; green, 0; blue, 0 }  ,opacity=1 ] [align=left] {\begin{minipage}[lt]{76.13pt}\setlength\topsep{0pt}
\begin{center}
no decidible\\
\small{(\cref{sec:problemas-no-decidibles})}
\end{center}

\end{minipage}};
% Text Node
\draw (343.31,204.85) node  [color={rgb, 255:red, 0; green, 0; blue, 0 }  ,opacity=1 ] [align=left] {\begin{minipage}[lt]{76.13pt}\setlength\topsep{0pt}
\begin{center}
semidecidible\\
\small{(\cref{sec:problemas-semidecidibles})}
\end{center}

\end{minipage}};
% Text Node
\draw (110.4,314.95) node  [color={rgb, 255:red, 0; green, 0; blue, 0 }  ,opacity=1 ] [align=left] {\begin{minipage}[lt]{92.93pt}\setlength\topsep{0pt}
\begin{center}
universalidad\\
\small{(\cref{sec:universalidad})}
\end{center}

\end{minipage}};
% Text Node
\draw (330,314.85) node  [color={rgb, 255:red, 0; green, 0; blue, 0 }  ,opacity=1 ] [align=left] {\begin{minipage}[lt]{92.93pt}\setlength\topsep{0pt}
\begin{center}
reducción\\
\small{(\cref{sec:reducciones})}
\end{center}

\end{minipage}};
% Text Node
\draw (220,414.85) node  [color={rgb, 255:red, 0; green, 0; blue, 0 }  ,opacity=1 ] [align=left] {\begin{minipage}[lt]{107.26pt}\setlength\topsep{0pt}
\begin{center}
problema de la parada\\
\small{(\cref{sec:problema-parada})}
\end{center}

\end{minipage}};
% Text Node
\draw (540.33,415.05) node  [color={rgb, 255:red, 0; green, 0; blue, 0 }  ,opacity=1 ] [align=left] {\begin{minipage}[lt]{107.26pt}\setlength\topsep{0pt}
\begin{center}
aritmética de Peano\\
\small{(\cref{sec:aritmetica-peano})}
\end{center}

\end{minipage}};
% Text Node
\draw (698.8,314.95) node  [color={rgb, 255:red, 0; green, 0; blue, 0 }  ,opacity=1 ] [align=left] {\begin{minipage}[lt]{92.93pt}\setlength\topsep{0pt}
\begin{center}
sistema formal\\
\small{(\cref{sec:sistemas-formales})}
\end{center}

\end{minipage}};
% Text Node
\draw (299.6,544.95) node  [color={rgb, 255:red, 0; green, 0; blue, 0 }  ,opacity=1 ] [align=left] {\begin{minipage}[lt]{92.93pt}\setlength\topsep{0pt}
\begin{center}
versión semántica\\
\small{(\cref{sec:version-semantica})}
\end{center}

\end{minipage}};
% Text Node
\draw (459.6,545.05) node  [color={rgb, 255:red, 0; green, 0; blue, 0 }  ,opacity=1 ] [align=left] {\begin{minipage}[lt]{92.93pt}\setlength\topsep{0pt}
\begin{center}
versión sintáctica\\
\small{(\cref{sec:version-sintactica})}
\end{center}

\end{minipage}};
% Text Node
\draw (60.5,464.5) node  [font=\small,color={rgb, 255:red, 0; green, 0; blue, 0 }  ,opacity=1 ] [align=left] {\begin{minipage}[lt]{41.55pt}\setlength\topsep{0pt}
\cref{ch:problema-parada}
\end{minipage}};
% Text Node
\draw (749.5,465.5) node  [font=\small,color={rgb, 255:red, 0; green, 0; blue, 0 }  ,opacity=1 ] [align=left] {\begin{minipage}[lt]{41.55pt}\setlength\topsep{0pt}
\begin{flushright}
\cref{ch:sistemas-logicos}
\end{flushright}

\end{minipage}};
% Text Node
\draw (170.5,604.5) node  [font=\small,color={rgb, 255:red, 0; green, 0; blue, 0 }  ,opacity=1 ] [align=left] {\begin{minipage}[lt]{41.55pt}\setlength\topsep{0pt}
\cref{ch:teorema-incompletitud}
\end{minipage}};
% Text Node
\draw (643.76,195.55) node  [color={rgb, 255:red, 0; green, 0; blue, 0 }  ,opacity=1 ] [align=left] {\begin{minipage}[lt]{217.92pt}\setlength\topsep{0pt}
-- sólido {\small{(\cref{sec:solidez-completitud-decidibilidad})}}\\
-- decidible {\small{(\cref{sec:solidez-completitud-decidibilidad})}}\\
-- (semánticamente) completo {\small{(\cref{sec:solidez-completitud-decidibilidad})}}\\
-- sintácticamente completo {\small{(\cref{sec:consistencia})}}\\
-- sem./sint. consistente {\small{(\cref{sec:consistencia})}}
\end{minipage}};
% Text Node
\draw (414.84,645.25) node  [color={rgb, 255:red, 0; green, 0; blue, 0 }  ,opacity=1 ] [align=left] {\begin{minipage}[lt]{534.7pt}\setlength\topsep{0pt}
\begin{center}
Esquema resumen de este Trabajo de Fin de Grado
\end{center}
\end{minipage}};
% Text Node
\draw (380.15,407.8) node  [font=\small,color={rgb, 255:red, 0; green, 0; blue, 0 }  ,opacity=1 ] [align=left] {\begin{minipage}[lt]{108.93pt}\setlength\topsep{0pt}
\begin{center}
\cref{sec:primera-aproximacion}
\end{center}

\end{minipage}};

\end{tikzpicture}}
}
%\caption{Esquema resumen (extendido)}

%\label{fig:resumen-extendido}
\end{figure}
% ====================
\end{landscape}

\endinput
}
%\caption{Esquema resumen}
%\label{fig:resumen}
%\end{figure}
% ====================
%\vfill

\begin{landscape}
% ====================
\begin{figure}[H]
%\captionsetup{margin=200cm}

\centering
%\vspace{-1.5cm}
%\hspace{-1cm}
%

\tikzset{every picture/.style={line width=0.75pt}} %set default line width to 0.75pt        

\begin{tikzpicture}[x=0.75pt,y=0.75pt,yscale=-1,xscale=1]
%uncomment if require: \path (0,694); %set diagram left start at 0, and has height of 694

%Straight Lines [id:da2056540562354816] 
\draw  [dash pattern={on 4.5pt off 4.5pt}]  (672,64.47) -- (810,64.47) ;
%Straight Lines [id:da23044041717189923] 
\draw  [dash pattern={on 4.5pt off 4.5pt}]  (412.67,64.47) -- (610,64.47) ;
\draw [shift={(410.67,64.47)}, rotate = 0] [color={rgb, 255:red, 0; green, 0; blue, 0 }  ][line width=0.75]    (10.93,-3.29) .. controls (6.95,-1.4) and (3.31,-0.3) .. (0,0) .. controls (3.31,0.3) and (6.95,1.4) .. (10.93,3.29)   ;
%Straight Lines [id:da8398435684999113] 
\draw  [dash pattern={on 4.5pt off 4.5pt}]  (810,64.47) -- (810,530) ;
%Rounded Rect [id:dp027018508739026448] 
\draw  [draw opacity=0][fill={rgb, 255:red, 230; green, 230; blue, 230 }  ,fill opacity=0.5 ] (20,439.33) .. controls (20,445.22) and (24.78,450) .. (30.67,450) -- (409.33,450) .. controls (415.22,450) and (420,445.22) .. (420,439.33) -- (420,140.67) .. controls (420,134.78) and (415.22,130) .. (409.33,130) -- (30.67,130) .. controls (24.78,130) and (20,134.78) .. (20,140.67) -- cycle ;
%Rounded Rect [id:dp9415679498487528] 
\draw  [dash pattern={on 4.5pt off 4.5pt}] (29.83,292.77) .. controls (29.83,285.88) and (35.41,280.3) .. (42.3,280.3) -- (398.54,280.3) .. controls (405.42,280.3) and (411,285.88) .. (411,292.77) -- (411,337.17) .. controls (411,344.05) and (405.42,349.64) .. (398.54,349.64) -- (42.3,349.64) .. controls (35.41,349.64) and (29.83,344.05) .. (29.83,337.17) -- cycle ;
%Rounded Rect [id:dp05100022295548956] 
\draw  [draw opacity=0][fill={rgb, 255:red, 230; green, 230; blue, 230 }  ,fill opacity=0.5 ] (130.5,575) .. controls (130.5,583.28) and (137.22,590) .. (145.5,590) -- (615,590) .. controls (623.28,590) and (630,583.28) .. (630,575) -- (630,395) .. controls (630,386.72) and (623.28,380) .. (615,380) -- (145.5,380) .. controls (137.22,380) and (130.5,386.72) .. (130.5,395) -- cycle ;
%Rounded Rect [id:dp6496809750851564] 
\draw  [draw opacity=0][fill={rgb, 255:red, 230; green, 230; blue, 230 }  ,fill opacity=0.5 ] (450.61,440.79) .. controls (450.61,445.82) and (454.69,449.9) .. (459.72,449.9) -- (780.89,449.9) .. controls (785.92,449.9) and (790,445.82) .. (790,440.79) -- (790,139.11) .. controls (790,134.08) and (785.92,130) .. (780.89,130) -- (459.72,130) .. controls (454.69,130) and (450.61,134.08) .. (450.61,139.11) -- cycle ;
%Rounded Rect [id:dp549636265291094] 
\draw  [fill={rgb, 255:red, 230; green, 230; blue, 230 }  ,fill opacity=0.5 ][dash pattern={on 4.5pt off 4.5pt}] (29.83,42.77) .. controls (29.83,35.88) and (35.41,30.3) .. (42.3,30.3) -- (398.54,30.3) .. controls (405.42,30.3) and (411,35.88) .. (411,42.77) -- (411,87.17) .. controls (411,94.05) and (405.42,99.64) .. (398.54,99.64) -- (42.3,99.64) .. controls (35.41,99.64) and (29.83,94.05) .. (29.83,87.17) -- cycle ;
%Rounded Rect [id:dp1378581119146769] 
\draw  [fill={rgb, 255:red, 255; green, 255; blue, 255 }  ,fill opacity=1 ] (40.4,49.82) .. controls (40.4,44.29) and (44.89,39.8) .. (50.42,39.8) -- (169.58,39.8) .. controls (175.11,39.8) and (179.6,44.29) .. (179.6,49.82) -- (179.6,79.88) .. controls (179.6,85.41) and (175.11,89.9) .. (169.58,89.9) -- (50.42,89.9) .. controls (44.89,89.9) and (40.4,85.41) .. (40.4,79.88) -- cycle ;
%Straight Lines [id:da8729153203483555] 
\draw    (188.42,63.77) -- (250.99,63.77)(188.42,66.77) -- (250.99,66.77) ;
\draw [shift={(257.99,65.27)}, rotate = 180] [color={rgb, 255:red, 0; green, 0; blue, 0 }  ][line width=0.75]    (10.93,-4.9) .. controls (6.95,-2.3) and (3.31,-0.67) .. (0,0) .. controls (3.31,0.67) and (6.95,2.3) .. (10.93,4.9)   ;
\draw [shift={(181.42,65.27)}, rotate = 0] [color={rgb, 255:red, 0; green, 0; blue, 0 }  ][line width=0.75]    (10.93,-4.9) .. controls (6.95,-2.3) and (3.31,-0.67) .. (0,0) .. controls (3.31,0.67) and (6.95,2.3) .. (10.93,4.9)   ;
%Straight Lines [id:da845243590120559] 
\draw    (220.17,99.97) -- (220.17,137.64) ;
\draw [shift={(220.17,139.64)}, rotate = 270] [color={rgb, 255:red, 0; green, 0; blue, 0 }  ][line width=0.75]    (10.93,-3.29) .. controls (6.95,-1.4) and (3.31,-0.3) .. (0,0) .. controls (3.31,0.3) and (6.95,1.4) .. (10.93,3.29)   ;
%Straight Lines [id:da7497001546502899] 
\draw    (380,420.33) -- (380,478) ;
\draw [shift={(380,480)}, rotate = 270] [color={rgb, 255:red, 0; green, 0; blue, 0 }  ][line width=0.75]    (10.93,-3.29) .. controls (6.95,-1.4) and (3.31,-0.3) .. (0,0) .. controls (3.31,0.3) and (6.95,1.4) .. (10.93,3.29)   ;
%Straight Lines [id:da5689893955334542] 
\draw    (300,420) -- (460,420) ;
%Rounded Rect [id:dp18582062333838767] 
\draw  [fill={rgb, 255:red, 255; green, 255; blue, 255 }  ,fill opacity=1 ] (260.4,49.82) .. controls (260.4,44.29) and (264.89,39.8) .. (270.42,39.8) -- (389.58,39.8) .. controls (395.11,39.8) and (399.6,44.29) .. (399.6,49.82) -- (399.6,79.88) .. controls (399.6,85.41) and (395.11,89.9) .. (389.58,89.9) -- (270.42,89.9) .. controls (264.89,89.9) and (260.4,85.41) .. (260.4,79.88) -- cycle ;
%Rounded Rect [id:dp9523169763067778] 
\draw  [fill={rgb, 255:red, 255; green, 255; blue, 255 }  ,fill opacity=1 ] (470,299.92) .. controls (470,294.39) and (474.49,289.9) .. (480.02,289.9) -- (599.18,289.9) .. controls (604.71,289.9) and (609.2,294.39) .. (609.2,299.92) -- (609.2,329.98) .. controls (609.2,335.51) and (604.71,340) .. (599.18,340) -- (480.02,340) .. controls (474.49,340) and (470,335.51) .. (470,329.98) -- cycle ;
%Rounded Rect [id:dp48544030815291705] 
\draw  [fill={rgb, 255:red, 255; green, 255; blue, 255 }  ,fill opacity=1 ] (220,490) .. controls (220,484.48) and (224.48,480) .. (230,480) -- (530,480) .. controls (535.52,480) and (540,484.48) .. (540,490) -- (540,570) .. controls (540,575.52) and (535.52,580) .. (530,580) -- (230,580) .. controls (224.48,580) and (220,575.52) .. (220,570) -- cycle ;
%Rounded Rect [id:dp08222627508455749] 
\draw  [fill={rgb, 255:red, 255; green, 255; blue, 255 }  ,fill opacity=1 ] (540.4,49.82) .. controls (540.4,44.29) and (544.89,39.8) .. (550.42,39.8) -- (690.65,39.8) .. controls (696.18,39.8) and (700.67,44.29) .. (700.67,49.82) -- (700.67,79.88) .. controls (700.67,85.41) and (696.18,89.9) .. (690.65,89.9) -- (550.42,89.9) .. controls (544.89,89.9) and (540.4,85.41) .. (540.4,79.88) -- cycle ;
%Straight Lines [id:da5868346364154593] 
\draw  [dash pattern={on 4.5pt off 4.5pt}]  (810,530) -- (542,530) ;
\draw [shift={(540,530)}, rotate = 360] [color={rgb, 255:red, 0; green, 0; blue, 0 }  ][line width=0.75]    (10.93,-3.29) .. controls (6.95,-1.4) and (3.31,-0.3) .. (0,0) .. controls (3.31,0.3) and (6.95,1.4) .. (10.93,3.29)   ;
%Rounded Rect [id:dp7392784471609453] 
\draw  [fill={rgb, 255:red, 255; green, 255; blue, 255 }  ,fill opacity=1 ] (30.33,150.47) .. controls (30.33,144.58) and (35.11,139.8) .. (41,139.8) -- (400.34,139.8) .. controls (406.23,139.8) and (411,144.58) .. (411,150.47) -- (411,229.34) .. controls (411,235.23) and (406.23,240.01) .. (400.34,240.01) -- (41,240.01) .. controls (35.11,240.01) and (30.33,235.23) .. (30.33,229.34) -- cycle ;
%Rounded Rect [id:dp90901637181607] 
\draw  [fill={rgb, 255:red, 255; green, 255; blue, 255 }  ,fill opacity=1 ] (41,189.82) .. controls (41,184.29) and (45.49,179.8) .. (51.02,179.8) -- (145.02,179.8) .. controls (150.55,179.8) and (155.04,184.29) .. (155.04,189.82) -- (155.04,219.88) .. controls (155.04,225.41) and (150.55,229.9) .. (145.02,229.9) -- (51.02,229.9) .. controls (45.49,229.9) and (41,225.41) .. (41,219.88) -- cycle ;
%Rounded Rect [id:dp06115018179625542] 
\draw  [fill={rgb, 255:red, 255; green, 255; blue, 255 }  ,fill opacity=1 ] (163.65,189.82) .. controls (163.65,184.29) and (168.13,179.8) .. (173.67,179.8) -- (267.67,179.8) .. controls (273.2,179.8) and (277.69,184.29) .. (277.69,189.82) -- (277.69,219.88) .. controls (277.69,225.41) and (273.2,229.9) .. (267.67,229.9) -- (173.67,229.9) .. controls (168.13,229.9) and (163.65,225.41) .. (163.65,219.88) -- cycle ;
%Rounded Rect [id:dp6195275548153432] 
\draw  [fill={rgb, 255:red, 255; green, 255; blue, 255 }  ,fill opacity=1 ] (286.3,189.82) .. controls (286.3,184.29) and (290.78,179.8) .. (296.32,179.8) -- (390.31,179.8) .. controls (395.85,179.8) and (400.33,184.29) .. (400.33,189.82) -- (400.33,219.88) .. controls (400.33,225.41) and (395.85,229.9) .. (390.31,229.9) -- (296.32,229.9) .. controls (290.78,229.9) and (286.3,225.41) .. (286.3,219.88) -- cycle ;
%Rounded Rect [id:dp41283415100159493] 
\draw  [fill={rgb, 255:red, 255; green, 255; blue, 255 }  ,fill opacity=1 ] (40.8,299.92) .. controls (40.8,294.39) and (45.29,289.9) .. (50.82,289.9) -- (169.98,289.9) .. controls (175.51,289.9) and (180,294.39) .. (180,299.92) -- (180,329.98) .. controls (180,335.51) and (175.51,340) .. (169.98,340) -- (50.82,340) .. controls (45.29,340) and (40.8,335.51) .. (40.8,329.98) -- cycle ;
%Rounded Rect [id:dp40154854661776396] 
\draw  [fill={rgb, 255:red, 255; green, 255; blue, 255 }  ,fill opacity=1 ] (260.4,299.82) .. controls (260.4,294.29) and (264.89,289.8) .. (270.42,289.8) -- (389.58,289.8) .. controls (395.11,289.8) and (399.6,294.29) .. (399.6,299.82) -- (399.6,329.88) .. controls (399.6,335.41) and (395.11,339.9) .. (389.58,339.9) -- (270.42,339.9) .. controls (264.89,339.9) and (260.4,335.41) .. (260.4,329.88) -- cycle ;
%Straight Lines [id:da6602630444249074] 
\draw    (180.17,315.8) -- (260,315.8) ;
%Straight Lines [id:da30905866347109745] 
\draw    (219.86,315.97) -- (219.86,386.81) ;
\draw [shift={(219.86,388.81)}, rotate = 270] [color={rgb, 255:red, 0; green, 0; blue, 0 }  ][line width=0.75]    (10.93,-3.29) .. controls (6.95,-1.4) and (3.31,-0.3) .. (0,0) .. controls (3.31,0.3) and (6.95,1.4) .. (10.93,3.29)   ;
%Rounded Rect [id:dp44627062039563214] 
\draw  [fill={rgb, 255:red, 255; green, 255; blue, 255 }  ,fill opacity=1 ] (139.67,399.82) .. controls (139.67,394.29) and (144.15,389.8) .. (149.69,389.8) -- (290.31,389.8) .. controls (295.85,389.8) and (300.33,394.29) .. (300.33,399.82) -- (300.33,429.88) .. controls (300.33,435.41) and (295.85,439.9) .. (290.31,439.9) -- (149.69,439.9) .. controls (144.15,439.9) and (139.67,435.41) .. (139.67,429.88) -- cycle ;
%Straight Lines [id:da8375870493480149] 
\draw    (220,240) -- (220,277.67) ;
\draw [shift={(220,279.67)}, rotate = 270] [color={rgb, 255:red, 0; green, 0; blue, 0 }  ][line width=0.75]    (10.93,-3.29) .. controls (6.95,-1.4) and (3.31,-0.3) .. (0,0) .. controls (3.31,0.3) and (6.95,1.4) .. (10.93,3.29)   ;
%Rounded Rect [id:dp37285624114093485] 
\draw  [fill={rgb, 255:red, 255; green, 255; blue, 255 }  ,fill opacity=1 ] (460,400.02) .. controls (460,394.49) and (464.49,390) .. (470.02,390) -- (610.65,390) .. controls (616.18,390) and (620.67,394.49) .. (620.67,400.02) -- (620.67,430.08) .. controls (620.67,435.61) and (616.18,440.1) .. (610.65,440.1) -- (470.02,440.1) .. controls (464.49,440.1) and (460,435.61) .. (460,430.08) -- cycle ;
%Rounded Rect [id:dp6883616335272991] 
\draw  [fill={rgb, 255:red, 255; green, 255; blue, 255 }  ,fill opacity=1 ] (629.2,299.92) .. controls (629.2,294.39) and (633.69,289.9) .. (639.22,289.9) -- (758.38,289.9) .. controls (763.91,289.9) and (768.4,294.39) .. (768.4,299.92) -- (768.4,329.98) .. controls (768.4,335.51) and (763.91,340) .. (758.38,340) -- (639.22,340) .. controls (633.69,340) and (629.2,335.51) .. (629.2,329.98) -- cycle ;
%Straight Lines [id:da19870419404011752] 
\draw    (540,340) -- (540,388) ;
\draw [shift={(540,390)}, rotate = 270] [color={rgb, 255:red, 0; green, 0; blue, 0 }  ][line width=0.75]    (10.93,-3.29) .. controls (6.95,-1.4) and (3.31,-0.3) .. (0,0) .. controls (3.31,0.3) and (6.95,1.4) .. (10.93,3.29)   ;
%Rounded Rect [id:dp9692784599420237] 
\draw  [fill={rgb, 255:red, 255; green, 255; blue, 255 }  ,fill opacity=1 ] (230,529.92) .. controls (230,524.39) and (234.49,519.9) .. (240.02,519.9) -- (359.18,519.9) .. controls (364.71,519.9) and (369.2,524.39) .. (369.2,529.92) -- (369.2,559.98) .. controls (369.2,565.51) and (364.71,570) .. (359.18,570) -- (240.02,570) .. controls (234.49,570) and (230,565.51) .. (230,559.98) -- cycle ;
%Rounded Rect [id:dp749684165723072] 
\draw  [fill={rgb, 255:red, 255; green, 255; blue, 255 }  ,fill opacity=1 ] (390,530.02) .. controls (390,524.49) and (394.49,520) .. (400.02,520) -- (519.18,520) .. controls (524.71,520) and (529.2,524.49) .. (529.2,530.02) -- (529.2,560.08) .. controls (529.2,565.61) and (524.71,570.1) .. (519.18,570.1) -- (400.02,570.1) .. controls (394.49,570.1) and (390,565.61) .. (390,560.08) -- cycle ;
%Rounded Rect [id:dp28001422897762485] 
\draw  [fill={rgb, 255:red, 255; green, 255; blue, 255 }  ,fill opacity=1 ][dash pattern={on 4.5pt off 4.5pt}] (460,148.92) .. controls (460,143.99) and (463.99,140) .. (468.92,140) -- (771.08,140) .. controls (776.01,140) and (780,143.99) .. (780,148.92) -- (780,241.55) .. controls (780,246.48) and (776.01,250.47) .. (771.08,250.47) -- (468.92,250.47) .. controls (463.99,250.47) and (460,246.48) .. (460,241.55) -- cycle ;
%Shape: Brace [id:dp9411387904179627] 
\draw   (774.67,286.8) .. controls (774.67,282.13) and (772.34,279.8) .. (767.67,279.8) -- (629.96,279.8) .. controls (623.29,279.8) and (619.96,277.47) .. (619.96,272.8) .. controls (619.96,277.47) and (616.63,279.8) .. (609.96,279.8)(612.96,279.8) -- (472.25,279.8) .. controls (467.58,279.8) and (465.25,282.13) .. (465.25,286.8) ;
%Straight Lines [id:da9902399507274253] 
\draw    (620,273) -- (620,253.14) ;
\draw [shift={(620,251.14)}, rotate = 90] [color={rgb, 255:red, 0; green, 0; blue, 0 }  ][line width=0.75]    (10.93,-3.29) .. controls (6.95,-1.4) and (3.31,-0.3) .. (0,0) .. controls (3.31,0.3) and (6.95,1.4) .. (10.93,3.29)   ;

% Text Node
\draw (110,64.85) node  [color={rgb, 255:red, 0; green, 0; blue, 0 }  ,opacity=1 ] [align=left] {\begin{minipage}[lt]{92.93pt}\setlength\topsep{0pt}
\begin{center}
programa en Python\\
\small{(\cref{sec:programas-python})}
\end{center}

\end{minipage}};
% Text Node
\draw (219.89,51.97) node  [font=\small,color={rgb, 255:red, 0; green, 0; blue, 0 }  ,opacity=1 ] [align=left] {\begin{minipage}[lt]{54.86pt}\setlength\topsep{0pt}
\begin{center}
\Cref{sec:equivalencia}
\end{center}

\end{minipage}};
% Text Node
\draw (69.5,15.5) node  [font=\small,color={rgb, 255:red, 0; green, 0; blue, 0 }  ,opacity=1 ] [align=left] {\begin{minipage}[lt]{41.55pt}\setlength\topsep{0pt}
\cref{ch:programas-python-maquinas-turing}
\end{minipage}};
% Text Node
\draw (330,64.85) node  [color={rgb, 255:red, 0; green, 0; blue, 0 }  ,opacity=1 ] [align=left] {\begin{minipage}[lt]{92.93pt}\setlength\topsep{0pt}
\begin{center}
máquina de Turing\\
\small{(\cref{sec:maquinas-turing})}
\end{center}

\end{minipage}};
% Text Node
\draw (539.6,314.95) node  [color={rgb, 255:red, 0; green, 0; blue, 0 }  ,opacity=1 ] [align=left] {\begin{minipage}[lt]{92.93pt}\setlength\topsep{0pt}
\begin{center}
sistema lógico\\
\small{(\cref{sec:sistemas-logicos})}
\end{center}

\end{minipage}};
% Text Node
\draw (380.5,500.11) node  [color={rgb, 255:red, 0; green, 0; blue, 0 }  ,opacity=1 ] [align=left] {\begin{minipage}[lt]{159.55pt}\setlength\topsep{0pt}
\begin{center}
\textbf{Primer Teorema de Incompletitud}
\end{center}

\end{minipage}};
% Text Node
\draw (620.53,64.85) node  [color={rgb, 255:red, 0; green, 0; blue, 0 }  ,opacity=1 ] [align=left] {\begin{minipage}[lt]{106.99pt}\setlength\topsep{0pt}
\begin{center}
tesis de Church-Turing\\
\small{(\cref{sec:church-turing})}
\end{center}

\end{minipage}};
% Text Node
\draw (480.28,52.8) node  [font=\small,color={rgb, 255:red, 0; green, 0; blue, 0 }  ,opacity=1 ] [align=left] {\begin{minipage}[lt]{81.9pt}\setlength\topsep{0pt}
\begin{center}
\cref{sec:church-turing}
\end{center}

\end{minipage}};
% Text Node
\draw (589.5,518.3) node  [font=\small,color={rgb, 255:red, 0; green, 0; blue, 0 }  ,opacity=1 ] [align=left] {\begin{minipage}[lt]{54.86pt}\setlength\topsep{0pt}
\begin{center}
\Cref{subsec:mente-humana-tesis-church-turing}
\end{center}

\end{minipage}};
% Text Node
\draw (220.2,159.62) node  [color={rgb, 255:red, 0; green, 0; blue, 0 }  ,opacity=1 ] [align=left] {\begin{minipage}[lt]{230.11pt}\setlength\topsep{0pt}
\begin{center}
problema computacional
\end{center}

\end{minipage}};
% Text Node
\draw (98.02,204.85) node  [color={rgb, 255:red, 0; green, 0; blue, 0 }  ,opacity=1 ] [align=left] {\begin{minipage}[lt]{76.13pt}\setlength\topsep{0pt}
\begin{center}
decidible\\
\small{(\cref{sec:problemas-decidibles})}
\end{center}

\end{minipage}};
% Text Node
\draw (220.67,204.85) node  [color={rgb, 255:red, 0; green, 0; blue, 0 }  ,opacity=1 ] [align=left] {\begin{minipage}[lt]{76.13pt}\setlength\topsep{0pt}
\begin{center}
no decidible\\
\small{(\cref{sec:problemas-no-decidibles})}
\end{center}

\end{minipage}};
% Text Node
\draw (343.31,204.85) node  [color={rgb, 255:red, 0; green, 0; blue, 0 }  ,opacity=1 ] [align=left] {\begin{minipage}[lt]{76.13pt}\setlength\topsep{0pt}
\begin{center}
semidecidible\\
\small{(\cref{sec:problemas-semidecidibles})}
\end{center}

\end{minipage}};
% Text Node
\draw (110.4,314.95) node  [color={rgb, 255:red, 0; green, 0; blue, 0 }  ,opacity=1 ] [align=left] {\begin{minipage}[lt]{92.93pt}\setlength\topsep{0pt}
\begin{center}
universalidad\\
\small{(\cref{sec:universalidad})}
\end{center}

\end{minipage}};
% Text Node
\draw (330,314.85) node  [color={rgb, 255:red, 0; green, 0; blue, 0 }  ,opacity=1 ] [align=left] {\begin{minipage}[lt]{92.93pt}\setlength\topsep{0pt}
\begin{center}
reducción\\
\small{(\cref{sec:reducciones})}
\end{center}

\end{minipage}};
% Text Node
\draw (220,414.85) node  [color={rgb, 255:red, 0; green, 0; blue, 0 }  ,opacity=1 ] [align=left] {\begin{minipage}[lt]{107.26pt}\setlength\topsep{0pt}
\begin{center}
problema de la parada\\
\small{(\cref{sec:problema-parada})}
\end{center}

\end{minipage}};
% Text Node
\draw (540.33,415.05) node  [color={rgb, 255:red, 0; green, 0; blue, 0 }  ,opacity=1 ] [align=left] {\begin{minipage}[lt]{107.26pt}\setlength\topsep{0pt}
\begin{center}
aritmética de Peano\\
\small{(\cref{sec:aritmetica-peano})}
\end{center}

\end{minipage}};
% Text Node
\draw (698.8,314.95) node  [color={rgb, 255:red, 0; green, 0; blue, 0 }  ,opacity=1 ] [align=left] {\begin{minipage}[lt]{92.93pt}\setlength\topsep{0pt}
\begin{center}
sistema formal\\
\small{(\cref{sec:sistemas-formales})}
\end{center}

\end{minipage}};
% Text Node
\draw (299.6,544.95) node  [color={rgb, 255:red, 0; green, 0; blue, 0 }  ,opacity=1 ] [align=left] {\begin{minipage}[lt]{92.93pt}\setlength\topsep{0pt}
\begin{center}
versión semántica\\
\small{(\cref{sec:version-semantica})}
\end{center}

\end{minipage}};
% Text Node
\draw (459.6,545.05) node  [color={rgb, 255:red, 0; green, 0; blue, 0 }  ,opacity=1 ] [align=left] {\begin{minipage}[lt]{92.93pt}\setlength\topsep{0pt}
\begin{center}
versión sintáctica\\
\small{(\cref{sec:version-sintactica})}
\end{center}

\end{minipage}};
% Text Node
\draw (60.5,464.5) node  [font=\small,color={rgb, 255:red, 0; green, 0; blue, 0 }  ,opacity=1 ] [align=left] {\begin{minipage}[lt]{41.55pt}\setlength\topsep{0pt}
\cref{ch:problema-parada}
\end{minipage}};
% Text Node
\draw (749.5,465.5) node  [font=\small,color={rgb, 255:red, 0; green, 0; blue, 0 }  ,opacity=1 ] [align=left] {\begin{minipage}[lt]{41.55pt}\setlength\topsep{0pt}
\begin{flushright}
\cref{ch:sistemas-logicos}
\end{flushright}

\end{minipage}};
% Text Node
\draw (170.5,604.5) node  [font=\small,color={rgb, 255:red, 0; green, 0; blue, 0 }  ,opacity=1 ] [align=left] {\begin{minipage}[lt]{41.55pt}\setlength\topsep{0pt}
\cref{ch:teorema-incompletitud}
\end{minipage}};
% Text Node
\draw (643.76,195.55) node  [color={rgb, 255:red, 0; green, 0; blue, 0 }  ,opacity=1 ] [align=left] {\begin{minipage}[lt]{217.92pt}\setlength\topsep{0pt}
-- sólido {\small{(\cref{sec:solidez-completitud-decidibilidad})}}\\
-- decidible {\small{(\cref{sec:solidez-completitud-decidibilidad})}}\\
-- (semánticamente) completo {\small{(\cref{sec:solidez-completitud-decidibilidad})}}\\
-- sintácticamente completo {\small{(\cref{sec:consistencia})}}\\
-- sem./sint. consistente {\small{(\cref{sec:consistencia})}}
\end{minipage}};
% Text Node
\draw (414.84,645.25) node  [color={rgb, 255:red, 0; green, 0; blue, 0 }  ,opacity=1 ] [align=left] {\begin{minipage}[lt]{534.7pt}\setlength\topsep{0pt}
\begin{center}
Esquema resumen de este Trabajo de Fin de Grado
\end{center}
\end{minipage}};
% Text Node
\draw (380.15,407.8) node  [font=\small,color={rgb, 255:red, 0; green, 0; blue, 0 }  ,opacity=1 ] [align=left] {\begin{minipage}[lt]{108.93pt}\setlength\topsep{0pt}
\begin{center}
\cref{sec:primera-aproximacion}
\end{center}

\end{minipage}};

\end{tikzpicture}
\raisebox{-15.5cm}[0pt][0pt]{%
\resizebox{\linewidth}{!}{

\tikzset{every picture/.style={line width=0.75pt}} %set default line width to 0.75pt        

\begin{tikzpicture}[x=0.75pt,y=0.75pt,yscale=-1,xscale=1]
%uncomment if require: \path (0,694); %set diagram left start at 0, and has height of 694

%Straight Lines [id:da2056540562354816] 
\draw  [dash pattern={on 4.5pt off 4.5pt}]  (672,64.47) -- (810,64.47) ;
%Straight Lines [id:da23044041717189923] 
\draw  [dash pattern={on 4.5pt off 4.5pt}]  (412.67,64.47) -- (610,64.47) ;
\draw [shift={(410.67,64.47)}, rotate = 0] [color={rgb, 255:red, 0; green, 0; blue, 0 }  ][line width=0.75]    (10.93,-3.29) .. controls (6.95,-1.4) and (3.31,-0.3) .. (0,0) .. controls (3.31,0.3) and (6.95,1.4) .. (10.93,3.29)   ;
%Straight Lines [id:da8398435684999113] 
\draw  [dash pattern={on 4.5pt off 4.5pt}]  (810,64.47) -- (810,530) ;
%Rounded Rect [id:dp027018508739026448] 
\draw  [draw opacity=0][fill={rgb, 255:red, 230; green, 230; blue, 230 }  ,fill opacity=0.5 ] (20,439.33) .. controls (20,445.22) and (24.78,450) .. (30.67,450) -- (409.33,450) .. controls (415.22,450) and (420,445.22) .. (420,439.33) -- (420,140.67) .. controls (420,134.78) and (415.22,130) .. (409.33,130) -- (30.67,130) .. controls (24.78,130) and (20,134.78) .. (20,140.67) -- cycle ;
%Rounded Rect [id:dp9415679498487528] 
\draw  [dash pattern={on 4.5pt off 4.5pt}] (29.83,292.77) .. controls (29.83,285.88) and (35.41,280.3) .. (42.3,280.3) -- (398.54,280.3) .. controls (405.42,280.3) and (411,285.88) .. (411,292.77) -- (411,337.17) .. controls (411,344.05) and (405.42,349.64) .. (398.54,349.64) -- (42.3,349.64) .. controls (35.41,349.64) and (29.83,344.05) .. (29.83,337.17) -- cycle ;
%Rounded Rect [id:dp05100022295548956] 
\draw  [draw opacity=0][fill={rgb, 255:red, 230; green, 230; blue, 230 }  ,fill opacity=0.5 ] (130.5,575) .. controls (130.5,583.28) and (137.22,590) .. (145.5,590) -- (615,590) .. controls (623.28,590) and (630,583.28) .. (630,575) -- (630,395) .. controls (630,386.72) and (623.28,380) .. (615,380) -- (145.5,380) .. controls (137.22,380) and (130.5,386.72) .. (130.5,395) -- cycle ;
%Rounded Rect [id:dp6496809750851564] 
\draw  [draw opacity=0][fill={rgb, 255:red, 230; green, 230; blue, 230 }  ,fill opacity=0.5 ] (450.61,440.79) .. controls (450.61,445.82) and (454.69,449.9) .. (459.72,449.9) -- (780.89,449.9) .. controls (785.92,449.9) and (790,445.82) .. (790,440.79) -- (790,139.11) .. controls (790,134.08) and (785.92,130) .. (780.89,130) -- (459.72,130) .. controls (454.69,130) and (450.61,134.08) .. (450.61,139.11) -- cycle ;
%Rounded Rect [id:dp549636265291094] 
\draw  [fill={rgb, 255:red, 230; green, 230; blue, 230 }  ,fill opacity=0.5 ][dash pattern={on 4.5pt off 4.5pt}] (29.83,42.77) .. controls (29.83,35.88) and (35.41,30.3) .. (42.3,30.3) -- (398.54,30.3) .. controls (405.42,30.3) and (411,35.88) .. (411,42.77) -- (411,87.17) .. controls (411,94.05) and (405.42,99.64) .. (398.54,99.64) -- (42.3,99.64) .. controls (35.41,99.64) and (29.83,94.05) .. (29.83,87.17) -- cycle ;
%Rounded Rect [id:dp1378581119146769] 
\draw  [fill={rgb, 255:red, 255; green, 255; blue, 255 }  ,fill opacity=1 ] (40.4,49.82) .. controls (40.4,44.29) and (44.89,39.8) .. (50.42,39.8) -- (169.58,39.8) .. controls (175.11,39.8) and (179.6,44.29) .. (179.6,49.82) -- (179.6,79.88) .. controls (179.6,85.41) and (175.11,89.9) .. (169.58,89.9) -- (50.42,89.9) .. controls (44.89,89.9) and (40.4,85.41) .. (40.4,79.88) -- cycle ;
%Straight Lines [id:da8729153203483555] 
\draw    (188.42,63.77) -- (250.99,63.77)(188.42,66.77) -- (250.99,66.77) ;
\draw [shift={(257.99,65.27)}, rotate = 180] [color={rgb, 255:red, 0; green, 0; blue, 0 }  ][line width=0.75]    (10.93,-4.9) .. controls (6.95,-2.3) and (3.31,-0.67) .. (0,0) .. controls (3.31,0.67) and (6.95,2.3) .. (10.93,4.9)   ;
\draw [shift={(181.42,65.27)}, rotate = 0] [color={rgb, 255:red, 0; green, 0; blue, 0 }  ][line width=0.75]    (10.93,-4.9) .. controls (6.95,-2.3) and (3.31,-0.67) .. (0,0) .. controls (3.31,0.67) and (6.95,2.3) .. (10.93,4.9)   ;
%Straight Lines [id:da845243590120559] 
\draw    (220.17,99.97) -- (220.17,137.64) ;
\draw [shift={(220.17,139.64)}, rotate = 270] [color={rgb, 255:red, 0; green, 0; blue, 0 }  ][line width=0.75]    (10.93,-3.29) .. controls (6.95,-1.4) and (3.31,-0.3) .. (0,0) .. controls (3.31,0.3) and (6.95,1.4) .. (10.93,3.29)   ;
%Straight Lines [id:da7497001546502899] 
\draw    (380,420.33) -- (380,478) ;
\draw [shift={(380,480)}, rotate = 270] [color={rgb, 255:red, 0; green, 0; blue, 0 }  ][line width=0.75]    (10.93,-3.29) .. controls (6.95,-1.4) and (3.31,-0.3) .. (0,0) .. controls (3.31,0.3) and (6.95,1.4) .. (10.93,3.29)   ;
%Straight Lines [id:da5689893955334542] 
\draw    (300,420) -- (460,420) ;
%Rounded Rect [id:dp18582062333838767] 
\draw  [fill={rgb, 255:red, 255; green, 255; blue, 255 }  ,fill opacity=1 ] (260.4,49.82) .. controls (260.4,44.29) and (264.89,39.8) .. (270.42,39.8) -- (389.58,39.8) .. controls (395.11,39.8) and (399.6,44.29) .. (399.6,49.82) -- (399.6,79.88) .. controls (399.6,85.41) and (395.11,89.9) .. (389.58,89.9) -- (270.42,89.9) .. controls (264.89,89.9) and (260.4,85.41) .. (260.4,79.88) -- cycle ;
%Rounded Rect [id:dp9523169763067778] 
\draw  [fill={rgb, 255:red, 255; green, 255; blue, 255 }  ,fill opacity=1 ] (470,299.92) .. controls (470,294.39) and (474.49,289.9) .. (480.02,289.9) -- (599.18,289.9) .. controls (604.71,289.9) and (609.2,294.39) .. (609.2,299.92) -- (609.2,329.98) .. controls (609.2,335.51) and (604.71,340) .. (599.18,340) -- (480.02,340) .. controls (474.49,340) and (470,335.51) .. (470,329.98) -- cycle ;
%Rounded Rect [id:dp48544030815291705] 
\draw  [fill={rgb, 255:red, 255; green, 255; blue, 255 }  ,fill opacity=1 ] (220,490) .. controls (220,484.48) and (224.48,480) .. (230,480) -- (530,480) .. controls (535.52,480) and (540,484.48) .. (540,490) -- (540,570) .. controls (540,575.52) and (535.52,580) .. (530,580) -- (230,580) .. controls (224.48,580) and (220,575.52) .. (220,570) -- cycle ;
%Rounded Rect [id:dp08222627508455749] 
\draw  [fill={rgb, 255:red, 255; green, 255; blue, 255 }  ,fill opacity=1 ] (540.4,49.82) .. controls (540.4,44.29) and (544.89,39.8) .. (550.42,39.8) -- (690.65,39.8) .. controls (696.18,39.8) and (700.67,44.29) .. (700.67,49.82) -- (700.67,79.88) .. controls (700.67,85.41) and (696.18,89.9) .. (690.65,89.9) -- (550.42,89.9) .. controls (544.89,89.9) and (540.4,85.41) .. (540.4,79.88) -- cycle ;
%Straight Lines [id:da5868346364154593] 
\draw  [dash pattern={on 4.5pt off 4.5pt}]  (810,530) -- (542,530) ;
\draw [shift={(540,530)}, rotate = 360] [color={rgb, 255:red, 0; green, 0; blue, 0 }  ][line width=0.75]    (10.93,-3.29) .. controls (6.95,-1.4) and (3.31,-0.3) .. (0,0) .. controls (3.31,0.3) and (6.95,1.4) .. (10.93,3.29)   ;
%Rounded Rect [id:dp7392784471609453] 
\draw  [fill={rgb, 255:red, 255; green, 255; blue, 255 }  ,fill opacity=1 ] (30.33,150.47) .. controls (30.33,144.58) and (35.11,139.8) .. (41,139.8) -- (400.34,139.8) .. controls (406.23,139.8) and (411,144.58) .. (411,150.47) -- (411,229.34) .. controls (411,235.23) and (406.23,240.01) .. (400.34,240.01) -- (41,240.01) .. controls (35.11,240.01) and (30.33,235.23) .. (30.33,229.34) -- cycle ;
%Rounded Rect [id:dp90901637181607] 
\draw  [fill={rgb, 255:red, 255; green, 255; blue, 255 }  ,fill opacity=1 ] (41,189.82) .. controls (41,184.29) and (45.49,179.8) .. (51.02,179.8) -- (145.02,179.8) .. controls (150.55,179.8) and (155.04,184.29) .. (155.04,189.82) -- (155.04,219.88) .. controls (155.04,225.41) and (150.55,229.9) .. (145.02,229.9) -- (51.02,229.9) .. controls (45.49,229.9) and (41,225.41) .. (41,219.88) -- cycle ;
%Rounded Rect [id:dp06115018179625542] 
\draw  [fill={rgb, 255:red, 255; green, 255; blue, 255 }  ,fill opacity=1 ] (163.65,189.82) .. controls (163.65,184.29) and (168.13,179.8) .. (173.67,179.8) -- (267.67,179.8) .. controls (273.2,179.8) and (277.69,184.29) .. (277.69,189.82) -- (277.69,219.88) .. controls (277.69,225.41) and (273.2,229.9) .. (267.67,229.9) -- (173.67,229.9) .. controls (168.13,229.9) and (163.65,225.41) .. (163.65,219.88) -- cycle ;
%Rounded Rect [id:dp6195275548153432] 
\draw  [fill={rgb, 255:red, 255; green, 255; blue, 255 }  ,fill opacity=1 ] (286.3,189.82) .. controls (286.3,184.29) and (290.78,179.8) .. (296.32,179.8) -- (390.31,179.8) .. controls (395.85,179.8) and (400.33,184.29) .. (400.33,189.82) -- (400.33,219.88) .. controls (400.33,225.41) and (395.85,229.9) .. (390.31,229.9) -- (296.32,229.9) .. controls (290.78,229.9) and (286.3,225.41) .. (286.3,219.88) -- cycle ;
%Rounded Rect [id:dp41283415100159493] 
\draw  [fill={rgb, 255:red, 255; green, 255; blue, 255 }  ,fill opacity=1 ] (40.8,299.92) .. controls (40.8,294.39) and (45.29,289.9) .. (50.82,289.9) -- (169.98,289.9) .. controls (175.51,289.9) and (180,294.39) .. (180,299.92) -- (180,329.98) .. controls (180,335.51) and (175.51,340) .. (169.98,340) -- (50.82,340) .. controls (45.29,340) and (40.8,335.51) .. (40.8,329.98) -- cycle ;
%Rounded Rect [id:dp40154854661776396] 
\draw  [fill={rgb, 255:red, 255; green, 255; blue, 255 }  ,fill opacity=1 ] (260.4,299.82) .. controls (260.4,294.29) and (264.89,289.8) .. (270.42,289.8) -- (389.58,289.8) .. controls (395.11,289.8) and (399.6,294.29) .. (399.6,299.82) -- (399.6,329.88) .. controls (399.6,335.41) and (395.11,339.9) .. (389.58,339.9) -- (270.42,339.9) .. controls (264.89,339.9) and (260.4,335.41) .. (260.4,329.88) -- cycle ;
%Straight Lines [id:da6602630444249074] 
\draw    (180.17,315.8) -- (260,315.8) ;
%Straight Lines [id:da30905866347109745] 
\draw    (219.86,315.97) -- (219.86,386.81) ;
\draw [shift={(219.86,388.81)}, rotate = 270] [color={rgb, 255:red, 0; green, 0; blue, 0 }  ][line width=0.75]    (10.93,-3.29) .. controls (6.95,-1.4) and (3.31,-0.3) .. (0,0) .. controls (3.31,0.3) and (6.95,1.4) .. (10.93,3.29)   ;
%Rounded Rect [id:dp44627062039563214] 
\draw  [fill={rgb, 255:red, 255; green, 255; blue, 255 }  ,fill opacity=1 ] (139.67,399.82) .. controls (139.67,394.29) and (144.15,389.8) .. (149.69,389.8) -- (290.31,389.8) .. controls (295.85,389.8) and (300.33,394.29) .. (300.33,399.82) -- (300.33,429.88) .. controls (300.33,435.41) and (295.85,439.9) .. (290.31,439.9) -- (149.69,439.9) .. controls (144.15,439.9) and (139.67,435.41) .. (139.67,429.88) -- cycle ;
%Straight Lines [id:da8375870493480149] 
\draw    (220,240) -- (220,277.67) ;
\draw [shift={(220,279.67)}, rotate = 270] [color={rgb, 255:red, 0; green, 0; blue, 0 }  ][line width=0.75]    (10.93,-3.29) .. controls (6.95,-1.4) and (3.31,-0.3) .. (0,0) .. controls (3.31,0.3) and (6.95,1.4) .. (10.93,3.29)   ;
%Rounded Rect [id:dp37285624114093485] 
\draw  [fill={rgb, 255:red, 255; green, 255; blue, 255 }  ,fill opacity=1 ] (460,400.02) .. controls (460,394.49) and (464.49,390) .. (470.02,390) -- (610.65,390) .. controls (616.18,390) and (620.67,394.49) .. (620.67,400.02) -- (620.67,430.08) .. controls (620.67,435.61) and (616.18,440.1) .. (610.65,440.1) -- (470.02,440.1) .. controls (464.49,440.1) and (460,435.61) .. (460,430.08) -- cycle ;
%Rounded Rect [id:dp6883616335272991] 
\draw  [fill={rgb, 255:red, 255; green, 255; blue, 255 }  ,fill opacity=1 ] (629.2,299.92) .. controls (629.2,294.39) and (633.69,289.9) .. (639.22,289.9) -- (758.38,289.9) .. controls (763.91,289.9) and (768.4,294.39) .. (768.4,299.92) -- (768.4,329.98) .. controls (768.4,335.51) and (763.91,340) .. (758.38,340) -- (639.22,340) .. controls (633.69,340) and (629.2,335.51) .. (629.2,329.98) -- cycle ;
%Straight Lines [id:da19870419404011752] 
\draw    (540,340) -- (540,388) ;
\draw [shift={(540,390)}, rotate = 270] [color={rgb, 255:red, 0; green, 0; blue, 0 }  ][line width=0.75]    (10.93,-3.29) .. controls (6.95,-1.4) and (3.31,-0.3) .. (0,0) .. controls (3.31,0.3) and (6.95,1.4) .. (10.93,3.29)   ;
%Rounded Rect [id:dp9692784599420237] 
\draw  [fill={rgb, 255:red, 255; green, 255; blue, 255 }  ,fill opacity=1 ] (230,529.92) .. controls (230,524.39) and (234.49,519.9) .. (240.02,519.9) -- (359.18,519.9) .. controls (364.71,519.9) and (369.2,524.39) .. (369.2,529.92) -- (369.2,559.98) .. controls (369.2,565.51) and (364.71,570) .. (359.18,570) -- (240.02,570) .. controls (234.49,570) and (230,565.51) .. (230,559.98) -- cycle ;
%Rounded Rect [id:dp749684165723072] 
\draw  [fill={rgb, 255:red, 255; green, 255; blue, 255 }  ,fill opacity=1 ] (390,530.02) .. controls (390,524.49) and (394.49,520) .. (400.02,520) -- (519.18,520) .. controls (524.71,520) and (529.2,524.49) .. (529.2,530.02) -- (529.2,560.08) .. controls (529.2,565.61) and (524.71,570.1) .. (519.18,570.1) -- (400.02,570.1) .. controls (394.49,570.1) and (390,565.61) .. (390,560.08) -- cycle ;
%Rounded Rect [id:dp28001422897762485] 
\draw  [fill={rgb, 255:red, 255; green, 255; blue, 255 }  ,fill opacity=1 ][dash pattern={on 4.5pt off 4.5pt}] (460,148.92) .. controls (460,143.99) and (463.99,140) .. (468.92,140) -- (771.08,140) .. controls (776.01,140) and (780,143.99) .. (780,148.92) -- (780,241.55) .. controls (780,246.48) and (776.01,250.47) .. (771.08,250.47) -- (468.92,250.47) .. controls (463.99,250.47) and (460,246.48) .. (460,241.55) -- cycle ;
%Shape: Brace [id:dp9411387904179627] 
\draw   (774.67,286.8) .. controls (774.67,282.13) and (772.34,279.8) .. (767.67,279.8) -- (629.96,279.8) .. controls (623.29,279.8) and (619.96,277.47) .. (619.96,272.8) .. controls (619.96,277.47) and (616.63,279.8) .. (609.96,279.8)(612.96,279.8) -- (472.25,279.8) .. controls (467.58,279.8) and (465.25,282.13) .. (465.25,286.8) ;
%Straight Lines [id:da9902399507274253] 
\draw    (620,273) -- (620,253.14) ;
\draw [shift={(620,251.14)}, rotate = 90] [color={rgb, 255:red, 0; green, 0; blue, 0 }  ][line width=0.75]    (10.93,-3.29) .. controls (6.95,-1.4) and (3.31,-0.3) .. (0,0) .. controls (3.31,0.3) and (6.95,1.4) .. (10.93,3.29)   ;

% Text Node
\draw (110,64.85) node  [color={rgb, 255:red, 0; green, 0; blue, 0 }  ,opacity=1 ] [align=left] {\begin{minipage}[lt]{92.93pt}\setlength\topsep{0pt}
\begin{center}
programa en Python\\
\small{(\cref{sec:programas-python})}
\end{center}

\end{minipage}};
% Text Node
\draw (219.89,51.97) node  [font=\small,color={rgb, 255:red, 0; green, 0; blue, 0 }  ,opacity=1 ] [align=left] {\begin{minipage}[lt]{54.86pt}\setlength\topsep{0pt}
\begin{center}
\Cref{sec:equivalencia}
\end{center}

\end{minipage}};
% Text Node
\draw (69.5,15.5) node  [font=\small,color={rgb, 255:red, 0; green, 0; blue, 0 }  ,opacity=1 ] [align=left] {\begin{minipage}[lt]{41.55pt}\setlength\topsep{0pt}
\cref{ch:programas-python-maquinas-turing}
\end{minipage}};
% Text Node
\draw (330,64.85) node  [color={rgb, 255:red, 0; green, 0; blue, 0 }  ,opacity=1 ] [align=left] {\begin{minipage}[lt]{92.93pt}\setlength\topsep{0pt}
\begin{center}
máquina de Turing\\
\small{(\cref{sec:maquinas-turing})}
\end{center}

\end{minipage}};
% Text Node
\draw (539.6,314.95) node  [color={rgb, 255:red, 0; green, 0; blue, 0 }  ,opacity=1 ] [align=left] {\begin{minipage}[lt]{92.93pt}\setlength\topsep{0pt}
\begin{center}
sistema lógico\\
\small{(\cref{sec:sistemas-logicos})}
\end{center}

\end{minipage}};
% Text Node
\draw (380.5,500.11) node  [color={rgb, 255:red, 0; green, 0; blue, 0 }  ,opacity=1 ] [align=left] {\begin{minipage}[lt]{159.55pt}\setlength\topsep{0pt}
\begin{center}
\textbf{Primer Teorema de Incompletitud}
\end{center}

\end{minipage}};
% Text Node
\draw (620.53,64.85) node  [color={rgb, 255:red, 0; green, 0; blue, 0 }  ,opacity=1 ] [align=left] {\begin{minipage}[lt]{106.99pt}\setlength\topsep{0pt}
\begin{center}
tesis de Church-Turing\\
\small{(\cref{sec:church-turing})}
\end{center}

\end{minipage}};
% Text Node
\draw (480.28,52.8) node  [font=\small,color={rgb, 255:red, 0; green, 0; blue, 0 }  ,opacity=1 ] [align=left] {\begin{minipage}[lt]{81.9pt}\setlength\topsep{0pt}
\begin{center}
\cref{sec:church-turing}
\end{center}

\end{minipage}};
% Text Node
\draw (589.5,518.3) node  [font=\small,color={rgb, 255:red, 0; green, 0; blue, 0 }  ,opacity=1 ] [align=left] {\begin{minipage}[lt]{54.86pt}\setlength\topsep{0pt}
\begin{center}
\Cref{subsec:mente-humana-tesis-church-turing}
\end{center}

\end{minipage}};
% Text Node
\draw (220.2,159.62) node  [color={rgb, 255:red, 0; green, 0; blue, 0 }  ,opacity=1 ] [align=left] {\begin{minipage}[lt]{230.11pt}\setlength\topsep{0pt}
\begin{center}
problema computacional
\end{center}

\end{minipage}};
% Text Node
\draw (98.02,204.85) node  [color={rgb, 255:red, 0; green, 0; blue, 0 }  ,opacity=1 ] [align=left] {\begin{minipage}[lt]{76.13pt}\setlength\topsep{0pt}
\begin{center}
decidible\\
\small{(\cref{sec:problemas-decidibles})}
\end{center}

\end{minipage}};
% Text Node
\draw (220.67,204.85) node  [color={rgb, 255:red, 0; green, 0; blue, 0 }  ,opacity=1 ] [align=left] {\begin{minipage}[lt]{76.13pt}\setlength\topsep{0pt}
\begin{center}
no decidible\\
\small{(\cref{sec:problemas-no-decidibles})}
\end{center}

\end{minipage}};
% Text Node
\draw (343.31,204.85) node  [color={rgb, 255:red, 0; green, 0; blue, 0 }  ,opacity=1 ] [align=left] {\begin{minipage}[lt]{76.13pt}\setlength\topsep{0pt}
\begin{center}
semidecidible\\
\small{(\cref{sec:problemas-semidecidibles})}
\end{center}

\end{minipage}};
% Text Node
\draw (110.4,314.95) node  [color={rgb, 255:red, 0; green, 0; blue, 0 }  ,opacity=1 ] [align=left] {\begin{minipage}[lt]{92.93pt}\setlength\topsep{0pt}
\begin{center}
universalidad\\
\small{(\cref{sec:universalidad})}
\end{center}

\end{minipage}};
% Text Node
\draw (330,314.85) node  [color={rgb, 255:red, 0; green, 0; blue, 0 }  ,opacity=1 ] [align=left] {\begin{minipage}[lt]{92.93pt}\setlength\topsep{0pt}
\begin{center}
reducción\\
\small{(\cref{sec:reducciones})}
\end{center}

\end{minipage}};
% Text Node
\draw (220,414.85) node  [color={rgb, 255:red, 0; green, 0; blue, 0 }  ,opacity=1 ] [align=left] {\begin{minipage}[lt]{107.26pt}\setlength\topsep{0pt}
\begin{center}
problema de la parada\\
\small{(\cref{sec:problema-parada})}
\end{center}

\end{minipage}};
% Text Node
\draw (540.33,415.05) node  [color={rgb, 255:red, 0; green, 0; blue, 0 }  ,opacity=1 ] [align=left] {\begin{minipage}[lt]{107.26pt}\setlength\topsep{0pt}
\begin{center}
aritmética de Peano\\
\small{(\cref{sec:aritmetica-peano})}
\end{center}

\end{minipage}};
% Text Node
\draw (698.8,314.95) node  [color={rgb, 255:red, 0; green, 0; blue, 0 }  ,opacity=1 ] [align=left] {\begin{minipage}[lt]{92.93pt}\setlength\topsep{0pt}
\begin{center}
sistema formal\\
\small{(\cref{sec:sistemas-formales})}
\end{center}

\end{minipage}};
% Text Node
\draw (299.6,544.95) node  [color={rgb, 255:red, 0; green, 0; blue, 0 }  ,opacity=1 ] [align=left] {\begin{minipage}[lt]{92.93pt}\setlength\topsep{0pt}
\begin{center}
versión semántica\\
\small{(\cref{sec:version-semantica})}
\end{center}

\end{minipage}};
% Text Node
\draw (459.6,545.05) node  [color={rgb, 255:red, 0; green, 0; blue, 0 }  ,opacity=1 ] [align=left] {\begin{minipage}[lt]{92.93pt}\setlength\topsep{0pt}
\begin{center}
versión sintáctica\\
\small{(\cref{sec:version-sintactica})}
\end{center}

\end{minipage}};
% Text Node
\draw (60.5,464.5) node  [font=\small,color={rgb, 255:red, 0; green, 0; blue, 0 }  ,opacity=1 ] [align=left] {\begin{minipage}[lt]{41.55pt}\setlength\topsep{0pt}
\cref{ch:problema-parada}
\end{minipage}};
% Text Node
\draw (749.5,465.5) node  [font=\small,color={rgb, 255:red, 0; green, 0; blue, 0 }  ,opacity=1 ] [align=left] {\begin{minipage}[lt]{41.55pt}\setlength\topsep{0pt}
\begin{flushright}
\cref{ch:sistemas-logicos}
\end{flushright}

\end{minipage}};
% Text Node
\draw (170.5,604.5) node  [font=\small,color={rgb, 255:red, 0; green, 0; blue, 0 }  ,opacity=1 ] [align=left] {\begin{minipage}[lt]{41.55pt}\setlength\topsep{0pt}
\cref{ch:teorema-incompletitud}
\end{minipage}};
% Text Node
\draw (643.76,195.55) node  [color={rgb, 255:red, 0; green, 0; blue, 0 }  ,opacity=1 ] [align=left] {\begin{minipage}[lt]{217.92pt}\setlength\topsep{0pt}
-- sólido {\small{(\cref{sec:solidez-completitud-decidibilidad})}}\\
-- decidible {\small{(\cref{sec:solidez-completitud-decidibilidad})}}\\
-- (semánticamente) completo {\small{(\cref{sec:solidez-completitud-decidibilidad})}}\\
-- sintácticamente completo {\small{(\cref{sec:consistencia})}}\\
-- sem./sint. consistente {\small{(\cref{sec:consistencia})}}
\end{minipage}};
% Text Node
\draw (414.84,645.25) node  [color={rgb, 255:red, 0; green, 0; blue, 0 }  ,opacity=1 ] [align=left] {\begin{minipage}[lt]{534.7pt}\setlength\topsep{0pt}
\begin{center}
Esquema resumen de este Trabajo de Fin de Grado
\end{center}
\end{minipage}};
% Text Node
\draw (380.15,407.8) node  [font=\small,color={rgb, 255:red, 0; green, 0; blue, 0 }  ,opacity=1 ] [align=left] {\begin{minipage}[lt]{108.93pt}\setlength\topsep{0pt}
\begin{center}
\cref{sec:primera-aproximacion}
\end{center}

\end{minipage}};

\end{tikzpicture}}
}
%\caption{Esquema resumen (extendido)}

%\label{fig:resumen-extendido}
\end{figure}
% ====================
\end{landscape}

\endinput

\end{adjustbox}

\end{frame}

\section[Programas en Python y\\máquinas de Turing]{Programas en Python y máquinas de Turing}

\begin{frame}<presentation:0>{Programas en Python}

\textbf{Definición 3.1} (Función SISO)\textbf{.}

Definimos una \hle{función SISO} como una función de Python que acepta un único parámetro de tipo \emph{string} como entrada, y que devuelve un \emph{string}.

\vspace{5mm}

\textbf{Definición 3.2} (Programa en Python de tipo SISO)\textbf{.}

Respecto a un sistema informático de referencia $C$, definimos un \hle{programa en Python de tipo SISO} como una cadena de caracteres $P$, de modo que:

\begin{itemize}
    \item $P$ es código en Python sintácticamente correcto;
    \item $P$ define al menos una función, y la primera de ellas es de \hle{tipo SISO}. A esta función la llamaremos \hle{función main}.
\end{itemize}

\end{frame}

\begin{frame}<presentation:0>{Programas en Python}

\textbf{Definición 3.3} (Salida de un programa)\textbf{.}

Sea $P$ un programa en Python respecto a un sistema informático $C$. Sea $M$ la función \emph{main} de $P$, e $I$ la \hle{entrada} de $P$, una cadena de caracteres. La \hle{salida de $P$ respecto a la entrada $I$}, $P(I)$, se produce usando $C$ para ejecutar $P$ con entrada $I$, y se define como:

\begin{itemize}
    \item Si $M$ devuelve una \ul{cadena de caracteres $O$}, entonces \ul{$P(I)=O$}.
    \item Si $M$ devuelve \ul{cualquier otro objeto} en Python, entonces \ul{$P(I)$ no está definido}.
    \item Si $M$ lanza una \ul{excepción}, entonces \ul{$P(I)$ no está definido}.
    \item Si $M$ \ul{no termina}, entonces \ul{$P(I)$ no está definido}.
\end{itemize}

\end{frame}

\begin{frame}<presentation:0>{Máquinas de Turing}

\textbf{Definición 3.5} (Máquina de Turing)\textbf{.}

Una \hle{máquina de Turing} es una séptupla $M=(Q, A, B, \delta, q_0, \Vtextvisiblespace\:, F)$, en la que:

\begin{itemize}
    \item $Q = \{q_0, ...\}$ es un conjunto (finito) de \hle{estados}.
    \item $A$ es un \hle{alfabeto de entrada}.
    \item $B$ es un \emph{alfabeto de símbolos de la cinta} (o \hle{alfabeto de trabajo}), con \ul{$B\subset A$}.
    \item $\delta : Q \times B \longrightarrow (Q \times B \times \{I, D, S\}) \cup \emptyset$ es la \hle{función de transición} que asigna a cada estado $q\in Q$ y símbolo $b\in B$, el valor $\delta(q, b)$ que puede ser:
    \begin{itemize}
        \item $\emptyset$ en caso de no estar definido, o
        \item una tripleta $(p, c, M)$, donde $p\in Q$, $c\in B$, $M \in \{I, D, S\}$.
    \end{itemize}
    \item $q_0\in Q$ es el \hle{estado inicial}.
    \item $\Vtextvisiblespace\in B \setminus A$ es el \hle{símbolo blanco}.
    \item $F \subseteq Q$ es el conjunto de \hle{estados finales}.
\end{itemize}

\end{frame}

\begin{frame}{Máquinas de Turing}

\vspace{5mm}
\begin{adjustbox}{max height={.67\textheight},center}


\tikzset{every picture/.style={line width=0.75pt}} %set default line width to 0.75pt        

\begin{tikzpicture}[x=0.75pt,y=0.75pt,yscale=-1,xscale=1]
%uncomment if require: \path (0,221); %set diagram left start at 0, and has height of 221

%Flowchart: Direct Access Storage [id:dp5745083904317894] 
\draw   (45.92,32.5) -- (91.42,32.5) .. controls (98.18,32.5) and (103.67,54.66) .. (103.67,82) .. controls (103.67,109.34) and (98.18,131.5) .. (91.42,131.5) -- (45.92,131.5)(33.67,82) .. controls (33.67,54.66) and (39.15,32.5) .. (45.92,32.5) .. controls (52.68,32.5) and (58.17,54.66) .. (58.17,82) .. controls (58.17,109.34) and (52.68,131.5) .. (45.92,131.5) .. controls (39.15,131.5) and (33.67,109.34) .. (33.67,82) ;
%Shape: Cube [id:dp9677041927446011] 
\draw   (132.67,13.07) -- (121.59,2) -- (2,2) -- (2,140.43) -- (13.07,151.5) -- (132.67,151.5) -- cycle ; \draw   (2,2) -- (13.07,13.07) -- (132.67,13.07) ; \draw   (13.07,13.07) -- (13.07,151.5) ;
%Shape: Ellipse [id:dp3944720200277354] 
\draw  [fill={rgb, 255:red, 0; green, 0; blue, 0 }  ,fill opacity=1 ] (42.81,82.14) .. controls (42.81,72.36) and (44.22,64.43) .. (45.95,64.43) .. controls (47.69,64.43) and (49.1,72.36) .. (49.1,82.14) .. controls (49.1,91.93) and (47.69,99.86) .. (45.95,99.86) .. controls (44.22,99.86) and (42.81,91.93) .. (42.81,82.14) -- cycle ;
%Straight Lines [id:da7811026207895715] 
\draw    (57.67,95) -- (103.1,95) ;
%Straight Lines [id:da7715257519556658] 
\draw    (54.52,116.86) -- (99.95,116.86) ;
%Straight Lines [id:da7719478908390764] 
\draw    (55.1,47.86) -- (100.52,47.86) ;
%Straight Lines [id:da44878361199349515] 
\draw    (57.67,69) -- (103.1,69) ;
%Flowchart: Merge [id:dp7750110622325217] 
\draw  [fill={rgb, 255:red, 0; green, 0; blue, 0 }  ,fill opacity=1 ] (111.95,76.14) -- (111.95,88) -- (100.24,82.07) -- cycle ;
%Shape: Rectangle [id:dp3100328904459395] 
\draw   (232.67,152.17) -- (262.67,152.17) -- (262.67,182.17) -- (232.67,182.17) -- cycle ;
%Shape: Rectangle [id:dp1692041741743724] 
\draw  [line width=2.25]  (262.67,152.17) -- (292.67,152.17) -- (292.67,182.17) -- (262.67,182.17) -- cycle ;
%Shape: Rectangle [id:dp5247156403937852] 
\draw   (292.67,152.17) -- (322.67,152.17) -- (322.67,182.17) -- (292.67,182.17) -- cycle ;
%Shape: Rectangle [id:dp1939162798122891] 
\draw   (322.67,152.17) -- (352.67,152.17) -- (352.67,182.17) -- (322.67,182.17) -- cycle ;
%Shape: Rectangle [id:dp955944006123826] 
\draw   (352.67,152.17) -- (382.67,152.17) -- (382.67,182.17) -- (352.67,182.17) -- cycle ;
%Straight Lines [id:da2297213302864911] 
\draw    (382.67,152.17) -- (402.67,152.17) ;
%Straight Lines [id:da0017979141787713981] 
\draw    (382.67,182.17) -- (402.67,182.17) ;
%Straight Lines [id:da5196954519515742] 
\draw    (212.67,182.17) -- (232.67,182.17) ;
%Straight Lines [id:da6250621041798603] 
\draw    (212.67,152.17) -- (232.67,152.17) ;
%Straight Lines [id:da4041928331954272] 
\draw [color={rgb, 255:red, 128; green, 128; blue, 128 }  ,draw opacity=1 ]   (262.67,195.17) -- (352.67,195.17) ;
%Straight Lines [id:da40491651879492707] 
\draw [color={rgb, 255:red, 128; green, 128; blue, 128 }  ,draw opacity=1 ]   (262.67,188.17) -- (262.67,195.17) ;
%Straight Lines [id:da9178129384480556] 
\draw [color={rgb, 255:red, 128; green, 128; blue, 128 }  ,draw opacity=1 ]   (352.67,188.17) -- (352.67,195.17) ;
%Shape: Right Angle [id:dp42415196847546355] 
\draw  [line width=1.5]  (132,82) -- (278,82) -- (278,146) ;
%Straight Lines [id:da05158428581441665] 
\draw    (277.67,128) -- (277.67,149.17) ;
\draw [shift={(277.67,152.17)}, rotate = 270] [fill={rgb, 255:red, 0; green, 0; blue, 0 }  ][line width=0.08]  [draw opacity=0] (8.93,-4.29) -- (0,0) -- (8.93,4.29) -- cycle    ;
%Straight Lines [id:da7132164152395672] 
\draw [color={rgb, 255:red, 128; green, 128; blue, 128 }  ,draw opacity=1 ]   (155.33,46) -- (118.18,78.03) ;
\draw [shift={(116.67,79.33)}, rotate = 319.24] [color={rgb, 255:red, 128; green, 128; blue, 128 }  ,draw opacity=1 ][line width=0.75]    (10.93,-3.29) .. controls (6.95,-1.4) and (3.31,-0.3) .. (0,0) .. controls (3.31,0.3) and (6.95,1.4) .. (10.93,3.29)   ;
%Straight Lines [id:da42945985402945897] 
\draw [color={rgb, 255:red, 128; green, 128; blue, 128 }  ,draw opacity=1 ]   (457.33,167.33) -- (418,167.33) ;
\draw [shift={(416,167.33)}, rotate = 360] [color={rgb, 255:red, 128; green, 128; blue, 128 }  ,draw opacity=1 ][line width=0.75]    (10.93,-3.29) .. controls (6.95,-1.4) and (3.31,-0.3) .. (0,0) .. controls (3.31,0.3) and (6.95,1.4) .. (10.93,3.29)   ;
%Straight Lines [id:da3859416199015755] 
\draw [color={rgb, 255:red, 128; green, 128; blue, 128 }  ,draw opacity=1 ]   (328,109.33) -- (290.85,141.36) ;
\draw [shift={(289.33,142.67)}, rotate = 319.24] [color={rgb, 255:red, 128; green, 128; blue, 128 }  ,draw opacity=1 ][line width=0.75]    (10.93,-3.29) .. controls (6.95,-1.4) and (3.31,-0.3) .. (0,0) .. controls (3.31,0.3) and (6.95,1.4) .. (10.93,3.29)   ;

% Text Node
\draw (73.26,40.18) node  [font=\footnotesize] [align=left] {\begin{minipage}[lt]{24.7pt}\setlength\topsep{0pt}
\begin{center}
$q_{n-1}$
\end{center}

\end{minipage}};
% Text Node
\draw (79.1,58.43) node  [font=\small] [align=left] {\begin{minipage}[lt]{32.64pt}\setlength\topsep{0pt}
\begin{center}
$q_n$
\end{center}

\end{minipage}};
% Text Node
\draw (80.38,82) node  [font=\normalsize] [align=left] {\begin{minipage}[lt]{30.89pt}\setlength\topsep{0pt}
\begin{center}
$q_0$
\end{center}

\end{minipage}};
% Text Node
\draw (78.81,105.93) node  [font=\small] [align=left] {\begin{minipage}[lt]{28.75pt}\setlength\topsep{0pt}
\begin{center}
$q_1$
\end{center}

\end{minipage}};
% Text Node
\draw (72.97,124.18) node  [font=\footnotesize] [align=left] {\begin{minipage}[lt]{25.09pt}\setlength\topsep{0pt}
\begin{center}
$q_2$
\end{center}

\end{minipage}};
% Text Node
\draw (247.67,167.17) node  [font=\normalsize] [align=left] {\begin{minipage}[lt]{20.4pt}\setlength\topsep{0pt}
\begin{center}
\Vtextvisiblespace
\end{center}

\end{minipage}};
% Text Node
\draw (277.67,167.17) node  [font=\normalsize] [align=left] {\begin{minipage}[lt]{20.4pt}\setlength\topsep{0pt}
\begin{center}
$0$
\end{center}

\end{minipage}};
% Text Node
\draw (307.67,167.17) node  [font=\normalsize] [align=left] {\begin{minipage}[lt]{20.4pt}\setlength\topsep{0pt}
\begin{center}
$1$
\end{center}

\end{minipage}};
% Text Node
\draw (337.67,167.17) node  [font=\normalsize] [align=left] {\begin{minipage}[lt]{20.4pt}\setlength\topsep{0pt}
\begin{center}
$0$
\end{center}

\end{minipage}};
% Text Node
\draw (367.67,167.17) node  [font=\normalsize] [align=left] {\begin{minipage}[lt]{20.4pt}\setlength\topsep{0pt}
\begin{center}
\Vtextvisiblespace
\end{center}

\end{minipage}};
% Text Node
\draw (397.67,167.17) node  [font=\normalsize] [align=left] {\begin{minipage}[lt]{20.4pt}\setlength\topsep{0pt}
\begin{center}
$...$
\end{center}

\end{minipage}};
% Text Node
\draw (217.67,167.17) node  [font=\normalsize] [align=left] {\begin{minipage}[lt]{20.4pt}\setlength\topsep{0pt}
\begin{center}
$...$
\end{center}

\end{minipage}};
% Text Node
\draw (307.67,205.17) node  [font=\normalsize,color={rgb, 255:red, 128; green, 128; blue, 128 }  ,opacity=1 ] [align=left] {\begin{minipage}[lt]{61.2pt}\setlength\topsep{0pt}
\begin{center}
entrada
\end{center}

\end{minipage}};
% Text Node
\draw (73.33,166.5) node  [font=\normalsize,color={rgb, 255:red, 128; green, 128; blue, 128 }  ,opacity=1 ] [align=left] {\begin{minipage}[lt]{91.57pt}\setlength\topsep{0pt}
\begin{center}
unidad de control
\end{center}

\end{minipage}};
% Text Node
\draw (214.69,41.83) node  [font=\normalsize,color={rgb, 255:red, 128; green, 128; blue, 128 }  ,opacity=1 ] [align=left] {\begin{minipage}[lt]{74.32pt}\setlength\topsep{0pt}
estado actual
\end{minipage}};
% Text Node
\draw (485.02,167.17) node  [font=\normalsize,color={rgb, 255:red, 128; green, 128; blue, 128 }  ,opacity=1 ] [align=left] {\begin{minipage}[lt]{30.35pt}\setlength\topsep{0pt}
cinta
\end{minipage}};
% Text Node
\draw (387.35,105.17) node  [font=\normalsize,color={rgb, 255:red, 128; green, 128; blue, 128 }  ,opacity=1 ] [align=left] {\begin{minipage}[lt]{74.32pt}\setlength\topsep{0pt}
cabezal
\end{minipage}};


\end{tikzpicture}
\end{adjustbox}

\begin{center}
{\small Figura 3.1.: Configuración inicial de una máquina de Turing para la entrada $010$}
\end{center}

\end{frame}

\begin{frame}<presentation:0>{Máquinas de Turing}

\textbf{Definición 3.6} (Configuración de una máquina de Turing)\textbf{.}

Especificaremos la \hle{configuración} de una máquina de Turing mediante dos parámetros:

\begin{itemize}
    \item El estado $q\in Q$ en el que se encuentra la máquina.
    \item La palabra $w \in B^*$ presente en la \emph{cinta} y la posición del \emph{cabezal de lectura}, que representaremos recuadrando el símbolo donde se encuentra el cabezal. Dado que el cabezal se extiende infinitamente a ambos lados, eliminamos los caracteres blancos a la izquierda y a la derecha hasta llegar a la posición del cabezal.
\end{itemize}

Representaremos una configuración mediante $$q\;:\;w$$ con uno de los símbolos de $w$ debidamente recuadrado.

\end{frame}

\begin{frame}<presentation:0>{Máquinas de Turing}

\vspace{5mm}
\begin{adjustbox}{max height={.67\textheight},center}


\tikzset{every picture/.style={line width=0.75pt}} %set default line width to 0.75pt        

\begin{tikzpicture}[x=0.75pt,y=0.75pt,yscale=-1,xscale=1]
%uncomment if require: \path (0,221); %set diagram left start at 0, and has height of 221

%Flowchart: Direct Access Storage [id:dp5745083904317894] 
\draw   (45.92,32.5) -- (91.42,32.5) .. controls (98.18,32.5) and (103.67,54.66) .. (103.67,82) .. controls (103.67,109.34) and (98.18,131.5) .. (91.42,131.5) -- (45.92,131.5)(33.67,82) .. controls (33.67,54.66) and (39.15,32.5) .. (45.92,32.5) .. controls (52.68,32.5) and (58.17,54.66) .. (58.17,82) .. controls (58.17,109.34) and (52.68,131.5) .. (45.92,131.5) .. controls (39.15,131.5) and (33.67,109.34) .. (33.67,82) ;
%Shape: Cube [id:dp9677041927446011] 
\draw   (132.67,13.07) -- (121.59,2) -- (2,2) -- (2,140.43) -- (13.07,151.5) -- (132.67,151.5) -- cycle ; \draw   (2,2) -- (13.07,13.07) -- (132.67,13.07) ; \draw   (13.07,13.07) -- (13.07,151.5) ;
%Shape: Ellipse [id:dp3944720200277354] 
\draw  [fill={rgb, 255:red, 0; green, 0; blue, 0 }  ,fill opacity=1 ] (42.81,82.14) .. controls (42.81,72.36) and (44.22,64.43) .. (45.95,64.43) .. controls (47.69,64.43) and (49.1,72.36) .. (49.1,82.14) .. controls (49.1,91.93) and (47.69,99.86) .. (45.95,99.86) .. controls (44.22,99.86) and (42.81,91.93) .. (42.81,82.14) -- cycle ;
%Straight Lines [id:da7811026207895715] 
\draw    (57.67,95) -- (103.1,95) ;
%Straight Lines [id:da7715257519556658] 
\draw    (54.52,116.86) -- (99.95,116.86) ;
%Straight Lines [id:da7719478908390764] 
\draw    (55.1,47.86) -- (100.52,47.86) ;
%Straight Lines [id:da44878361199349515] 
\draw    (57.67,69) -- (103.1,69) ;
%Flowchart: Merge [id:dp7750110622325217] 
\draw  [fill={rgb, 255:red, 0; green, 0; blue, 0 }  ,fill opacity=1 ] (111.95,76.14) -- (111.95,88) -- (100.24,82.07) -- cycle ;
%Shape: Rectangle [id:dp3100328904459395] 
\draw   (232.67,152.17) -- (262.67,152.17) -- (262.67,182.17) -- (232.67,182.17) -- cycle ;
%Shape: Rectangle [id:dp1692041741743724] 
\draw  [line width=2.25]  (262.67,152.17) -- (292.67,152.17) -- (292.67,182.17) -- (262.67,182.17) -- cycle ;
%Shape: Rectangle [id:dp5247156403937852] 
\draw   (292.67,152.17) -- (322.67,152.17) -- (322.67,182.17) -- (292.67,182.17) -- cycle ;
%Shape: Rectangle [id:dp1939162798122891] 
\draw   (322.67,152.17) -- (352.67,152.17) -- (352.67,182.17) -- (322.67,182.17) -- cycle ;
%Shape: Rectangle [id:dp955944006123826] 
\draw   (352.67,152.17) -- (382.67,152.17) -- (382.67,182.17) -- (352.67,182.17) -- cycle ;
%Straight Lines [id:da2297213302864911] 
\draw    (382.67,152.17) -- (402.67,152.17) ;
%Straight Lines [id:da0017979141787713981] 
\draw    (382.67,182.17) -- (402.67,182.17) ;
%Straight Lines [id:da5196954519515742] 
\draw    (212.67,182.17) -- (232.67,182.17) ;
%Straight Lines [id:da6250621041798603] 
\draw    (212.67,152.17) -- (232.67,152.17) ;
%Straight Lines [id:da4041928331954272] 
\draw [color={rgb, 255:red, 128; green, 128; blue, 128 }  ,draw opacity=1 ]   (262.67,195.17) -- (352.67,195.17) ;
%Straight Lines [id:da40491651879492707] 
\draw [color={rgb, 255:red, 128; green, 128; blue, 128 }  ,draw opacity=1 ]   (262.67,188.17) -- (262.67,195.17) ;
%Straight Lines [id:da9178129384480556] 
\draw [color={rgb, 255:red, 128; green, 128; blue, 128 }  ,draw opacity=1 ]   (352.67,188.17) -- (352.67,195.17) ;
%Shape: Right Angle [id:dp42415196847546355] 
\draw  [line width=1.5]  (132,82) -- (278,82) -- (278,146) ;
%Straight Lines [id:da05158428581441665] 
\draw    (277.67,128) -- (277.67,149.17) ;
\draw [shift={(277.67,152.17)}, rotate = 270] [fill={rgb, 255:red, 0; green, 0; blue, 0 }  ][line width=0.08]  [draw opacity=0] (8.93,-4.29) -- (0,0) -- (8.93,4.29) -- cycle    ;
%Straight Lines [id:da7132164152395672] 
\draw [color={rgb, 255:red, 128; green, 128; blue, 128 }  ,draw opacity=1 ]   (155.33,46) -- (118.18,78.03) ;
\draw [shift={(116.67,79.33)}, rotate = 319.24] [color={rgb, 255:red, 128; green, 128; blue, 128 }  ,draw opacity=1 ][line width=0.75]    (10.93,-3.29) .. controls (6.95,-1.4) and (3.31,-0.3) .. (0,0) .. controls (3.31,0.3) and (6.95,1.4) .. (10.93,3.29)   ;
%Straight Lines [id:da42945985402945897] 
\draw [color={rgb, 255:red, 128; green, 128; blue, 128 }  ,draw opacity=1 ]   (457.33,167.33) -- (418,167.33) ;
\draw [shift={(416,167.33)}, rotate = 360] [color={rgb, 255:red, 128; green, 128; blue, 128 }  ,draw opacity=1 ][line width=0.75]    (10.93,-3.29) .. controls (6.95,-1.4) and (3.31,-0.3) .. (0,0) .. controls (3.31,0.3) and (6.95,1.4) .. (10.93,3.29)   ;
%Straight Lines [id:da3859416199015755] 
\draw [color={rgb, 255:red, 128; green, 128; blue, 128 }  ,draw opacity=1 ]   (328,109.33) -- (290.85,141.36) ;
\draw [shift={(289.33,142.67)}, rotate = 319.24] [color={rgb, 255:red, 128; green, 128; blue, 128 }  ,draw opacity=1 ][line width=0.75]    (10.93,-3.29) .. controls (6.95,-1.4) and (3.31,-0.3) .. (0,0) .. controls (3.31,0.3) and (6.95,1.4) .. (10.93,3.29)   ;

% Text Node
\draw (73.26,40.18) node  [font=\footnotesize] [align=left] {\begin{minipage}[lt]{24.7pt}\setlength\topsep{0pt}
\begin{center}
$q_{n-1}$
\end{center}

\end{minipage}};
% Text Node
\draw (79.1,58.43) node  [font=\small] [align=left] {\begin{minipage}[lt]{32.64pt}\setlength\topsep{0pt}
\begin{center}
$q_n$
\end{center}

\end{minipage}};
% Text Node
\draw (80.38,82) node  [font=\normalsize] [align=left] {\begin{minipage}[lt]{30.89pt}\setlength\topsep{0pt}
\begin{center}
$q_0$
\end{center}

\end{minipage}};
% Text Node
\draw (78.81,105.93) node  [font=\small] [align=left] {\begin{minipage}[lt]{28.75pt}\setlength\topsep{0pt}
\begin{center}
$q_1$
\end{center}

\end{minipage}};
% Text Node
\draw (72.97,124.18) node  [font=\footnotesize] [align=left] {\begin{minipage}[lt]{25.09pt}\setlength\topsep{0pt}
\begin{center}
$q_2$
\end{center}

\end{minipage}};
% Text Node
\draw (247.67,167.17) node  [font=\normalsize] [align=left] {\begin{minipage}[lt]{20.4pt}\setlength\topsep{0pt}
\begin{center}
\Vtextvisiblespace
\end{center}

\end{minipage}};
% Text Node
\draw (277.67,167.17) node  [font=\normalsize] [align=left] {\begin{minipage}[lt]{20.4pt}\setlength\topsep{0pt}
\begin{center}
$0$
\end{center}

\end{minipage}};
% Text Node
\draw (307.67,167.17) node  [font=\normalsize] [align=left] {\begin{minipage}[lt]{20.4pt}\setlength\topsep{0pt}
\begin{center}
$1$
\end{center}

\end{minipage}};
% Text Node
\draw (337.67,167.17) node  [font=\normalsize] [align=left] {\begin{minipage}[lt]{20.4pt}\setlength\topsep{0pt}
\begin{center}
$0$
\end{center}

\end{minipage}};
% Text Node
\draw (367.67,167.17) node  [font=\normalsize] [align=left] {\begin{minipage}[lt]{20.4pt}\setlength\topsep{0pt}
\begin{center}
\Vtextvisiblespace
\end{center}

\end{minipage}};
% Text Node
\draw (397.67,167.17) node  [font=\normalsize] [align=left] {\begin{minipage}[lt]{20.4pt}\setlength\topsep{0pt}
\begin{center}
$...$
\end{center}

\end{minipage}};
% Text Node
\draw (217.67,167.17) node  [font=\normalsize] [align=left] {\begin{minipage}[lt]{20.4pt}\setlength\topsep{0pt}
\begin{center}
$...$
\end{center}

\end{minipage}};
% Text Node
\draw (307.67,205.17) node  [font=\normalsize,color={rgb, 255:red, 128; green, 128; blue, 128 }  ,opacity=1 ] [align=left] {\begin{minipage}[lt]{61.2pt}\setlength\topsep{0pt}
\begin{center}
entrada
\end{center}

\end{minipage}};
% Text Node
\draw (73.33,166.5) node  [font=\normalsize,color={rgb, 255:red, 128; green, 128; blue, 128 }  ,opacity=1 ] [align=left] {\begin{minipage}[lt]{91.57pt}\setlength\topsep{0pt}
\begin{center}
unidad de control
\end{center}

\end{minipage}};
% Text Node
\draw (214.69,41.83) node  [font=\normalsize,color={rgb, 255:red, 128; green, 128; blue, 128 }  ,opacity=1 ] [align=left] {\begin{minipage}[lt]{74.32pt}\setlength\topsep{0pt}
estado actual
\end{minipage}};
% Text Node
\draw (485.02,167.17) node  [font=\normalsize,color={rgb, 255:red, 128; green, 128; blue, 128 }  ,opacity=1 ] [align=left] {\begin{minipage}[lt]{30.35pt}\setlength\topsep{0pt}
cinta
\end{minipage}};
% Text Node
\draw (387.35,105.17) node  [font=\normalsize,color={rgb, 255:red, 128; green, 128; blue, 128 }  ,opacity=1 ] [align=left] {\begin{minipage}[lt]{74.32pt}\setlength\topsep{0pt}
cabezal
\end{minipage}};


\draw (400.69,41.83) node  [font=\normalsize,color={rgb, 255:red, 0; green, 0; blue, 0 }  ,opacity=1 ] [align=left] {\begin{minipage}[lt]{74.32pt}\setlength\topsep{0pt}
\Large{$q_0\;:\;\fbox{$0$}\:1\:0$}
\end{minipage}};

\end{tikzpicture}
\end{adjustbox}

\begin{center}
{\small Figura 3.1.: Configuración inicial de una máquina de Turing para la entrada $010$}
\end{center}

\end{frame}

\begin{frame}<presentation:0>{Máquinas de Turing}

\textbf{Definición 3.7} (Proceso de cálculo de una máquina de Turing)\textbf{.}

$M=(Q, A, B, \delta, q_0, \Vtextvisiblespace\:, F)$ y \hle{entrada inicial} $u = u_1\:u_2\:...\:u_n\in A^*$

\begin{itemize}
    \item \hle{Configuración inicial:} $$q_0\;:\;\fbox{$u_1$}\:u_2\:...\:u_n$$
    \item \hle{Paso de cálculo:} $$q\;:\;b_{-n}\:...\:b_{-2}\:b_{-1}\:\fbox{$b_0$}\:b_1\:b_2\:...\:b_m$$

    \begin{itemize}
        \item Si $\delta(q,b_0)=\emptyset$, \ul{parar}.
        \item Si $\delta(q,b_0)=(p,c,M)$, \ul{continuar}:
        $$
        \begin{matrix}
            M=I&\Rightarrow&p\;:\;b_{-n}\:...\:b_{-2}\:\fbox{$b_{-1}$}\:c\:b_1\:b_2\:...\:b_m\\
            M=D&\Rightarrow&p\;:\;b_{-n}\:...\:b_{-2}\:b_{-1}\:c\:\fbox{$b_1$}\:b_2\:...\:b_m\\
            M=S&\Rightarrow&p\;:\;b_{-n}\:...\:b_{-2}\:b_{-1}\:\fbox{$c$}\:b_1\:b_2\:...\:b_m
        \end{matrix}
    $$
    \end{itemize}
\end{itemize}

\end{frame}

\begin{frame}{Máquinas de Turing}

\textbf{Ejemplo 3.1}\textbf{.}
$$
    M_{\#\text{a}>\#\text{b}}=(\{q_0, q_\text{a}, q_\text{b}, q_\text{R}, q_\text{F}\}, \{\text{a}, \text{b}\}, \{\text{a}, \text{b}, \text{X}, \Vtextvisiblespace\}, \delta, q_0, \Vtextvisiblespace\:,\{q_\text{F}\})
$$
$$
    \begin{matrix}
        \delta(q_0, \text{a}) = (q_\text{a},\text{X},D) & \delta(q_0, \text{b}) = (q_\text{b},\text{X},D) & \delta(q_0, \text{X}) = (q_0,\text{X},D) &  \\
        \delta(q_\text{a}, \text{a}) = (q_\text{a},\text{a},D) & \delta(q_\text{a}, \text{b}) = (q_\text{R},\text{X},I) & \delta(q_\text{a}, \text{X}) = (q_\text{a},\text{X},D) & \delta(q_\text{a}, \Vtextvisiblespace) = (q_\text{F},\Vtextvisiblespace,S) \\
        \delta(q_\text{b}, \text{a}) = (q_\text{R},\text{X},I) & \delta(q_\text{b}, \text{b}) = (q_\text{b},\text{b},D) & \delta(q_\text{b}, \text{X}) = (q_\text{b},\text{X},D) &  \\
        \delta(q_\text{R}, \text{a}) = (q_\text{R},\text{a},I) & \delta(q_\text{R}, \text{b}) = (q_\text{R},\text{b},I) & \delta(q_\text{R}, \text{X}) = (q_\text{R},\text{X},I) & \delta(q_\text{R}, \Vtextvisiblespace) = (q_0,\Vtextvisiblespace,D) \\
    \end{matrix}
$$

\vspace{5mm}

\begin{adjustbox}{max height={.67\textheight},center}


\tikzset{every picture/.style={line width=0.75pt}} %set default line width to 0.75pt        

\begin{tikzpicture}[x=0.75pt,y=0.75pt,yscale=-1,xscale=1]
%uncomment if require: \path (0,446); %set diagram left start at 0, and has height of 446

%Shape: Triangle [id:dp9549836358247636] 
\draw  [fill={rgb, 255:red, 0; green, 0; blue, 0 }  ,fill opacity=1 ] (134.4,85) -- (141.6,82.2) -- (141.6,87.8) -- cycle ;
%Shape: Circle [id:dp39324764126664746] 
\draw   (101.2,85) .. controls (101.2,75) and (109.31,66.9) .. (119.3,66.9) .. controls (129.3,66.9) and (137.4,75) .. (137.4,85) .. controls (137.4,95) and (129.3,103.1) .. (119.3,103.1) .. controls (109.31,103.1) and (101.2,95) .. (101.2,85) -- cycle ;
%Straight Lines [id:da279972863636345] 
\draw    (379.6,70.2) -- (248.8,70.2) -- (189.6,70.2) ;
%Straight Lines [id:da8063028764122899] 
\draw    (379.2,99.6) -- (189.4,99.6) ;
%Straight Lines [id:da6194955219154272] 
\draw    (210,70.4) -- (210,99.6) ;
%Straight Lines [id:da9291793988668682] 
\draw    (239.6,70.8) -- (239.6,100) ;
%Straight Lines [id:da1753997056206802] 
\draw    (269.2,70.4) -- (269.2,99.6) ;
%Straight Lines [id:da7326138590955908] 
\draw    (299.2,70.4) -- (299.2,99.6) ;
%Straight Lines [id:da8827634409952425] 
\draw    (329.2,70.4) -- (329.2,99.6) ;
%Straight Lines [id:da08240500809926843] 
\draw    (359.2,70.8) -- (359.2,100) ;
%Shape: Rectangle [id:dp8505881042891044] 
\draw  [line width=2.25]  (269.2,99.6) -- (239.6,99.6) -- (239.6,70.4) -- (269.2,70.4) -- cycle ;
%Straight Lines [id:da2295782410185434] 
\draw    (254.8,117.75) -- (254.8,103.55) ;
\draw [shift={(254.8,100.55)}, rotate = 90] [fill={rgb, 255:red, 0; green, 0; blue, 0 }  ][line width=0.08]  [draw opacity=0] (8.93,-4.29) -- (0,0) -- (8.93,4.29) -- cycle    ;
%Straight Lines [id:da2131860495885145] 
\draw [line width=1.5]    (254.8,123.4) -- (254.8,106.15) ;


% Text Node
\draw (107.4,79.2) node [anchor=north west][inner sep=0.75pt]  [font=\small] [align=left] {\begin{minipage}[lt]{18.04pt}\setlength\topsep{0pt}
\begin{center}
$q_0$
\end{center}

\end{minipage}};
% Text Node
\draw (224.8,85.4) node   [align=left] {\begin{minipage}[lt]{20.13pt}\setlength\topsep{0pt}
\begin{center}
\Vtextvisiblespace
\end{center}

\end{minipage}};
% Text Node
\draw (254.4,86.4) node   [align=left] {\begin{minipage}[lt]{20.13pt}\setlength\topsep{0pt}
\begin{center}
a
\end{center}

\end{minipage}};
% Text Node
\draw (284.4,85) node   [align=left] {\begin{minipage}[lt]{20.13pt}\setlength\topsep{0pt}
\begin{center}
b
\end{center}

\end{minipage}};
% Text Node
\draw (314,86.4) node   [align=left] {\begin{minipage}[lt]{20.13pt}\setlength\topsep{0pt}
\begin{center}
a
\end{center}

\end{minipage}};
% Text Node
\draw (344.4,85.4) node   [align=left] {\begin{minipage}[lt]{20.13pt}\setlength\topsep{0pt}
\begin{center}
\Vtextvisiblespace
\end{center}

\end{minipage}};
% Text Node
\draw (374,85.4) node   [align=left] {\begin{minipage}[lt]{20.13pt}\setlength\topsep{0pt}
\begin{center}
...
\end{center}

\end{minipage}};
% Text Node
\draw (194.2,86) node   [align=left] {\begin{minipage}[lt]{20.13pt}\setlength\topsep{0pt}
\begin{center}
...
\end{center}

\end{minipage}};


\end{tikzpicture}

\end{adjustbox}

\end{frame}


\begin{frame}{Máquinas de Turing}

\textbf{Ejemplo 3.1}\textbf{.}
$$
    M_{\#\text{a}>\#\text{b}}=(\{q_0, q_\text{a}, q_\text{b}, q_\text{R}, q_\text{F}\}, \{\text{a}, \text{b}\}, \{\text{a}, \text{b}, \text{X}, \Vtextvisiblespace\}, \delta, q_0, \Vtextvisiblespace\:,\{q_\text{F}\})
$$
$$
    \begin{matrix}
        \underline{\delta(q_0, \text{a}) = (q_\text{a},\text{X},D)} & \delta(q_0, \text{b}) = (q_\text{b},\text{X},D) & \delta(q_0, \text{X}) = (q_0,\text{X},D) &  \\
        \delta(q_\text{a}, \text{a}) = (q_\text{a},\text{a},D) & \delta(q_\text{a}, \text{b}) = (q_\text{R},\text{X},I) & \delta(q_\text{a}, \text{X}) = (q_\text{a},\text{X},D) & \delta(q_\text{a}, \Vtextvisiblespace) = (q_\text{F},\Vtextvisiblespace,S) \\
        \delta(q_\text{b}, \text{a}) = (q_\text{R},\text{X},I) & \delta(q_\text{b}, \text{b}) = (q_\text{b},\text{b},D) & \delta(q_\text{b}, \text{X}) = (q_\text{b},\text{X},D) &  \\
        \delta(q_\text{R}, \text{a}) = (q_\text{R},\text{a},I) & \delta(q_\text{R}, \text{b}) = (q_\text{R},\text{b},I) & \delta(q_\text{R}, \text{X}) = (q_\text{R},\text{X},I) & \delta(q_\text{R}, \Vtextvisiblespace) = (q_0,\Vtextvisiblespace,D) \\
    \end{matrix}
$$

\vspace{5mm}

\begin{adjustbox}{max height={.67\textheight},center}


\tikzset{every picture/.style={line width=0.75pt}} %set default line width to 0.75pt        

\begin{tikzpicture}[x=0.75pt,y=0.75pt,yscale=-1,xscale=1]
%uncomment if require: \path (0,446); %set diagram left start at 0, and has height of 446

%Shape: Triangle [id:dp9549836358247636] 
\draw  [fill={rgb, 255:red, 0; green, 0; blue, 0 }  ,fill opacity=1 ] (134.4,85) -- (141.6,82.2) -- (141.6,87.8) -- cycle ;
%Shape: Circle [id:dp39324764126664746] 
\draw   (101.2,85) .. controls (101.2,75) and (109.31,66.9) .. (119.3,66.9) .. controls (129.3,66.9) and (137.4,75) .. (137.4,85) .. controls (137.4,95) and (129.3,103.1) .. (119.3,103.1) .. controls (109.31,103.1) and (101.2,95) .. (101.2,85) -- cycle ;
%Straight Lines [id:da279972863636345] 
\draw    (379.6,70.2) -- (248.8,70.2) -- (189.6,70.2) ;
%Straight Lines [id:da8063028764122899] 
\draw    (379.2,99.6) -- (189.4,99.6) ;
%Straight Lines [id:da6194955219154272] 
\draw    (210,70.4) -- (210,99.6) ;
%Straight Lines [id:da9291793988668682] 
\draw    (239.6,70.8) -- (239.6,100) ;
%Straight Lines [id:da1753997056206802] 
\draw    (269.2,70.4) -- (269.2,99.6) ;
%Straight Lines [id:da7326138590955908] 
\draw    (299.2,70.4) -- (299.2,99.6) ;
%Straight Lines [id:da8827634409952425] 
\draw    (329.2,70.4) -- (329.2,99.6) ;
%Straight Lines [id:da08240500809926843] 
\draw    (359.2,70.8) -- (359.2,100) ;
%Shape: Rectangle [id:dp8505881042891044] 
\draw  [line width=2.25]  (269.2,99.6) -- (239.6,99.6) -- (239.6,70.4) -- (269.2,70.4) -- cycle ;
%Straight Lines [id:da2295782410185434] 
\draw    (254.8,117.75) -- (254.8,103.55) ;
\draw [shift={(254.8,100.55)}, rotate = 90] [fill={rgb, 255:red, 0; green, 0; blue, 0 }  ][line width=0.08]  [draw opacity=0] (8.93,-4.29) -- (0,0) -- (8.93,4.29) -- cycle    ;
%Straight Lines [id:da2131860495885145] 
\draw [line width=1.5]    (254.8,123.4) -- (254.8,106.15) ;


% Text Node
\draw (107.4,79.2) node [anchor=north west][inner sep=0.75pt]  [font=\small] [align=left] {\begin{minipage}[lt]{18.04pt}\setlength\topsep{0pt}
\begin{center}
$q_0$
\end{center}

\end{minipage}};
% Text Node
\draw (224.8,85.4) node   [align=left] {\begin{minipage}[lt]{20.13pt}\setlength\topsep{0pt}
\begin{center}
\Vtextvisiblespace
\end{center}

\end{minipage}};
% Text Node
\draw (254.4,86.4) node   [align=left] {\begin{minipage}[lt]{20.13pt}\setlength\topsep{0pt}
\begin{center}
a
\end{center}

\end{minipage}};
% Text Node
\draw (284.4,85) node   [align=left] {\begin{minipage}[lt]{20.13pt}\setlength\topsep{0pt}
\begin{center}
b
\end{center}

\end{minipage}};
% Text Node
\draw (314,86.4) node   [align=left] {\begin{minipage}[lt]{20.13pt}\setlength\topsep{0pt}
\begin{center}
a
\end{center}

\end{minipage}};
% Text Node
\draw (344.4,85.4) node   [align=left] {\begin{minipage}[lt]{20.13pt}\setlength\topsep{0pt}
\begin{center}
\Vtextvisiblespace
\end{center}

\end{minipage}};
% Text Node
\draw (374,85.4) node   [align=left] {\begin{minipage}[lt]{20.13pt}\setlength\topsep{0pt}
\begin{center}
...
\end{center}

\end{minipage}};
% Text Node
\draw (194.2,86) node   [align=left] {\begin{minipage}[lt]{20.13pt}\setlength\topsep{0pt}
\begin{center}
...
\end{center}

\end{minipage}};


\end{tikzpicture}

\end{adjustbox}

\end{frame}


\begin{frame}{Máquinas de Turing}

\textbf{Ejemplo 3.1}\textbf{.}
$$
    M_{\#\text{a}>\#\text{b}}=(\{q_0, q_\text{a}, q_\text{b}, q_\text{R}, q_\text{F}\}, \{\text{a}, \text{b}\}, \{\text{a}, \text{b}, \text{X}, \Vtextvisiblespace\}, \delta, q_0, \Vtextvisiblespace\:,\{q_\text{F}\})
$$
$$
    \begin{matrix}
        \delta(q_0, \text{a}) = (q_\text{a},\text{X},D) & \delta(q_0, \text{b}) = (q_\text{b},\text{X},D) & \delta(q_0, \text{X}) = (q_0,\text{X},D) &  \\
        \delta(q_\text{a}, \text{a}) = (q_\text{a},\text{a},D) & \delta(q_\text{a}, \text{b}) = (q_\text{R},\text{X},I) & \delta(q_\text{a}, \text{X}) = (q_\text{a},\text{X},D) & \delta(q_\text{a}, \Vtextvisiblespace) = (q_\text{F},\Vtextvisiblespace,S) \\
        \delta(q_\text{b}, \text{a}) = (q_\text{R},\text{X},I) & \delta(q_\text{b}, \text{b}) = (q_\text{b},\text{b},D) & \delta(q_\text{b}, \text{X}) = (q_\text{b},\text{X},D) &  \\
        \delta(q_\text{R}, \text{a}) = (q_\text{R},\text{a},I) & \delta(q_\text{R}, \text{b}) = (q_\text{R},\text{b},I) & \delta(q_\text{R}, \text{X}) = (q_\text{R},\text{X},I) & \delta(q_\text{R}, \Vtextvisiblespace) = (q_0,\Vtextvisiblespace,D) \\
    \end{matrix}
$$

\vspace{5mm}

\begin{adjustbox}{max height={.67\textheight},center}


\tikzset{every picture/.style={line width=0.75pt}} %set default line width to 0.75pt        

\begin{tikzpicture}[x=0.75pt,y=0.75pt,yscale=-1,xscale=1]
%uncomment if require: \path (0,446); %set diagram left start at 0, and has height of 446

%Shape: Triangle [id:dp9549836358247636] 
\draw  [fill={rgb, 255:red, 0; green, 0; blue, 0 }  ,fill opacity=1 ] (134.4,85) -- (141.6,82.2) -- (141.6,87.8) -- cycle ;
%Shape: Circle [id:dp39324764126664746] 
\draw   (101.2,85) .. controls (101.2,75) and (109.31,66.9) .. (119.3,66.9) .. controls (129.3,66.9) and (137.4,75) .. (137.4,85) .. controls (137.4,95) and (129.3,103.1) .. (119.3,103.1) .. controls (109.31,103.1) and (101.2,95) .. (101.2,85) -- cycle ;
%Straight Lines [id:da279972863636345] 
\draw    (379.6,70.2) -- (248.8,70.2) -- (189.6,70.2) ;
%Straight Lines [id:da8063028764122899] 
\draw    (379.2,99.6) -- (189.4,99.6) ;
%Straight Lines [id:da6194955219154272] 
\draw    (210,70.4) -- (210,99.6) ;
%Straight Lines [id:da9291793988668682] 
\draw    (239.6,70.8) -- (239.6,100) ;
%Straight Lines [id:da1753997056206802] 
\draw    (269.2,70.4) -- (269.2,99.6) ;
%Straight Lines [id:da7326138590955908] 
\draw    (299.2,70.4) -- (299.2,99.6) ;
%Straight Lines [id:da8827634409952425] 
\draw    (329.2,70.4) -- (329.2,99.6) ;
%Straight Lines [id:da08240500809926843] 
\draw    (359.2,70.8) -- (359.2,100) ;
%Shape: Rectangle [id:dp8505881042891044] 
\draw  [line width=2.25]  (299.2,99.6) -- (269.6,99.6) -- (269.6,70.4) -- (299.2,70.4) -- cycle ;
%Straight Lines [id:da2295782410185434] 
\draw    (284.8,117.75) -- (284.8,103.55) ;
\draw [shift={(284.8,100.55)}, rotate = 90] [fill={rgb, 255:red, 0; green, 0; blue, 0 }  ][line width=0.08]  [draw opacity=0] (8.93,-4.29) -- (0,0) -- (8.93,4.29) -- cycle    ;
%Straight Lines [id:da2131860495885145] 
\draw [line width=1.5]    (284.8,123.4) -- (284.8,106.15) ;


% Text Node
\draw (107.4,79.2) node [anchor=north west][inner sep=0.75pt]  [font=\small] [align=left] {\begin{minipage}[lt]{18.04pt}\setlength\topsep{0pt}
\begin{center}
$q_\text{a}$
\end{center}

\end{minipage}};
% Text Node
\draw (224.8,85.4) node   [align=left] {\begin{minipage}[lt]{20.13pt}\setlength\topsep{0pt}
\begin{center}
\Vtextvisiblespace
\end{center}

\end{minipage}};
% Text Node
\draw (254.4,85) node   [align=left] {\begin{minipage}[lt]{20.13pt}\setlength\topsep{0pt}
\begin{center}
X
\end{center}

\end{minipage}};
% Text Node
\draw (284.4,85) node   [align=left] {\begin{minipage}[lt]{20.13pt}\setlength\topsep{0pt}
\begin{center}
b
\end{center}

\end{minipage}};
% Text Node
\draw (314,86.4) node   [align=left] {\begin{minipage}[lt]{20.13pt}\setlength\topsep{0pt}
\begin{center}
a
\end{center}

\end{minipage}};
% Text Node
\draw (344.4,85.4) node   [align=left] {\begin{minipage}[lt]{20.13pt}\setlength\topsep{0pt}
\begin{center}
\Vtextvisiblespace
\end{center}

\end{minipage}};
% Text Node
\draw (374,85.4) node   [align=left] {\begin{minipage}[lt]{20.13pt}\setlength\topsep{0pt}
\begin{center}
...
\end{center}

\end{minipage}};
% Text Node
\draw (194.2,86) node   [align=left] {\begin{minipage}[lt]{20.13pt}\setlength\topsep{0pt}
\begin{center}
...
\end{center}

\end{minipage}};


\end{tikzpicture}

\end{adjustbox}

\end{frame}


\begin{frame}{Máquinas de Turing}

\textbf{Ejemplo 3.1}\textbf{.}
$$
    M_{\#\text{a}>\#\text{b}}=(\{q_0, q_\text{a}, q_\text{b}, q_\text{R}, q_\text{F}\}, \{\text{a}, \text{b}\}, \{\text{a}, \text{b}, \text{X}, \Vtextvisiblespace\}, \delta, q_0, \Vtextvisiblespace\:,\{q_\text{F}\})
$$
$$
    \begin{matrix}
        \delta(q_0, \text{a}) = (q_\text{a},\text{X},D) & \delta(q_0, \text{b}) = (q_\text{b},\text{X},D) & \delta(q_0, \text{X}) = (q_0,\text{X},D) &  \\
        \delta(q_\text{a}, \text{a}) = (q_\text{a},\text{a},D) & \underline{\delta(q_\text{a}, \text{b}) = (q_\text{R},\text{X},I)} & \delta(q_\text{a}, \text{X}) = (q_\text{a},\text{X},D) & \delta(q_\text{a}, \Vtextvisiblespace) = (q_\text{F},\Vtextvisiblespace,S) \\
        \delta(q_\text{b}, \text{a}) = (q_\text{R},\text{X},I) & \delta(q_\text{b}, \text{b}) = (q_\text{b},\text{b},D) & \delta(q_\text{b}, \text{X}) = (q_\text{b},\text{X},D) &  \\
        \delta(q_\text{R}, \text{a}) = (q_\text{R},\text{a},I) & \delta(q_\text{R}, \text{b}) = (q_\text{R},\text{b},I) & \delta(q_\text{R}, \text{X}) = (q_\text{R},\text{X},I) & \delta(q_\text{R}, \Vtextvisiblespace) = (q_0,\Vtextvisiblespace,D) \\
    \end{matrix}
$$

\vspace{5mm}

\begin{adjustbox}{max height={.67\textheight},center}


\tikzset{every picture/.style={line width=0.75pt}} %set default line width to 0.75pt        

\begin{tikzpicture}[x=0.75pt,y=0.75pt,yscale=-1,xscale=1]
%uncomment if require: \path (0,446); %set diagram left start at 0, and has height of 446

%Shape: Triangle [id:dp9549836358247636] 
\draw  [fill={rgb, 255:red, 0; green, 0; blue, 0 }  ,fill opacity=1 ] (134.4,85) -- (141.6,82.2) -- (141.6,87.8) -- cycle ;
%Shape: Circle [id:dp39324764126664746] 
\draw   (101.2,85) .. controls (101.2,75) and (109.31,66.9) .. (119.3,66.9) .. controls (129.3,66.9) and (137.4,75) .. (137.4,85) .. controls (137.4,95) and (129.3,103.1) .. (119.3,103.1) .. controls (109.31,103.1) and (101.2,95) .. (101.2,85) -- cycle ;
%Straight Lines [id:da279972863636345] 
\draw    (379.6,70.2) -- (248.8,70.2) -- (189.6,70.2) ;
%Straight Lines [id:da8063028764122899] 
\draw    (379.2,99.6) -- (189.4,99.6) ;
%Straight Lines [id:da6194955219154272] 
\draw    (210,70.4) -- (210,99.6) ;
%Straight Lines [id:da9291793988668682] 
\draw    (239.6,70.8) -- (239.6,100) ;
%Straight Lines [id:da1753997056206802] 
\draw    (269.2,70.4) -- (269.2,99.6) ;
%Straight Lines [id:da7326138590955908] 
\draw    (299.2,70.4) -- (299.2,99.6) ;
%Straight Lines [id:da8827634409952425] 
\draw    (329.2,70.4) -- (329.2,99.6) ;
%Straight Lines [id:da08240500809926843] 
\draw    (359.2,70.8) -- (359.2,100) ;
%Shape: Rectangle [id:dp8505881042891044] 
\draw  [line width=2.25]  (299.2,99.6) -- (269.6,99.6) -- (269.6,70.4) -- (299.2,70.4) -- cycle ;
%Straight Lines [id:da2295782410185434] 
\draw    (284.8,117.75) -- (284.8,103.55) ;
\draw [shift={(284.8,100.55)}, rotate = 90] [fill={rgb, 255:red, 0; green, 0; blue, 0 }  ][line width=0.08]  [draw opacity=0] (8.93,-4.29) -- (0,0) -- (8.93,4.29) -- cycle    ;
%Straight Lines [id:da2131860495885145] 
\draw [line width=1.5]    (284.8,123.4) -- (284.8,106.15) ;


% Text Node
\draw (107.4,79.2) node [anchor=north west][inner sep=0.75pt]  [font=\small] [align=left] {\begin{minipage}[lt]{18.04pt}\setlength\topsep{0pt}
\begin{center}
$q_\text{a}$
\end{center}

\end{minipage}};
% Text Node
\draw (224.8,85.4) node   [align=left] {\begin{minipage}[lt]{20.13pt}\setlength\topsep{0pt}
\begin{center}
\Vtextvisiblespace
\end{center}

\end{minipage}};
% Text Node
\draw (254.4,85) node   [align=left] {\begin{minipage}[lt]{20.13pt}\setlength\topsep{0pt}
\begin{center}
X
\end{center}

\end{minipage}};
% Text Node
\draw (284.4,85) node   [align=left] {\begin{minipage}[lt]{20.13pt}\setlength\topsep{0pt}
\begin{center}
b
\end{center}

\end{minipage}};
% Text Node
\draw (314,86.4) node   [align=left] {\begin{minipage}[lt]{20.13pt}\setlength\topsep{0pt}
\begin{center}
a
\end{center}

\end{minipage}};
% Text Node
\draw (344.4,85.4) node   [align=left] {\begin{minipage}[lt]{20.13pt}\setlength\topsep{0pt}
\begin{center}
\Vtextvisiblespace
\end{center}

\end{minipage}};
% Text Node
\draw (374,85.4) node   [align=left] {\begin{minipage}[lt]{20.13pt}\setlength\topsep{0pt}
\begin{center}
...
\end{center}

\end{minipage}};
% Text Node
\draw (194.2,86) node   [align=left] {\begin{minipage}[lt]{20.13pt}\setlength\topsep{0pt}
\begin{center}
...
\end{center}

\end{minipage}};


\end{tikzpicture}

\end{adjustbox}

\end{frame}


\begin{frame}{Máquinas de Turing}

\textbf{Ejemplo 3.1}\textbf{.}
$$
    M_{\#\text{a}>\#\text{b}}=(\{q_0, q_\text{a}, q_\text{b}, q_\text{R}, q_\text{F}\}, \{\text{a}, \text{b}\}, \{\text{a}, \text{b}, \text{X}, \Vtextvisiblespace\}, \delta, q_0, \Vtextvisiblespace\:,\{q_\text{F}\})
$$
$$
    \begin{matrix}
        \delta(q_0, \text{a}) = (q_\text{a},\text{X},D) & \delta(q_0, \text{b}) = (q_\text{b},\text{X},D) & \delta(q_0, \text{X}) = (q_0,\text{X},D) &  \\
        \delta(q_\text{a}, \text{a}) = (q_\text{a},\text{a},D) & \delta(q_\text{a}, \text{b}) = (q_\text{R},\text{X},I) & \delta(q_\text{a}, \text{X}) = (q_\text{a},\text{X},D) & \delta(q_\text{a}, \Vtextvisiblespace) = (q_\text{F},\Vtextvisiblespace,S) \\
        \delta(q_\text{b}, \text{a}) = (q_\text{R},\text{X},I) & \delta(q_\text{b}, \text{b}) = (q_\text{b},\text{b},D) & \delta(q_\text{b}, \text{X}) = (q_\text{b},\text{X},D) &  \\
        \delta(q_\text{R}, \text{a}) = (q_\text{R},\text{a},I) & \delta(q_\text{R}, \text{b}) = (q_\text{R},\text{b},I) & \delta(q_\text{R}, \text{X}) = (q_\text{R},\text{X},I) & \delta(q_\text{R}, \Vtextvisiblespace) = (q_0,\Vtextvisiblespace,D) \\
    \end{matrix}
$$

\vspace{5mm}

\begin{adjustbox}{max height={.67\textheight},center}


\tikzset{every picture/.style={line width=0.75pt}} %set default line width to 0.75pt        

\begin{tikzpicture}[x=0.75pt,y=0.75pt,yscale=-1,xscale=1]
%uncomment if require: \path (0,446); %set diagram left start at 0, and has height of 446

%Shape: Triangle [id:dp9549836358247636] 
\draw  [fill={rgb, 255:red, 0; green, 0; blue, 0 }  ,fill opacity=1 ] (134.4,85) -- (141.6,82.2) -- (141.6,87.8) -- cycle ;
%Shape: Circle [id:dp39324764126664746] 
\draw   (101.2,85) .. controls (101.2,75) and (109.31,66.9) .. (119.3,66.9) .. controls (129.3,66.9) and (137.4,75) .. (137.4,85) .. controls (137.4,95) and (129.3,103.1) .. (119.3,103.1) .. controls (109.31,103.1) and (101.2,95) .. (101.2,85) -- cycle ;
%Straight Lines [id:da279972863636345] 
\draw    (379.6,70.2) -- (248.8,70.2) -- (189.6,70.2) ;
%Straight Lines [id:da8063028764122899] 
\draw    (379.2,99.6) -- (189.4,99.6) ;
%Straight Lines [id:da6194955219154272] 
\draw    (210,70.4) -- (210,99.6) ;
%Straight Lines [id:da9291793988668682] 
\draw    (239.6,70.8) -- (239.6,100) ;
%Straight Lines [id:da1753997056206802] 
\draw    (269.2,70.4) -- (269.2,99.6) ;
%Straight Lines [id:da7326138590955908] 
\draw    (299.2,70.4) -- (299.2,99.6) ;
%Straight Lines [id:da8827634409952425] 
\draw    (329.2,70.4) -- (329.2,99.6) ;
%Straight Lines [id:da08240500809926843] 
\draw    (359.2,70.8) -- (359.2,100) ;
%Shape: Rectangle [id:dp8505881042891044] 
\draw  [line width=2.25]  (269.2,99.6) -- (239.6,99.6) -- (239.6,70.4) -- (269.2,70.4) -- cycle ;
%Straight Lines [id:da2295782410185434] 
\draw    (254.8,117.75) -- (254.8,103.55) ;
\draw [shift={(254.8,100.55)}, rotate = 90] [fill={rgb, 255:red, 0; green, 0; blue, 0 }  ][line width=0.08]  [draw opacity=0] (8.93,-4.29) -- (0,0) -- (8.93,4.29) -- cycle    ;
%Straight Lines [id:da2131860495885145] 
\draw [line width=1.5]    (254.8,123.4) -- (254.8,106.15) ;


% Text Node
\draw (107.4,79.2) node [anchor=north west][inner sep=0.75pt]  [font=\small] [align=left] {\begin{minipage}[lt]{18.04pt}\setlength\topsep{0pt}
\begin{center}
$q_\text{R}$
\end{center}

\end{minipage}};
% Text Node
\draw (224.8,85.4) node   [align=left] {\begin{minipage}[lt]{20.13pt}\setlength\topsep{0pt}
\begin{center}
\Vtextvisiblespace
\end{center}

\end{minipage}};
% Text Node
\draw (254.4,85) node   [align=left] {\begin{minipage}[lt]{20.13pt}\setlength\topsep{0pt}
\begin{center}
X
\end{center}

\end{minipage}};
% Text Node
\draw (284.4,85) node   [align=left] {\begin{minipage}[lt]{20.13pt}\setlength\topsep{0pt}
\begin{center}
X
\end{center}

\end{minipage}};
% Text Node
\draw (314,86.4) node   [align=left] {\begin{minipage}[lt]{20.13pt}\setlength\topsep{0pt}
\begin{center}
a
\end{center}

\end{minipage}};
% Text Node
\draw (344.4,85.4) node   [align=left] {\begin{minipage}[lt]{20.13pt}\setlength\topsep{0pt}
\begin{center}
\Vtextvisiblespace
\end{center}

\end{minipage}};
% Text Node
\draw (374,85.4) node   [align=left] {\begin{minipage}[lt]{20.13pt}\setlength\topsep{0pt}
\begin{center}
...
\end{center}

\end{minipage}};
% Text Node
\draw (194.2,86) node   [align=left] {\begin{minipage}[lt]{20.13pt}\setlength\topsep{0pt}
\begin{center}
...
\end{center}

\end{minipage}};


\end{tikzpicture}

\end{adjustbox}

\end{frame}


\begin{frame}{Máquinas de Turing}

\textbf{Ejemplo 3.1}\textbf{.}
$$
    M_{\#\text{a}>\#\text{b}}=(\{q_0, q_\text{a}, q_\text{b}, q_\text{R}, q_\text{F}\}, \{\text{a}, \text{b}\}, \{\text{a}, \text{b}, \text{X}, \Vtextvisiblespace\}, \delta, q_0, \Vtextvisiblespace\:,\{q_\text{F}\})
$$
$$
    \begin{matrix}
        \delta(q_0, \text{a}) = (q_\text{a},\text{X},D) & \delta(q_0, \text{b}) = (q_\text{b},\text{X},D) & \delta(q_0, \text{X}) = (q_0,\text{X},D) &  \\
        \delta(q_\text{a}, \text{a}) = (q_\text{a},\text{a},D) & \delta(q_\text{a}, \text{b}) = (q_\text{R},\text{X},I) & \delta(q_\text{a}, \text{X}) = (q_\text{a},\text{X},D) & \delta(q_\text{a}, \Vtextvisiblespace) = (q_\text{F},\Vtextvisiblespace,S) \\
        \delta(q_\text{b}, \text{a}) = (q_\text{R},\text{X},I) & \delta(q_\text{b}, \text{b}) = (q_\text{b},\text{b},D) & \delta(q_\text{b}, \text{X}) = (q_\text{b},\text{X},D) &  \\
        \delta(q_\text{R}, \text{a}) = (q_\text{R},\text{a},I) & \delta(q_\text{R}, \text{b}) = (q_\text{R},\text{b},I) & \underline{\delta(q_\text{R}, \text{X}) = (q_\text{R},\text{X},I)} & \delta(q_\text{R}, \Vtextvisiblespace) = (q_0,\Vtextvisiblespace,D) \\
    \end{matrix}
$$

\vspace{5mm}

\begin{adjustbox}{max height={.67\textheight},center}


\tikzset{every picture/.style={line width=0.75pt}} %set default line width to 0.75pt        

\begin{tikzpicture}[x=0.75pt,y=0.75pt,yscale=-1,xscale=1]
%uncomment if require: \path (0,446); %set diagram left start at 0, and has height of 446

%Shape: Triangle [id:dp9549836358247636] 
\draw  [fill={rgb, 255:red, 0; green, 0; blue, 0 }  ,fill opacity=1 ] (134.4,85) -- (141.6,82.2) -- (141.6,87.8) -- cycle ;
%Shape: Circle [id:dp39324764126664746] 
\draw   (101.2,85) .. controls (101.2,75) and (109.31,66.9) .. (119.3,66.9) .. controls (129.3,66.9) and (137.4,75) .. (137.4,85) .. controls (137.4,95) and (129.3,103.1) .. (119.3,103.1) .. controls (109.31,103.1) and (101.2,95) .. (101.2,85) -- cycle ;
%Straight Lines [id:da279972863636345] 
\draw    (379.6,70.2) -- (248.8,70.2) -- (189.6,70.2) ;
%Straight Lines [id:da8063028764122899] 
\draw    (379.2,99.6) -- (189.4,99.6) ;
%Straight Lines [id:da6194955219154272] 
\draw    (210,70.4) -- (210,99.6) ;
%Straight Lines [id:da9291793988668682] 
\draw    (239.6,70.8) -- (239.6,100) ;
%Straight Lines [id:da1753997056206802] 
\draw    (269.2,70.4) -- (269.2,99.6) ;
%Straight Lines [id:da7326138590955908] 
\draw    (299.2,70.4) -- (299.2,99.6) ;
%Straight Lines [id:da8827634409952425] 
\draw    (329.2,70.4) -- (329.2,99.6) ;
%Straight Lines [id:da08240500809926843] 
\draw    (359.2,70.8) -- (359.2,100) ;
%Shape: Rectangle [id:dp8505881042891044] 
\draw  [line width=2.25]  (269.2,99.6) -- (239.6,99.6) -- (239.6,70.4) -- (269.2,70.4) -- cycle ;
%Straight Lines [id:da2295782410185434] 
\draw    (254.8,117.75) -- (254.8,103.55) ;
\draw [shift={(254.8,100.55)}, rotate = 90] [fill={rgb, 255:red, 0; green, 0; blue, 0 }  ][line width=0.08]  [draw opacity=0] (8.93,-4.29) -- (0,0) -- (8.93,4.29) -- cycle    ;
%Straight Lines [id:da2131860495885145] 
\draw [line width=1.5]    (254.8,123.4) -- (254.8,106.15) ;


% Text Node
\draw (107.4,79.2) node [anchor=north west][inner sep=0.75pt]  [font=\small] [align=left] {\begin{minipage}[lt]{18.04pt}\setlength\topsep{0pt}
\begin{center}
$q_\text{R}$
\end{center}

\end{minipage}};
% Text Node
\draw (224.8,85.4) node   [align=left] {\begin{minipage}[lt]{20.13pt}\setlength\topsep{0pt}
\begin{center}
\Vtextvisiblespace
\end{center}

\end{minipage}};
% Text Node
\draw (254.4,85) node   [align=left] {\begin{minipage}[lt]{20.13pt}\setlength\topsep{0pt}
\begin{center}
X
\end{center}

\end{minipage}};
% Text Node
\draw (284.4,85) node   [align=left] {\begin{minipage}[lt]{20.13pt}\setlength\topsep{0pt}
\begin{center}
X
\end{center}

\end{minipage}};
% Text Node
\draw (314,86.4) node   [align=left] {\begin{minipage}[lt]{20.13pt}\setlength\topsep{0pt}
\begin{center}
a
\end{center}

\end{minipage}};
% Text Node
\draw (344.4,85.4) node   [align=left] {\begin{minipage}[lt]{20.13pt}\setlength\topsep{0pt}
\begin{center}
\Vtextvisiblespace
\end{center}

\end{minipage}};
% Text Node
\draw (374,85.4) node   [align=left] {\begin{minipage}[lt]{20.13pt}\setlength\topsep{0pt}
\begin{center}
...
\end{center}

\end{minipage}};
% Text Node
\draw (194.2,86) node   [align=left] {\begin{minipage}[lt]{20.13pt}\setlength\topsep{0pt}
\begin{center}
...
\end{center}

\end{minipage}};


\end{tikzpicture}

\end{adjustbox}

\end{frame}

\begin{frame}{Máquinas de Turing}

\textbf{Ejemplo 3.1}\textbf{.}
$$
    M_{\#\text{a}>\#\text{b}}=(\{q_0, q_\text{a}, q_\text{b}, q_\text{R}, q_\text{F}\}, \{\text{a}, \text{b}\}, \{\text{a}, \text{b}, \text{X}, \Vtextvisiblespace\}, \delta, q_0, \Vtextvisiblespace\:,\{q_\text{F}\})
$$
$$
    \begin{matrix}
        \delta(q_0, \text{a}) = (q_\text{a},\text{X},D) & \delta(q_0, \text{b}) = (q_\text{b},\text{X},D) & \delta(q_0, \text{X}) = (q_0,\text{X},D) &  \\
        \delta(q_\text{a}, \text{a}) = (q_\text{a},\text{a},D) & \delta(q_\text{a}, \text{b}) = (q_\text{R},\text{X},I) & \delta(q_\text{a}, \text{X}) = (q_\text{a},\text{X},D) & \delta(q_\text{a}, \Vtextvisiblespace) = (q_\text{F},\Vtextvisiblespace,S) \\
        \delta(q_\text{b}, \text{a}) = (q_\text{R},\text{X},I) & \delta(q_\text{b}, \text{b}) = (q_\text{b},\text{b},D) & \delta(q_\text{b}, \text{X}) = (q_\text{b},\text{X},D) &  \\
        \delta(q_\text{R}, \text{a}) = (q_\text{R},\text{a},I) & \delta(q_\text{R}, \text{b}) = (q_\text{R},\text{b},I) & \delta(q_\text{R}, \text{X}) = (q_\text{R},\text{X},I) & \delta(q_\text{R}, \Vtextvisiblespace) = (q_0,\Vtextvisiblespace,D) \\
    \end{matrix}
$$

\vspace{5mm}

\begin{adjustbox}{max height={.67\textheight},center}


\tikzset{every picture/.style={line width=0.75pt}} %set default line width to 0.75pt        

\begin{tikzpicture}[x=0.75pt,y=0.75pt,yscale=-1,xscale=1]
%uncomment if require: \path (0,446); %set diagram left start at 0, and has height of 446

%Shape: Triangle [id:dp9549836358247636] 
\draw  [fill={rgb, 255:red, 0; green, 0; blue, 0 }  ,fill opacity=1 ] (134.4,85) -- (141.6,82.2) -- (141.6,87.8) -- cycle ;
%Shape: Circle [id:dp39324764126664746] 
\draw   (101.2,85) .. controls (101.2,75) and (109.31,66.9) .. (119.3,66.9) .. controls (129.3,66.9) and (137.4,75) .. (137.4,85) .. controls (137.4,95) and (129.3,103.1) .. (119.3,103.1) .. controls (109.31,103.1) and (101.2,95) .. (101.2,85) -- cycle ;
%Straight Lines [id:da279972863636345] 
\draw    (379.6,70.2) -- (248.8,70.2) -- (189.6,70.2) ;
%Straight Lines [id:da8063028764122899] 
\draw    (379.2,99.6) -- (189.4,99.6) ;
%Straight Lines [id:da6194955219154272] 
\draw    (210,70.4) -- (210,99.6) ;
%Straight Lines [id:da9291793988668682] 
\draw    (239.6,70.8) -- (239.6,100) ;
%Straight Lines [id:da1753997056206802] 
\draw    (269.2,70.4) -- (269.2,99.6) ;
%Straight Lines [id:da7326138590955908] 
\draw    (299.2,70.4) -- (299.2,99.6) ;
%Straight Lines [id:da8827634409952425] 
\draw    (329.2,70.4) -- (329.2,99.6) ;
%Straight Lines [id:da08240500809926843] 
\draw    (359.2,70.8) -- (359.2,100) ;
%Shape: Rectangle [id:dp8505881042891044] 
\draw  [line width=2.25]  (239.2,99.6) -- (209.6,99.6) -- (209.6,70.4) -- (239.2,70.4) -- cycle ;
%Straight Lines [id:da2295782410185434] 
\draw    (224.8,117.75) -- (224.8,103.55) ;
\draw [shift={(224.8,100.55)}, rotate = 90] [fill={rgb, 255:red, 0; green, 0; blue, 0 }  ][line width=0.08]  [draw opacity=0] (8.93,-4.29) -- (0,0) -- (8.93,4.29) -- cycle    ;
%Straight Lines [id:da2131860495885145] 
\draw [line width=1.5]    (224.8,123.4) -- (224.8,106.15) ;


% Text Node
\draw (107.4,79.2) node [anchor=north west][inner sep=0.75pt]  [font=\small] [align=left] {\begin{minipage}[lt]{18.04pt}\setlength\topsep{0pt}
\begin{center}
$q_\text{R}$
\end{center}

\end{minipage}};
% Text Node
\draw (224.8,85.4) node   [align=left] {\begin{minipage}[lt]{20.13pt}\setlength\topsep{0pt}
\begin{center}
\Vtextvisiblespace
\end{center}

\end{minipage}};
% Text Node
\draw (254.4,85) node   [align=left] {\begin{minipage}[lt]{20.13pt}\setlength\topsep{0pt}
\begin{center}
X
\end{center}

\end{minipage}};
% Text Node
\draw (284.4,85) node   [align=left] {\begin{minipage}[lt]{20.13pt}\setlength\topsep{0pt}
\begin{center}
X
\end{center}

\end{minipage}};
% Text Node
\draw (314,86.4) node   [align=left] {\begin{minipage}[lt]{20.13pt}\setlength\topsep{0pt}
\begin{center}
a
\end{center}

\end{minipage}};
% Text Node
\draw (344.4,85.4) node   [align=left] {\begin{minipage}[lt]{20.13pt}\setlength\topsep{0pt}
\begin{center}
\Vtextvisiblespace
\end{center}

\end{minipage}};
% Text Node
\draw (374,85.4) node   [align=left] {\begin{minipage}[lt]{20.13pt}\setlength\topsep{0pt}
\begin{center}
...
\end{center}

\end{minipage}};
% Text Node
\draw (194.2,86) node   [align=left] {\begin{minipage}[lt]{20.13pt}\setlength\topsep{0pt}
\begin{center}
...
\end{center}

\end{minipage}};


\end{tikzpicture}

\end{adjustbox}

\end{frame}

\begin{frame}{Máquinas de Turing}

\textbf{Ejemplo 3.1}\textbf{.}
$$
    M_{\#\text{a}>\#\text{b}}=(\{q_0, q_\text{a}, q_\text{b}, q_\text{R}, q_\text{F}\}, \{\text{a}, \text{b}\}, \{\text{a}, \text{b}, \text{X}, \Vtextvisiblespace\}, \delta, q_0, \Vtextvisiblespace\:,\{q_\text{F}\})
$$
$$
    \begin{matrix}
        \delta(q_0, \text{a}) = (q_\text{a},\text{X},D) & \delta(q_0, \text{b}) = (q_\text{b},\text{X},D) & \delta(q_0, \text{X}) = (q_0,\text{X},D) &  \\
        \delta(q_\text{a}, \text{a}) = (q_\text{a},\text{a},D) & \delta(q_\text{a}, \text{b}) = (q_\text{R},\text{X},I) & \delta(q_\text{a}, \text{X}) = (q_\text{a},\text{X},D) & \delta(q_\text{a}, \Vtextvisiblespace) = (q_\text{F},\Vtextvisiblespace,S) \\
        \delta(q_\text{b}, \text{a}) = (q_\text{R},\text{X},I) & \delta(q_\text{b}, \text{b}) = (q_\text{b},\text{b},D) & \delta(q_\text{b}, \text{X}) = (q_\text{b},\text{X},D) &  \\
        \delta(q_\text{R}, \text{a}) = (q_\text{R},\text{a},I) & \delta(q_\text{R}, \text{b}) = (q_\text{R},\text{b},I) & \delta(q_\text{R}, \text{X}) = (q_\text{R},\text{X},I) & \underline{\delta(q_\text{R}, \Vtextvisiblespace) = (q_0,\Vtextvisiblespace,D)} \\
    \end{matrix}
$$

\vspace{5mm}

\begin{adjustbox}{max height={.67\textheight},center}


\tikzset{every picture/.style={line width=0.75pt}} %set default line width to 0.75pt        

\begin{tikzpicture}[x=0.75pt,y=0.75pt,yscale=-1,xscale=1]
%uncomment if require: \path (0,446); %set diagram left start at 0, and has height of 446

%Shape: Triangle [id:dp9549836358247636] 
\draw  [fill={rgb, 255:red, 0; green, 0; blue, 0 }  ,fill opacity=1 ] (134.4,85) -- (141.6,82.2) -- (141.6,87.8) -- cycle ;
%Shape: Circle [id:dp39324764126664746] 
\draw   (101.2,85) .. controls (101.2,75) and (109.31,66.9) .. (119.3,66.9) .. controls (129.3,66.9) and (137.4,75) .. (137.4,85) .. controls (137.4,95) and (129.3,103.1) .. (119.3,103.1) .. controls (109.31,103.1) and (101.2,95) .. (101.2,85) -- cycle ;
%Straight Lines [id:da279972863636345] 
\draw    (379.6,70.2) -- (248.8,70.2) -- (189.6,70.2) ;
%Straight Lines [id:da8063028764122899] 
\draw    (379.2,99.6) -- (189.4,99.6) ;
%Straight Lines [id:da6194955219154272] 
\draw    (210,70.4) -- (210,99.6) ;
%Straight Lines [id:da9291793988668682] 
\draw    (239.6,70.8) -- (239.6,100) ;
%Straight Lines [id:da1753997056206802] 
\draw    (269.2,70.4) -- (269.2,99.6) ;
%Straight Lines [id:da7326138590955908] 
\draw    (299.2,70.4) -- (299.2,99.6) ;
%Straight Lines [id:da8827634409952425] 
\draw    (329.2,70.4) -- (329.2,99.6) ;
%Straight Lines [id:da08240500809926843] 
\draw    (359.2,70.8) -- (359.2,100) ;
%Shape: Rectangle [id:dp8505881042891044] 
\draw  [line width=2.25]  (239.2,99.6) -- (209.6,99.6) -- (209.6,70.4) -- (239.2,70.4) -- cycle ;
%Straight Lines [id:da2295782410185434] 
\draw    (224.8,117.75) -- (224.8,103.55) ;
\draw [shift={(224.8,100.55)}, rotate = 90] [fill={rgb, 255:red, 0; green, 0; blue, 0 }  ][line width=0.08]  [draw opacity=0] (8.93,-4.29) -- (0,0) -- (8.93,4.29) -- cycle    ;
%Straight Lines [id:da2131860495885145] 
\draw [line width=1.5]    (224.8,123.4) -- (224.8,106.15) ;


% Text Node
\draw (107.4,79.2) node [anchor=north west][inner sep=0.75pt]  [font=\small] [align=left] {\begin{minipage}[lt]{18.04pt}\setlength\topsep{0pt}
\begin{center}
$q_\text{R}$
\end{center}

\end{minipage}};
% Text Node
\draw (224.8,85.4) node   [align=left] {\begin{minipage}[lt]{20.13pt}\setlength\topsep{0pt}
\begin{center}
\Vtextvisiblespace
\end{center}

\end{minipage}};
% Text Node
\draw (254.4,85) node   [align=left] {\begin{minipage}[lt]{20.13pt}\setlength\topsep{0pt}
\begin{center}
X
\end{center}

\end{minipage}};
% Text Node
\draw (284.4,85) node   [align=left] {\begin{minipage}[lt]{20.13pt}\setlength\topsep{0pt}
\begin{center}
X
\end{center}

\end{minipage}};
% Text Node
\draw (314,86.4) node   [align=left] {\begin{minipage}[lt]{20.13pt}\setlength\topsep{0pt}
\begin{center}
a
\end{center}

\end{minipage}};
% Text Node
\draw (344.4,85.4) node   [align=left] {\begin{minipage}[lt]{20.13pt}\setlength\topsep{0pt}
\begin{center}
\Vtextvisiblespace
\end{center}

\end{minipage}};
% Text Node
\draw (374,85.4) node   [align=left] {\begin{minipage}[lt]{20.13pt}\setlength\topsep{0pt}
\begin{center}
...
\end{center}

\end{minipage}};
% Text Node
\draw (194.2,86) node   [align=left] {\begin{minipage}[lt]{20.13pt}\setlength\topsep{0pt}
\begin{center}
...
\end{center}

\end{minipage}};


\end{tikzpicture}

\end{adjustbox}

\end{frame}


\begin{frame}{Máquinas de Turing}

\textbf{Ejemplo 3.1}\textbf{.}
$$
    M_{\#\text{a}>\#\text{b}}=(\{q_0, q_\text{a}, q_\text{b}, q_\text{R}, q_\text{F}\}, \{\text{a}, \text{b}\}, \{\text{a}, \text{b}, \text{X}, \Vtextvisiblespace\}, \delta, q_0, \Vtextvisiblespace\:,\{q_\text{F}\})
$$
$$
    \begin{matrix}
        \delta(q_0, \text{a}) = (q_\text{a},\text{X},D) & \delta(q_0, \text{b}) = (q_\text{b},\text{X},D) & \delta(q_0, \text{X}) = (q_0,\text{X},D) &  \\
        \delta(q_\text{a}, \text{a}) = (q_\text{a},\text{a},D) & \delta(q_\text{a}, \text{b}) = (q_\text{R},\text{X},I) & \delta(q_\text{a}, \text{X}) = (q_\text{a},\text{X},D) & \delta(q_\text{a}, \Vtextvisiblespace) = (q_\text{F},\Vtextvisiblespace,S) \\
        \delta(q_\text{b}, \text{a}) = (q_\text{R},\text{X},I) & \delta(q_\text{b}, \text{b}) = (q_\text{b},\text{b},D) & \delta(q_\text{b}, \text{X}) = (q_\text{b},\text{X},D) &  \\
        \delta(q_\text{R}, \text{a}) = (q_\text{R},\text{a},I) & \delta(q_\text{R}, \text{b}) = (q_\text{R},\text{b},I) & \delta(q_\text{R}, \text{X}) = (q_\text{R},\text{X},I) & \delta(q_\text{R}, \Vtextvisiblespace) = (q_0,\Vtextvisiblespace,D) \\
    \end{matrix}
$$

\vspace{5mm}

\begin{adjustbox}{max height={.67\textheight},center}


\tikzset{every picture/.style={line width=0.75pt}} %set default line width to 0.75pt        

\begin{tikzpicture}[x=0.75pt,y=0.75pt,yscale=-1,xscale=1]
%uncomment if require: \path (0,446); %set diagram left start at 0, and has height of 446

%Shape: Triangle [id:dp9549836358247636] 
\draw  [fill={rgb, 255:red, 0; green, 0; blue, 0 }  ,fill opacity=1 ] (134.4,85) -- (141.6,82.2) -- (141.6,87.8) -- cycle ;
%Shape: Circle [id:dp39324764126664746] 
\draw   (101.2,85) .. controls (101.2,75) and (109.31,66.9) .. (119.3,66.9) .. controls (129.3,66.9) and (137.4,75) .. (137.4,85) .. controls (137.4,95) and (129.3,103.1) .. (119.3,103.1) .. controls (109.31,103.1) and (101.2,95) .. (101.2,85) -- cycle ;
%Straight Lines [id:da279972863636345] 
\draw    (379.6,70.2) -- (248.8,70.2) -- (189.6,70.2) ;
%Straight Lines [id:da8063028764122899] 
\draw    (379.2,99.6) -- (189.4,99.6) ;
%Straight Lines [id:da6194955219154272] 
\draw    (210,70.4) -- (210,99.6) ;
%Straight Lines [id:da9291793988668682] 
\draw    (239.6,70.8) -- (239.6,100) ;
%Straight Lines [id:da1753997056206802] 
\draw    (269.2,70.4) -- (269.2,99.6) ;
%Straight Lines [id:da7326138590955908] 
\draw    (299.2,70.4) -- (299.2,99.6) ;
%Straight Lines [id:da8827634409952425] 
\draw    (329.2,70.4) -- (329.2,99.6) ;
%Straight Lines [id:da08240500809926843] 
\draw    (359.2,70.8) -- (359.2,100) ;
%Shape: Rectangle [id:dp8505881042891044] 
\draw  [line width=2.25]  (269.2,99.6) -- (239.6,99.6) -- (239.6,70.4) -- (269.2,70.4) -- cycle ;
%Straight Lines [id:da2295782410185434] 
\draw    (254.8,117.75) -- (254.8,103.55) ;
\draw [shift={(254.8,100.55)}, rotate = 90] [fill={rgb, 255:red, 0; green, 0; blue, 0 }  ][line width=0.08]  [draw opacity=0] (8.93,-4.29) -- (0,0) -- (8.93,4.29) -- cycle    ;
%Straight Lines [id:da2131860495885145] 
\draw [line width=1.5]    (254.8,123.4) -- (254.8,106.15) ;


% Text Node
\draw (107.4,79.2) node [anchor=north west][inner sep=0.75pt]  [font=\small] [align=left] {\begin{minipage}[lt]{18.04pt}\setlength\topsep{0pt}
\begin{center}
$q_0$
\end{center}

\end{minipage}};
% Text Node
\draw (224.8,85.4) node   [align=left] {\begin{minipage}[lt]{20.13pt}\setlength\topsep{0pt}
\begin{center}
\Vtextvisiblespace
\end{center}

\end{minipage}};
% Text Node
\draw (254.4,85) node   [align=left] {\begin{minipage}[lt]{20.13pt}\setlength\topsep{0pt}
\begin{center}
X
\end{center}

\end{minipage}};
% Text Node
\draw (284.4,85) node   [align=left] {\begin{minipage}[lt]{20.13pt}\setlength\topsep{0pt}
\begin{center}
X
\end{center}

\end{minipage}};
% Text Node
\draw (314,86.4) node   [align=left] {\begin{minipage}[lt]{20.13pt}\setlength\topsep{0pt}
\begin{center}
a
\end{center}

\end{minipage}};
% Text Node
\draw (344.4,85.4) node   [align=left] {\begin{minipage}[lt]{20.13pt}\setlength\topsep{0pt}
\begin{center}
\Vtextvisiblespace
\end{center}

\end{minipage}};
% Text Node
\draw (374,85.4) node   [align=left] {\begin{minipage}[lt]{20.13pt}\setlength\topsep{0pt}
\begin{center}
...
\end{center}

\end{minipage}};
% Text Node
\draw (194.2,86) node   [align=left] {\begin{minipage}[lt]{20.13pt}\setlength\topsep{0pt}
\begin{center}
...
\end{center}

\end{minipage}};


\end{tikzpicture}

\end{adjustbox}

\end{frame}


\begin{frame}{Máquinas de Turing}

\textbf{Ejemplo 3.1}\textbf{.}
$$
    M_{\#\text{a}>\#\text{b}}=(\{q_0, q_\text{a}, q_\text{b}, q_\text{R}, q_\text{F}\}, \{\text{a}, \text{b}\}, \{\text{a}, \text{b}, \text{X}, \Vtextvisiblespace\}, \delta, q_0, \Vtextvisiblespace\:,\{q_\text{F}\})
$$
$$
    \begin{matrix}
        \delta(q_0, \text{a}) = (q_\text{a},\text{X},D) & \delta(q_0, \text{b}) = (q_\text{b},\text{X},D) & \underline{\delta(q_0, \text{X}) = (q_0,\text{X},D)} &  \\
        \delta(q_\text{a}, \text{a}) = (q_\text{a},\text{a},D) & \delta(q_\text{a}, \text{b}) = (q_\text{R},\text{X},I) & \delta(q_\text{a}, \text{X}) = (q_\text{a},\text{X},D) & \delta(q_\text{a}, \Vtextvisiblespace) = (q_\text{F},\Vtextvisiblespace,S) \\
        \delta(q_\text{b}, \text{a}) = (q_\text{R},\text{X},I) & \delta(q_\text{b}, \text{b}) = (q_\text{b},\text{b},D) & \delta(q_\text{b}, \text{X}) = (q_\text{b},\text{X},D) &  \\
        \delta(q_\text{R}, \text{a}) = (q_\text{R},\text{a},I) & \delta(q_\text{R}, \text{b}) = (q_\text{R},\text{b},I) & \delta(q_\text{R}, \text{X}) = (q_\text{R},\text{X},I) & \delta(q_\text{R}, \Vtextvisiblespace) = (q_0,\Vtextvisiblespace,D) \\
    \end{matrix}
$$

\vspace{5mm}

\begin{adjustbox}{max height={.67\textheight},center}


\tikzset{every picture/.style={line width=0.75pt}} %set default line width to 0.75pt        

\begin{tikzpicture}[x=0.75pt,y=0.75pt,yscale=-1,xscale=1]
%uncomment if require: \path (0,446); %set diagram left start at 0, and has height of 446

%Shape: Triangle [id:dp9549836358247636] 
\draw  [fill={rgb, 255:red, 0; green, 0; blue, 0 }  ,fill opacity=1 ] (134.4,85) -- (141.6,82.2) -- (141.6,87.8) -- cycle ;
%Shape: Circle [id:dp39324764126664746] 
\draw   (101.2,85) .. controls (101.2,75) and (109.31,66.9) .. (119.3,66.9) .. controls (129.3,66.9) and (137.4,75) .. (137.4,85) .. controls (137.4,95) and (129.3,103.1) .. (119.3,103.1) .. controls (109.31,103.1) and (101.2,95) .. (101.2,85) -- cycle ;
%Straight Lines [id:da279972863636345] 
\draw    (379.6,70.2) -- (248.8,70.2) -- (189.6,70.2) ;
%Straight Lines [id:da8063028764122899] 
\draw    (379.2,99.6) -- (189.4,99.6) ;
%Straight Lines [id:da6194955219154272] 
\draw    (210,70.4) -- (210,99.6) ;
%Straight Lines [id:da9291793988668682] 
\draw    (239.6,70.8) -- (239.6,100) ;
%Straight Lines [id:da1753997056206802] 
\draw    (269.2,70.4) -- (269.2,99.6) ;
%Straight Lines [id:da7326138590955908] 
\draw    (299.2,70.4) -- (299.2,99.6) ;
%Straight Lines [id:da8827634409952425] 
\draw    (329.2,70.4) -- (329.2,99.6) ;
%Straight Lines [id:da08240500809926843] 
\draw    (359.2,70.8) -- (359.2,100) ;
%Shape: Rectangle [id:dp8505881042891044] 
\draw  [line width=2.25]  (269.2,99.6) -- (239.6,99.6) -- (239.6,70.4) -- (269.2,70.4) -- cycle ;
%Straight Lines [id:da2295782410185434] 
\draw    (254.8,117.75) -- (254.8,103.55) ;
\draw [shift={(254.8,100.55)}, rotate = 90] [fill={rgb, 255:red, 0; green, 0; blue, 0 }  ][line width=0.08]  [draw opacity=0] (8.93,-4.29) -- (0,0) -- (8.93,4.29) -- cycle    ;
%Straight Lines [id:da2131860495885145] 
\draw [line width=1.5]    (254.8,123.4) -- (254.8,106.15) ;


% Text Node
\draw (107.4,79.2) node [anchor=north west][inner sep=0.75pt]  [font=\small] [align=left] {\begin{minipage}[lt]{18.04pt}\setlength\topsep{0pt}
\begin{center}
$q_0$
\end{center}

\end{minipage}};
% Text Node
\draw (224.8,85.4) node   [align=left] {\begin{minipage}[lt]{20.13pt}\setlength\topsep{0pt}
\begin{center}
\Vtextvisiblespace
\end{center}

\end{minipage}};
% Text Node
\draw (254.4,85) node   [align=left] {\begin{minipage}[lt]{20.13pt}\setlength\topsep{0pt}
\begin{center}
X
\end{center}

\end{minipage}};
% Text Node
\draw (284.4,85) node   [align=left] {\begin{minipage}[lt]{20.13pt}\setlength\topsep{0pt}
\begin{center}
X
\end{center}

\end{minipage}};
% Text Node
\draw (314,86.4) node   [align=left] {\begin{minipage}[lt]{20.13pt}\setlength\topsep{0pt}
\begin{center}
a
\end{center}

\end{minipage}};
% Text Node
\draw (344.4,85.4) node   [align=left] {\begin{minipage}[lt]{20.13pt}\setlength\topsep{0pt}
\begin{center}
\Vtextvisiblespace
\end{center}

\end{minipage}};
% Text Node
\draw (374,85.4) node   [align=left] {\begin{minipage}[lt]{20.13pt}\setlength\topsep{0pt}
\begin{center}
...
\end{center}

\end{minipage}};
% Text Node
\draw (194.2,86) node   [align=left] {\begin{minipage}[lt]{20.13pt}\setlength\topsep{0pt}
\begin{center}
...
\end{center}

\end{minipage}};


\end{tikzpicture}

\end{adjustbox}

\end{frame}


\begin{frame}{Máquinas de Turing}

\textbf{Ejemplo 3.1}\textbf{.}
$$
    M_{\#\text{a}>\#\text{b}}=(\{q_0, q_\text{a}, q_\text{b}, q_\text{R}, q_\text{F}\}, \{\text{a}, \text{b}\}, \{\text{a}, \text{b}, \text{X}, \Vtextvisiblespace\}, \delta, q_0, \Vtextvisiblespace\:,\{q_\text{F}\})
$$
$$
    \begin{matrix}
        \delta(q_0, \text{a}) = (q_\text{a},\text{X},D) & \delta(q_0, \text{b}) = (q_\text{b},\text{X},D) & \delta(q_0, \text{X}) = (q_0,\text{X},D) &  \\
        \delta(q_\text{a}, \text{a}) = (q_\text{a},\text{a},D) & \delta(q_\text{a}, \text{b}) = (q_\text{R},\text{X},I) & \delta(q_\text{a}, \text{X}) = (q_\text{a},\text{X},D) & \delta(q_\text{a}, \Vtextvisiblespace) = (q_\text{F},\Vtextvisiblespace,S) \\
        \delta(q_\text{b}, \text{a}) = (q_\text{R},\text{X},I) & \delta(q_\text{b}, \text{b}) = (q_\text{b},\text{b},D) & \delta(q_\text{b}, \text{X}) = (q_\text{b},\text{X},D) &  \\
        \delta(q_\text{R}, \text{a}) = (q_\text{R},\text{a},I) & \delta(q_\text{R}, \text{b}) = (q_\text{R},\text{b},I) & \delta(q_\text{R}, \text{X}) = (q_\text{R},\text{X},I) & \delta(q_\text{R}, \Vtextvisiblespace) = (q_0,\Vtextvisiblespace,D) \\
    \end{matrix}
$$

\vspace{5mm}

\begin{adjustbox}{max height={.67\textheight},center}


\tikzset{every picture/.style={line width=0.75pt}} %set default line width to 0.75pt        

\begin{tikzpicture}[x=0.75pt,y=0.75pt,yscale=-1,xscale=1]
%uncomment if require: \path (0,446); %set diagram left start at 0, and has height of 446

%Shape: Triangle [id:dp9549836358247636] 
\draw  [fill={rgb, 255:red, 0; green, 0; blue, 0 }  ,fill opacity=1 ] (134.4,85) -- (141.6,82.2) -- (141.6,87.8) -- cycle ;
%Shape: Circle [id:dp39324764126664746] 
\draw   (101.2,85) .. controls (101.2,75) and (109.31,66.9) .. (119.3,66.9) .. controls (129.3,66.9) and (137.4,75) .. (137.4,85) .. controls (137.4,95) and (129.3,103.1) .. (119.3,103.1) .. controls (109.31,103.1) and (101.2,95) .. (101.2,85) -- cycle ;
%Straight Lines [id:da279972863636345] 
\draw    (379.6,70.2) -- (248.8,70.2) -- (189.6,70.2) ;
%Straight Lines [id:da8063028764122899] 
\draw    (379.2,99.6) -- (189.4,99.6) ;
%Straight Lines [id:da6194955219154272] 
\draw    (210,70.4) -- (210,99.6) ;
%Straight Lines [id:da9291793988668682] 
\draw    (239.6,70.8) -- (239.6,100) ;
%Straight Lines [id:da1753997056206802] 
\draw    (269.2,70.4) -- (269.2,99.6) ;
%Straight Lines [id:da7326138590955908] 
\draw    (299.2,70.4) -- (299.2,99.6) ;
%Straight Lines [id:da8827634409952425] 
\draw    (329.2,70.4) -- (329.2,99.6) ;
%Straight Lines [id:da08240500809926843] 
\draw    (359.2,70.8) -- (359.2,100) ;
%Shape: Rectangle [id:dp8505881042891044] 
\draw  [line width=2.25]  (299.2,99.6) -- (269.6,99.6) -- (269.6,70.4) -- (299.2,70.4) -- cycle ;
%Straight Lines [id:da2295782410185434] 
\draw    (284.8,117.75) -- (284.8,103.55) ;
\draw [shift={(284.8,100.55)}, rotate = 90] [fill={rgb, 255:red, 0; green, 0; blue, 0 }  ][line width=0.08]  [draw opacity=0] (8.93,-4.29) -- (0,0) -- (8.93,4.29) -- cycle    ;
%Straight Lines [id:da2131860495885145] 
\draw [line width=1.5]    (284.8,123.4) -- (284.8,106.15) ;


% Text Node
\draw (107.4,79.2) node [anchor=north west][inner sep=0.75pt]  [font=\small] [align=left] {\begin{minipage}[lt]{18.04pt}\setlength\topsep{0pt}
\begin{center}
$q_0$
\end{center}

\end{minipage}};
% Text Node
\draw (224.8,85.4) node   [align=left] {\begin{minipage}[lt]{20.13pt}\setlength\topsep{0pt}
\begin{center}
\Vtextvisiblespace
\end{center}

\end{minipage}};
% Text Node
\draw (254.4,85) node   [align=left] {\begin{minipage}[lt]{20.13pt}\setlength\topsep{0pt}
\begin{center}
X
\end{center}

\end{minipage}};
% Text Node
\draw (284.4,85) node   [align=left] {\begin{minipage}[lt]{20.13pt}\setlength\topsep{0pt}
\begin{center}
X
\end{center}

\end{minipage}};
% Text Node
\draw (314,86.4) node   [align=left] {\begin{minipage}[lt]{20.13pt}\setlength\topsep{0pt}
\begin{center}
a
\end{center}

\end{minipage}};
% Text Node
\draw (344.4,85.4) node   [align=left] {\begin{minipage}[lt]{20.13pt}\setlength\topsep{0pt}
\begin{center}
\Vtextvisiblespace
\end{center}

\end{minipage}};
% Text Node
\draw (374,85.4) node   [align=left] {\begin{minipage}[lt]{20.13pt}\setlength\topsep{0pt}
\begin{center}
...
\end{center}

\end{minipage}};
% Text Node
\draw (194.2,86) node   [align=left] {\begin{minipage}[lt]{20.13pt}\setlength\topsep{0pt}
\begin{center}
...
\end{center}

\end{minipage}};


\end{tikzpicture}

\end{adjustbox}

\end{frame}


\begin{frame}{Máquinas de Turing}

\textbf{Ejemplo 3.1}\textbf{.}
$$
    M_{\#\text{a}>\#\text{b}}=(\{q_0, q_\text{a}, q_\text{b}, q_\text{R}, q_\text{F}\}, \{\text{a}, \text{b}\}, \{\text{a}, \text{b}, \text{X}, \Vtextvisiblespace\}, \delta, q_0, \Vtextvisiblespace\:,\{q_\text{F}\})
$$
$$
    \begin{matrix}
        \delta(q_0, \text{a}) = (q_\text{a},\text{X},D) & \delta(q_0, \text{b}) = (q_\text{b},\text{X},D) & \underline{\delta(q_0, \text{X}) = (q_0,\text{X},D)} &  \\
        \delta(q_\text{a}, \text{a}) = (q_\text{a},\text{a},D) & \delta(q_\text{a}, \text{b}) = (q_\text{R},\text{X},I) & \delta(q_\text{a}, \text{X}) = (q_\text{a},\text{X},D) & \delta(q_\text{a}, \Vtextvisiblespace) = (q_\text{F},\Vtextvisiblespace,S) \\
        \delta(q_\text{b}, \text{a}) = (q_\text{R},\text{X},I) & \delta(q_\text{b}, \text{b}) = (q_\text{b},\text{b},D) & \delta(q_\text{b}, \text{X}) = (q_\text{b},\text{X},D) &  \\
        \delta(q_\text{R}, \text{a}) = (q_\text{R},\text{a},I) & \delta(q_\text{R}, \text{b}) = (q_\text{R},\text{b},I) & \delta(q_\text{R}, \text{X}) = (q_\text{R},\text{X},I) & \delta(q_\text{R}, \Vtextvisiblespace) = (q_0,\Vtextvisiblespace,D) \\
    \end{matrix}
$$

\vspace{5mm}

\begin{adjustbox}{max height={.67\textheight},center}


\tikzset{every picture/.style={line width=0.75pt}} %set default line width to 0.75pt        

\begin{tikzpicture}[x=0.75pt,y=0.75pt,yscale=-1,xscale=1]
%uncomment if require: \path (0,446); %set diagram left start at 0, and has height of 446

%Shape: Triangle [id:dp9549836358247636] 
\draw  [fill={rgb, 255:red, 0; green, 0; blue, 0 }  ,fill opacity=1 ] (134.4,85) -- (141.6,82.2) -- (141.6,87.8) -- cycle ;
%Shape: Circle [id:dp39324764126664746] 
\draw   (101.2,85) .. controls (101.2,75) and (109.31,66.9) .. (119.3,66.9) .. controls (129.3,66.9) and (137.4,75) .. (137.4,85) .. controls (137.4,95) and (129.3,103.1) .. (119.3,103.1) .. controls (109.31,103.1) and (101.2,95) .. (101.2,85) -- cycle ;
%Straight Lines [id:da279972863636345] 
\draw    (379.6,70.2) -- (248.8,70.2) -- (189.6,70.2) ;
%Straight Lines [id:da8063028764122899] 
\draw    (379.2,99.6) -- (189.4,99.6) ;
%Straight Lines [id:da6194955219154272] 
\draw    (210,70.4) -- (210,99.6) ;
%Straight Lines [id:da9291793988668682] 
\draw    (239.6,70.8) -- (239.6,100) ;
%Straight Lines [id:da1753997056206802] 
\draw    (269.2,70.4) -- (269.2,99.6) ;
%Straight Lines [id:da7326138590955908] 
\draw    (299.2,70.4) -- (299.2,99.6) ;
%Straight Lines [id:da8827634409952425] 
\draw    (329.2,70.4) -- (329.2,99.6) ;
%Straight Lines [id:da08240500809926843] 
\draw    (359.2,70.8) -- (359.2,100) ;
%Shape: Rectangle [id:dp8505881042891044] 
\draw  [line width=2.25]  (299.2,99.6) -- (269.6,99.6) -- (269.6,70.4) -- (299.2,70.4) -- cycle ;
%Straight Lines [id:da2295782410185434] 
\draw    (284.8,117.75) -- (284.8,103.55) ;
\draw [shift={(284.8,100.55)}, rotate = 90] [fill={rgb, 255:red, 0; green, 0; blue, 0 }  ][line width=0.08]  [draw opacity=0] (8.93,-4.29) -- (0,0) -- (8.93,4.29) -- cycle    ;
%Straight Lines [id:da2131860495885145] 
\draw [line width=1.5]    (284.8,123.4) -- (284.8,106.15) ;


% Text Node
\draw (107.4,79.2) node [anchor=north west][inner sep=0.75pt]  [font=\small] [align=left] {\begin{minipage}[lt]{18.04pt}\setlength\topsep{0pt}
\begin{center}
$q_0$
\end{center}

\end{minipage}};
% Text Node
\draw (224.8,85.4) node   [align=left] {\begin{minipage}[lt]{20.13pt}\setlength\topsep{0pt}
\begin{center}
\Vtextvisiblespace
\end{center}

\end{minipage}};
% Text Node
\draw (254.4,85) node   [align=left] {\begin{minipage}[lt]{20.13pt}\setlength\topsep{0pt}
\begin{center}
X
\end{center}

\end{minipage}};
% Text Node
\draw (284.4,85) node   [align=left] {\begin{minipage}[lt]{20.13pt}\setlength\topsep{0pt}
\begin{center}
X
\end{center}

\end{minipage}};
% Text Node
\draw (314,86.4) node   [align=left] {\begin{minipage}[lt]{20.13pt}\setlength\topsep{0pt}
\begin{center}
a
\end{center}

\end{minipage}};
% Text Node
\draw (344.4,85.4) node   [align=left] {\begin{minipage}[lt]{20.13pt}\setlength\topsep{0pt}
\begin{center}
\Vtextvisiblespace
\end{center}

\end{minipage}};
% Text Node
\draw (374,85.4) node   [align=left] {\begin{minipage}[lt]{20.13pt}\setlength\topsep{0pt}
\begin{center}
...
\end{center}

\end{minipage}};
% Text Node
\draw (194.2,86) node   [align=left] {\begin{minipage}[lt]{20.13pt}\setlength\topsep{0pt}
\begin{center}
...
\end{center}

\end{minipage}};


\end{tikzpicture}

\end{adjustbox}

\end{frame}

\begin{frame}{Máquinas de Turing}

\textbf{Ejemplo 3.1}\textbf{.}
$$
    M_{\#\text{a}>\#\text{b}}=(\{q_0, q_\text{a}, q_\text{b}, q_\text{R}, q_\text{F}\}, \{\text{a}, \text{b}\}, \{\text{a}, \text{b}, \text{X}, \Vtextvisiblespace\}, \delta, q_0, \Vtextvisiblespace\:,\{q_\text{F}\})
$$
$$
    \begin{matrix}
        \delta(q_0, \text{a}) = (q_\text{a},\text{X},D) & \delta(q_0, \text{b}) = (q_\text{b},\text{X},D) & \delta(q_0, \text{X}) = (q_0,\text{X},D) &  \\
        \delta(q_\text{a}, \text{a}) = (q_\text{a},\text{a},D) & \delta(q_\text{a}, \text{b}) = (q_\text{R},\text{X},I) & \delta(q_\text{a}, \text{X}) = (q_\text{a},\text{X},D) & \delta(q_\text{a}, \Vtextvisiblespace) = (q_\text{F},\Vtextvisiblespace,S) \\
        \delta(q_\text{b}, \text{a}) = (q_\text{R},\text{X},I) & \delta(q_\text{b}, \text{b}) = (q_\text{b},\text{b},D) & \delta(q_\text{b}, \text{X}) = (q_\text{b},\text{X},D) &  \\
        \delta(q_\text{R}, \text{a}) = (q_\text{R},\text{a},I) & \delta(q_\text{R}, \text{b}) = (q_\text{R},\text{b},I) & \delta(q_\text{R}, \text{X}) = (q_\text{R},\text{X},I) & \delta(q_\text{R}, \Vtextvisiblespace) = (q_0,\Vtextvisiblespace,D) \\
    \end{matrix}
$$

\vspace{5mm}

\begin{adjustbox}{max height={.67\textheight},center}


\tikzset{every picture/.style={line width=0.75pt}} %set default line width to 0.75pt        

\begin{tikzpicture}[x=0.75pt,y=0.75pt,yscale=-1,xscale=1]
%uncomment if require: \path (0,446); %set diagram left start at 0, and has height of 446

%Shape: Triangle [id:dp9549836358247636] 
\draw  [fill={rgb, 255:red, 0; green, 0; blue, 0 }  ,fill opacity=1 ] (134.4,85) -- (141.6,82.2) -- (141.6,87.8) -- cycle ;
%Shape: Circle [id:dp39324764126664746] 
\draw   (101.2,85) .. controls (101.2,75) and (109.31,66.9) .. (119.3,66.9) .. controls (129.3,66.9) and (137.4,75) .. (137.4,85) .. controls (137.4,95) and (129.3,103.1) .. (119.3,103.1) .. controls (109.31,103.1) and (101.2,95) .. (101.2,85) -- cycle ;
%Straight Lines [id:da279972863636345] 
\draw    (379.6,70.2) -- (248.8,70.2) -- (189.6,70.2) ;
%Straight Lines [id:da8063028764122899] 
\draw    (379.2,99.6) -- (189.4,99.6) ;
%Straight Lines [id:da6194955219154272] 
\draw    (210,70.4) -- (210,99.6) ;
%Straight Lines [id:da9291793988668682] 
\draw    (239.6,70.8) -- (239.6,100) ;
%Straight Lines [id:da1753997056206802] 
\draw    (269.2,70.4) -- (269.2,99.6) ;
%Straight Lines [id:da7326138590955908] 
\draw    (299.2,70.4) -- (299.2,99.6) ;
%Straight Lines [id:da8827634409952425] 
\draw    (329.2,70.4) -- (329.2,99.6) ;
%Straight Lines [id:da08240500809926843] 
\draw    (359.2,70.8) -- (359.2,100) ;
%Shape: Rectangle [id:dp8505881042891044] 
\draw  [line width=2.25]  (329.2,99.6) -- (299.6,99.6) -- (299.6,70.4) -- (329.2,70.4) -- cycle ;
%Straight Lines [id:da2295782410185434] 
\draw    (314.8,117.75) -- (314.8,103.55) ;
\draw [shift={(314.8,100.55)}, rotate = 90] [fill={rgb, 255:red, 0; green, 0; blue, 0 }  ][line width=0.08]  [draw opacity=0] (8.93,-4.29) -- (0,0) -- (8.93,4.29) -- cycle    ;
%Straight Lines [id:da2131860495885145] 
\draw [line width=1.5]    (314.8,123.4) -- (314.8,106.15) ;


% Text Node
\draw (107.4,79.2) node [anchor=north west][inner sep=0.75pt]  [font=\small] [align=left] {\begin{minipage}[lt]{18.04pt}\setlength\topsep{0pt}
\begin{center}
$q_0$
\end{center}

\end{minipage}};
% Text Node
\draw (224.8,85.4) node   [align=left] {\begin{minipage}[lt]{20.13pt}\setlength\topsep{0pt}
\begin{center}
\Vtextvisiblespace
\end{center}

\end{minipage}};
% Text Node
\draw (254.4,85) node   [align=left] {\begin{minipage}[lt]{20.13pt}\setlength\topsep{0pt}
\begin{center}
X
\end{center}

\end{minipage}};
% Text Node
\draw (284.4,85) node   [align=left] {\begin{minipage}[lt]{20.13pt}\setlength\topsep{0pt}
\begin{center}
X
\end{center}

\end{minipage}};
% Text Node
\draw (314,86.4) node   [align=left] {\begin{minipage}[lt]{20.13pt}\setlength\topsep{0pt}
\begin{center}
a
\end{center}

\end{minipage}};
% Text Node
\draw (344.4,85.4) node   [align=left] {\begin{minipage}[lt]{20.13pt}\setlength\topsep{0pt}
\begin{center}
\Vtextvisiblespace
\end{center}

\end{minipage}};
% Text Node
\draw (374,85.4) node   [align=left] {\begin{minipage}[lt]{20.13pt}\setlength\topsep{0pt}
\begin{center}
...
\end{center}

\end{minipage}};
% Text Node
\draw (194.2,86) node   [align=left] {\begin{minipage}[lt]{20.13pt}\setlength\topsep{0pt}
\begin{center}
...
\end{center}

\end{minipage}};


\end{tikzpicture}

\end{adjustbox}

\end{frame}

\begin{frame}{Máquinas de Turing}

\textbf{Ejemplo 3.1}\textbf{.}
$$
    M_{\#\text{a}>\#\text{b}}=(\{q_0, q_\text{a}, q_\text{b}, q_\text{R}, q_\text{F}\}, \{\text{a}, \text{b}\}, \{\text{a}, \text{b}, \text{X}, \Vtextvisiblespace\}, \delta, q_0, \Vtextvisiblespace\:,\{q_\text{F}\})
$$
$$
    \begin{matrix}
        \underline{\delta(q_0, \text{a}) = (q_\text{a},\text{X},D)} & \delta(q_0, \text{b}) = (q_\text{b},\text{X},D) & \delta(q_0, \text{X}) = (q_0,\text{X},D) &  \\
        \delta(q_\text{a}, \text{a}) = (q_\text{a},\text{a},D) & \delta(q_\text{a}, \text{b}) = (q_\text{R},\text{X},I) & \delta(q_\text{a}, \text{X}) = (q_\text{a},\text{X},D) & \delta(q_\text{a}, \Vtextvisiblespace) = (q_\text{F},\Vtextvisiblespace,S) \\
        \delta(q_\text{b}, \text{a}) = (q_\text{R},\text{X},I) & \delta(q_\text{b}, \text{b}) = (q_\text{b},\text{b},D) & \delta(q_\text{b}, \text{X}) = (q_\text{b},\text{X},D) &  \\
        \delta(q_\text{R}, \text{a}) = (q_\text{R},\text{a},I) & \delta(q_\text{R}, \text{b}) = (q_\text{R},\text{b},I) & \delta(q_\text{R}, \text{X}) = (q_\text{R},\text{X},I) & \delta(q_\text{R}, \Vtextvisiblespace) = (q_0,\Vtextvisiblespace,D) \\
    \end{matrix}
$$

\vspace{5mm}

\begin{adjustbox}{max height={.67\textheight},center}


\tikzset{every picture/.style={line width=0.75pt}} %set default line width to 0.75pt        

\begin{tikzpicture}[x=0.75pt,y=0.75pt,yscale=-1,xscale=1]
%uncomment if require: \path (0,446); %set diagram left start at 0, and has height of 446

%Shape: Triangle [id:dp9549836358247636] 
\draw  [fill={rgb, 255:red, 0; green, 0; blue, 0 }  ,fill opacity=1 ] (134.4,85) -- (141.6,82.2) -- (141.6,87.8) -- cycle ;
%Shape: Circle [id:dp39324764126664746] 
\draw   (101.2,85) .. controls (101.2,75) and (109.31,66.9) .. (119.3,66.9) .. controls (129.3,66.9) and (137.4,75) .. (137.4,85) .. controls (137.4,95) and (129.3,103.1) .. (119.3,103.1) .. controls (109.31,103.1) and (101.2,95) .. (101.2,85) -- cycle ;
%Straight Lines [id:da279972863636345] 
\draw    (379.6,70.2) -- (248.8,70.2) -- (189.6,70.2) ;
%Straight Lines [id:da8063028764122899] 
\draw    (379.2,99.6) -- (189.4,99.6) ;
%Straight Lines [id:da6194955219154272] 
\draw    (210,70.4) -- (210,99.6) ;
%Straight Lines [id:da9291793988668682] 
\draw    (239.6,70.8) -- (239.6,100) ;
%Straight Lines [id:da1753997056206802] 
\draw    (269.2,70.4) -- (269.2,99.6) ;
%Straight Lines [id:da7326138590955908] 
\draw    (299.2,70.4) -- (299.2,99.6) ;
%Straight Lines [id:da8827634409952425] 
\draw    (329.2,70.4) -- (329.2,99.6) ;
%Straight Lines [id:da08240500809926843] 
\draw    (359.2,70.8) -- (359.2,100) ;
%Shape: Rectangle [id:dp8505881042891044] 
\draw  [line width=2.25]  (329.2,99.6) -- (299.6,99.6) -- (299.6,70.4) -- (329.2,70.4) -- cycle ;
%Straight Lines [id:da2295782410185434] 
\draw    (314.8,117.75) -- (314.8,103.55) ;
\draw [shift={(314.8,100.55)}, rotate = 90] [fill={rgb, 255:red, 0; green, 0; blue, 0 }  ][line width=0.08]  [draw opacity=0] (8.93,-4.29) -- (0,0) -- (8.93,4.29) -- cycle    ;
%Straight Lines [id:da2131860495885145] 
\draw [line width=1.5]    (314.8,123.4) -- (314.8,106.15) ;


% Text Node
\draw (107.4,79.2) node [anchor=north west][inner sep=0.75pt]  [font=\small] [align=left] {\begin{minipage}[lt]{18.04pt}\setlength\topsep{0pt}
\begin{center}
$q_0$
\end{center}

\end{minipage}};
% Text Node
\draw (224.8,85.4) node   [align=left] {\begin{minipage}[lt]{20.13pt}\setlength\topsep{0pt}
\begin{center}
\Vtextvisiblespace
\end{center}

\end{minipage}};
% Text Node
\draw (254.4,85) node   [align=left] {\begin{minipage}[lt]{20.13pt}\setlength\topsep{0pt}
\begin{center}
X
\end{center}

\end{minipage}};
% Text Node
\draw (284.4,85) node   [align=left] {\begin{minipage}[lt]{20.13pt}\setlength\topsep{0pt}
\begin{center}
X
\end{center}

\end{minipage}};
% Text Node
\draw (314,86.4) node   [align=left] {\begin{minipage}[lt]{20.13pt}\setlength\topsep{0pt}
\begin{center}
a
\end{center}

\end{minipage}};
% Text Node
\draw (344.4,85.4) node   [align=left] {\begin{minipage}[lt]{20.13pt}\setlength\topsep{0pt}
\begin{center}
\Vtextvisiblespace
\end{center}

\end{minipage}};
% Text Node
\draw (374,85.4) node   [align=left] {\begin{minipage}[lt]{20.13pt}\setlength\topsep{0pt}
\begin{center}
...
\end{center}

\end{minipage}};
% Text Node
\draw (194.2,86) node   [align=left] {\begin{minipage}[lt]{20.13pt}\setlength\topsep{0pt}
\begin{center}
...
\end{center}

\end{minipage}};


\end{tikzpicture}

\end{adjustbox}

\end{frame}


\begin{frame}{Máquinas de Turing}

\textbf{Ejemplo 3.1}\textbf{.}
$$
    M_{\#\text{a}>\#\text{b}}=(\{q_0, q_\text{a}, q_\text{b}, q_\text{R}, q_\text{F}\}, \{\text{a}, \text{b}\}, \{\text{a}, \text{b}, \text{X}, \Vtextvisiblespace\}, \delta, q_0, \Vtextvisiblespace\:,\{q_\text{F}\})
$$
$$
    \begin{matrix}
        \delta(q_0, \text{a}) = (q_\text{a},\text{X},D) & \delta(q_0, \text{b}) = (q_\text{b},\text{X},D) & \delta(q_0, \text{X}) = (q_0,\text{X},D) &  \\
        \delta(q_\text{a}, \text{a}) = (q_\text{a},\text{a},D) & \delta(q_\text{a}, \text{b}) = (q_\text{R},\text{X},I) & \delta(q_\text{a}, \text{X}) = (q_\text{a},\text{X},D) & \delta(q_\text{a}, \Vtextvisiblespace) = (q_\text{F},\Vtextvisiblespace,S) \\
        \delta(q_\text{b}, \text{a}) = (q_\text{R},\text{X},I) & \delta(q_\text{b}, \text{b}) = (q_\text{b},\text{b},D) & \delta(q_\text{b}, \text{X}) = (q_\text{b},\text{X},D) &  \\
        \delta(q_\text{R}, \text{a}) = (q_\text{R},\text{a},I) & \delta(q_\text{R}, \text{b}) = (q_\text{R},\text{b},I) & \delta(q_\text{R}, \text{X}) = (q_\text{R},\text{X},I) & \delta(q_\text{R}, \Vtextvisiblespace) = (q_0,\Vtextvisiblespace,D) \\
    \end{matrix}
$$

\vspace{5mm}

\begin{adjustbox}{max height={.67\textheight},center}


\tikzset{every picture/.style={line width=0.75pt}} %set default line width to 0.75pt        

\begin{tikzpicture}[x=0.75pt,y=0.75pt,yscale=-1,xscale=1]
%uncomment if require: \path (0,446); %set diagram left start at 0, and has height of 446

%Shape: Triangle [id:dp9549836358247636] 
\draw  [fill={rgb, 255:red, 0; green, 0; blue, 0 }  ,fill opacity=1 ] (134.4,85) -- (141.6,82.2) -- (141.6,87.8) -- cycle ;
%Shape: Circle [id:dp39324764126664746] 
\draw   (101.2,85) .. controls (101.2,75) and (109.31,66.9) .. (119.3,66.9) .. controls (129.3,66.9) and (137.4,75) .. (137.4,85) .. controls (137.4,95) and (129.3,103.1) .. (119.3,103.1) .. controls (109.31,103.1) and (101.2,95) .. (101.2,85) -- cycle ;
%Straight Lines [id:da279972863636345] 
\draw    (379.6,70.2) -- (248.8,70.2) -- (189.6,70.2) ;
%Straight Lines [id:da8063028764122899] 
\draw    (379.2,99.6) -- (189.4,99.6) ;
%Straight Lines [id:da6194955219154272] 
\draw    (210,70.4) -- (210,99.6) ;
%Straight Lines [id:da9291793988668682] 
\draw    (239.6,70.8) -- (239.6,100) ;
%Straight Lines [id:da1753997056206802] 
\draw    (269.2,70.4) -- (269.2,99.6) ;
%Straight Lines [id:da7326138590955908] 
\draw    (299.2,70.4) -- (299.2,99.6) ;
%Straight Lines [id:da8827634409952425] 
\draw    (329.2,70.4) -- (329.2,99.6) ;
%Straight Lines [id:da08240500809926843] 
\draw    (359.2,70.8) -- (359.2,100) ;
%Shape: Rectangle [id:dp8505881042891044] 
\draw  [line width=2.25]  (359.2,99.6) -- (329.6,99.6) -- (329.6,70.4) -- (359.2,70.4) -- cycle ;
%Straight Lines [id:da2295782410185434] 
\draw    (344.8,117.75) -- (344.8,103.55) ;
\draw [shift={(344.8,100.55)}, rotate = 90] [fill={rgb, 255:red, 0; green, 0; blue, 0 }  ][line width=0.08]  [draw opacity=0] (8.93,-4.29) -- (0,0) -- (8.93,4.29) -- cycle    ;
%Straight Lines [id:da2131860495885145] 
\draw [line width=1.5]    (344.8,123.4) -- (344.8,106.15) ;


% Text Node
\draw (107.4,79.2) node [anchor=north west][inner sep=0.75pt]  [font=\small] [align=left] {\begin{minipage}[lt]{18.04pt}\setlength\topsep{0pt}
\begin{center}
$q_\text{a}$
\end{center}

\end{minipage}};
% Text Node
\draw (224.8,85.4) node   [align=left] {\begin{minipage}[lt]{20.13pt}\setlength\topsep{0pt}
\begin{center}
\Vtextvisiblespace
\end{center}

\end{minipage}};
% Text Node
\draw (254.4,85) node   [align=left] {\begin{minipage}[lt]{20.13pt}\setlength\topsep{0pt}
\begin{center}
X
\end{center}

\end{minipage}};
% Text Node
\draw (284.4,85) node   [align=left] {\begin{minipage}[lt]{20.13pt}\setlength\topsep{0pt}
\begin{center}
X
\end{center}

\end{minipage}};
% Text Node
\draw (314,85) node   [align=left] {\begin{minipage}[lt]{20.13pt}\setlength\topsep{0pt}
\begin{center}
X
\end{center}

\end{minipage}};
% Text Node
\draw (344.4,85.4) node   [align=left] {\begin{minipage}[lt]{20.13pt}\setlength\topsep{0pt}
\begin{center}
\Vtextvisiblespace
\end{center}

\end{minipage}};
% Text Node
\draw (374,85.4) node   [align=left] {\begin{minipage}[lt]{20.13pt}\setlength\topsep{0pt}
\begin{center}
...
\end{center}

\end{minipage}};
% Text Node
\draw (194.2,86) node   [align=left] {\begin{minipage}[lt]{20.13pt}\setlength\topsep{0pt}
\begin{center}
...
\end{center}

\end{minipage}};


\end{tikzpicture}

\end{adjustbox}

\end{frame}



\begin{frame}{Máquinas de Turing}

\textbf{Ejemplo 3.1}\textbf{.}
$$
    M_{\#\text{a}>\#\text{b}}=(\{q_0, q_\text{a}, q_\text{b}, q_\text{R}, q_\text{F}\}, \{\text{a}, \text{b}\}, \{\text{a}, \text{b}, \text{X}, \Vtextvisiblespace\}, \delta, q_0, \Vtextvisiblespace\:,\{q_\text{F}\})
$$
$$
    \begin{matrix}
        \delta(q_0, \text{a}) = (q_\text{a},\text{X},D) & \delta(q_0, \text{b}) = (q_\text{b},\text{X},D) & \delta(q_0, \text{X}) = (q_0,\text{X},D) &  \\
        \delta(q_\text{a}, \text{a}) = (q_\text{a},\text{a},D) & \delta(q_\text{a}, \text{b}) = (q_\text{R},\text{X},I) & \delta(q_\text{a}, \text{X}) = (q_\text{a},\text{X},D) & \underline{\delta(q_\text{a}, \Vtextvisiblespace) = (q_\text{F},\Vtextvisiblespace,S)} \\
        \delta(q_\text{b}, \text{a}) = (q_\text{R},\text{X},I) & \delta(q_\text{b}, \text{b}) = (q_\text{b},\text{b},D) & \delta(q_\text{b}, \text{X}) = (q_\text{b},\text{X},D) &  \\
        \delta(q_\text{R}, \text{a}) = (q_\text{R},\text{a},I) & \delta(q_\text{R}, \text{b}) = (q_\text{R},\text{b},I) & \delta(q_\text{R}, \text{X}) = (q_\text{R},\text{X},I) & \delta(q_\text{R}, \Vtextvisiblespace) = (q_0,\Vtextvisiblespace,D) \\
    \end{matrix}
$$

\vspace{5mm}

\begin{adjustbox}{max height={.67\textheight},center}


\tikzset{every picture/.style={line width=0.75pt}} %set default line width to 0.75pt        

\begin{tikzpicture}[x=0.75pt,y=0.75pt,yscale=-1,xscale=1]
%uncomment if require: \path (0,446); %set diagram left start at 0, and has height of 446

%Shape: Triangle [id:dp9549836358247636] 
\draw  [fill={rgb, 255:red, 0; green, 0; blue, 0 }  ,fill opacity=1 ] (134.4,85) -- (141.6,82.2) -- (141.6,87.8) -- cycle ;
%Shape: Circle [id:dp39324764126664746] 
\draw   (101.2,85) .. controls (101.2,75) and (109.31,66.9) .. (119.3,66.9) .. controls (129.3,66.9) and (137.4,75) .. (137.4,85) .. controls (137.4,95) and (129.3,103.1) .. (119.3,103.1) .. controls (109.31,103.1) and (101.2,95) .. (101.2,85) -- cycle ;
%Straight Lines [id:da279972863636345] 
\draw    (379.6,70.2) -- (248.8,70.2) -- (189.6,70.2) ;
%Straight Lines [id:da8063028764122899] 
\draw    (379.2,99.6) -- (189.4,99.6) ;
%Straight Lines [id:da6194955219154272] 
\draw    (210,70.4) -- (210,99.6) ;
%Straight Lines [id:da9291793988668682] 
\draw    (239.6,70.8) -- (239.6,100) ;
%Straight Lines [id:da1753997056206802] 
\draw    (269.2,70.4) -- (269.2,99.6) ;
%Straight Lines [id:da7326138590955908] 
\draw    (299.2,70.4) -- (299.2,99.6) ;
%Straight Lines [id:da8827634409952425] 
\draw    (329.2,70.4) -- (329.2,99.6) ;
%Straight Lines [id:da08240500809926843] 
\draw    (359.2,70.8) -- (359.2,100) ;
%Shape: Rectangle [id:dp8505881042891044] 
\draw  [line width=2.25]  (359.2,99.6) -- (329.6,99.6) -- (329.6,70.4) -- (359.2,70.4) -- cycle ;
%Straight Lines [id:da2295782410185434] 
\draw    (344.8,117.75) -- (344.8,103.55) ;
\draw [shift={(344.8,100.55)}, rotate = 90] [fill={rgb, 255:red, 0; green, 0; blue, 0 }  ][line width=0.08]  [draw opacity=0] (8.93,-4.29) -- (0,0) -- (8.93,4.29) -- cycle    ;
%Straight Lines [id:da2131860495885145] 
\draw [line width=1.5]    (344.8,123.4) -- (344.8,106.15) ;


% Text Node
\draw (107.4,79.2) node [anchor=north west][inner sep=0.75pt]  [font=\small] [align=left] {\begin{minipage}[lt]{18.04pt}\setlength\topsep{0pt}
\begin{center}
$q_\text{a}$
\end{center}

\end{minipage}};
% Text Node
\draw (224.8,85.4) node   [align=left] {\begin{minipage}[lt]{20.13pt}\setlength\topsep{0pt}
\begin{center}
\Vtextvisiblespace
\end{center}

\end{minipage}};
% Text Node
\draw (254.4,85) node   [align=left] {\begin{minipage}[lt]{20.13pt}\setlength\topsep{0pt}
\begin{center}
X
\end{center}

\end{minipage}};
% Text Node
\draw (284.4,85) node   [align=left] {\begin{minipage}[lt]{20.13pt}\setlength\topsep{0pt}
\begin{center}
X
\end{center}

\end{minipage}};
% Text Node
\draw (314,85) node   [align=left] {\begin{minipage}[lt]{20.13pt}\setlength\topsep{0pt}
\begin{center}
X
\end{center}

\end{minipage}};
% Text Node
\draw (344.4,85.4) node   [align=left] {\begin{minipage}[lt]{20.13pt}\setlength\topsep{0pt}
\begin{center}
\Vtextvisiblespace
\end{center}

\end{minipage}};
% Text Node
\draw (374,85.4) node   [align=left] {\begin{minipage}[lt]{20.13pt}\setlength\topsep{0pt}
\begin{center}
...
\end{center}

\end{minipage}};
% Text Node
\draw (194.2,86) node   [align=left] {\begin{minipage}[lt]{20.13pt}\setlength\topsep{0pt}
\begin{center}
...
\end{center}

\end{minipage}};


\end{tikzpicture}

\end{adjustbox}

\end{frame}


\begin{frame}{Máquinas de Turing}

\textbf{Ejemplo 3.1}\textbf{.}
$$
    M_{\#\text{a}>\#\text{b}}=(\{q_0, q_\text{a}, q_\text{b}, q_\text{R}, q_\text{F}\}, \{\text{a}, \text{b}\}, \{\text{a}, \text{b}, \text{X}, \Vtextvisiblespace\}, \delta, q_0, \Vtextvisiblespace\:,\{q_\text{F}\})
$$
$$
    \begin{matrix}
        \delta(q_0, \text{a}) = (q_\text{a},\text{X},D) & \delta(q_0, \text{b}) = (q_\text{b},\text{X},D) & \delta(q_0, \text{X}) = (q_0,\text{X},D) &  \\
        \delta(q_\text{a}, \text{a}) = (q_\text{a},\text{a},D) & \delta(q_\text{a}, \text{b}) = (q_\text{R},\text{X},I) & \delta(q_\text{a}, \text{X}) = (q_\text{a},\text{X},D) & \delta(q_\text{a}, \Vtextvisiblespace) = (q_\text{F},\Vtextvisiblespace,S) \\
        \delta(q_\text{b}, \text{a}) = (q_\text{R},\text{X},I) & \delta(q_\text{b}, \text{b}) = (q_\text{b},\text{b},D) & \delta(q_\text{b}, \text{X}) = (q_\text{b},\text{X},D) &  \\
        \delta(q_\text{R}, \text{a}) = (q_\text{R},\text{a},I) & \delta(q_\text{R}, \text{b}) = (q_\text{R},\text{b},I) & \delta(q_\text{R}, \text{X}) = (q_\text{R},\text{X},I) & \delta(q_\text{R}, \Vtextvisiblespace) = (q_0,\Vtextvisiblespace,D) \\
    \end{matrix}
$$

\vspace{5mm}

\begin{adjustbox}{max height={.67\textheight},center}


\tikzset{every picture/.style={line width=0.75pt}} %set default line width to 0.75pt        

\begin{tikzpicture}[x=0.75pt,y=0.75pt,yscale=-1,xscale=1]
%uncomment if require: \path (0,446); %set diagram left start at 0, and has height of 446

%Shape: Triangle [id:dp9549836358247636] 
\draw  [fill={rgb, 255:red, 0; green, 0; blue, 0 }  ,fill opacity=1 ] (134.4,85) -- (141.6,82.2) -- (141.6,87.8) -- cycle ;
%Shape: Circle [id:dp39324764126664746] 
\draw   (101.2,85) .. controls (101.2,75) and (109.31,66.9) .. (119.3,66.9) .. controls (129.3,66.9) and (137.4,75) .. (137.4,85) .. controls (137.4,95) and (129.3,103.1) .. (119.3,103.1) .. controls (109.31,103.1) and (101.2,95) .. (101.2,85) -- cycle ;
%Straight Lines [id:da279972863636345] 
\draw    (379.6,70.2) -- (248.8,70.2) -- (189.6,70.2) ;
%Straight Lines [id:da8063028764122899] 
\draw    (379.2,99.6) -- (189.4,99.6) ;
%Straight Lines [id:da6194955219154272] 
\draw    (210,70.4) -- (210,99.6) ;
%Straight Lines [id:da9291793988668682] 
\draw    (239.6,70.8) -- (239.6,100) ;
%Straight Lines [id:da1753997056206802] 
\draw    (269.2,70.4) -- (269.2,99.6) ;
%Straight Lines [id:da7326138590955908] 
\draw    (299.2,70.4) -- (299.2,99.6) ;
%Straight Lines [id:da8827634409952425] 
\draw    (329.2,70.4) -- (329.2,99.6) ;
%Straight Lines [id:da08240500809926843] 
\draw    (359.2,70.8) -- (359.2,100) ;
%Shape: Rectangle [id:dp8505881042891044] 
\draw  [line width=2.25]  (359.2,99.6) -- (329.6,99.6) -- (329.6,70.4) -- (359.2,70.4) -- cycle ;
%Straight Lines [id:da2295782410185434] 
\draw    (344.8,117.75) -- (344.8,103.55) ;
\draw [shift={(344.8,100.55)}, rotate = 90] [fill={rgb, 255:red, 0; green, 0; blue, 0 }  ][line width=0.08]  [draw opacity=0] (8.93,-4.29) -- (0,0) -- (8.93,4.29) -- cycle    ;
%Straight Lines [id:da2131860495885145] 
\draw [line width=1.5]    (344.8,123.4) -- (344.8,106.15) ;


% Text Node
\draw (107.4,79.2) node [anchor=north west][inner sep=0.75pt]  [font=\small] [align=left] {\begin{minipage}[lt]{18.04pt}\setlength\topsep{0pt}
\begin{center}
$q_\text{F}$
\end{center}

\end{minipage}};
% Text Node
\draw (224.8,85.4) node   [align=left] {\begin{minipage}[lt]{20.13pt}\setlength\topsep{0pt}
\begin{center}
\Vtextvisiblespace
\end{center}

\end{minipage}};
% Text Node
\draw (254.4,85) node   [align=left] {\begin{minipage}[lt]{20.13pt}\setlength\topsep{0pt}
\begin{center}
X
\end{center}

\end{minipage}};
% Text Node
\draw (284.4,85) node   [align=left] {\begin{minipage}[lt]{20.13pt}\setlength\topsep{0pt}
\begin{center}
X
\end{center}

\end{minipage}};
% Text Node
\draw (314,85) node   [align=left] {\begin{minipage}[lt]{20.13pt}\setlength\topsep{0pt}
\begin{center}
X
\end{center}

\end{minipage}};
% Text Node
\draw (344.4,85.4) node   [align=left] {\begin{minipage}[lt]{20.13pt}\setlength\topsep{0pt}
\begin{center}
\Vtextvisiblespace
\end{center}

\end{minipage}};
% Text Node
\draw (374,85.4) node   [align=left] {\begin{minipage}[lt]{20.13pt}\setlength\topsep{0pt}
\begin{center}
...
\end{center}

\end{minipage}};
% Text Node
\draw (194.2,86) node   [align=left] {\begin{minipage}[lt]{20.13pt}\setlength\topsep{0pt}
\begin{center}
...
\end{center}

\end{minipage}};


\end{tikzpicture}

\end{adjustbox}

\end{frame}


\begin{frame}{Máquinas de Turing}

\textbf{Ejemplo 3.1}\textbf{.}
$$
    M_{\#\text{a}>\#\text{b}}=(\{q_0, q_\text{a}, q_\text{b}, q_\text{R}, q_\text{F}\}, \{\text{a}, \text{b}\}, \{\text{a}, \text{b}, \text{X}, \Vtextvisiblespace\}, \delta, q_0, \Vtextvisiblespace\:,\{\underline{q_\text{F}}\})
$$
$$
    \begin{matrix}
        \delta(q_0, \text{a}) = (q_\text{a},\text{X},D) & \delta(q_0, \text{b}) = (q_\text{b},\text{X},D) & \delta(q_0, \text{X}) = (q_0,\text{X},D) &  \\
        \delta(q_\text{a}, \text{a}) = (q_\text{a},\text{a},D) & \delta(q_\text{a}, \text{b}) = (q_\text{R},\text{X},I) & \delta(q_\text{a}, \text{X}) = (q_\text{a},\text{X},D) & \delta(q_\text{a}, \Vtextvisiblespace) = (q_\text{F},\Vtextvisiblespace,S) \\
        \delta(q_\text{b}, \text{a}) = (q_\text{R},\text{X},I) & \delta(q_\text{b}, \text{b}) = (q_\text{b},\text{b},D) & \delta(q_\text{b}, \text{X}) = (q_\text{b},\text{X},D) &  \\
        \delta(q_\text{R}, \text{a}) = (q_\text{R},\text{a},I) & \delta(q_\text{R}, \text{b}) = (q_\text{R},\text{b},I) & \delta(q_\text{R}, \text{X}) = (q_\text{R},\text{X},I) & \delta(q_\text{R}, \Vtextvisiblespace) = (q_0,\Vtextvisiblespace,D) \\
    \end{matrix}
$$

\vspace{5mm}

\begin{adjustbox}{max height={.67\textheight},center}


\tikzset{every picture/.style={line width=0.75pt}} %set default line width to 0.75pt        

\begin{tikzpicture}[x=0.75pt,y=0.75pt,yscale=-1,xscale=1]
%uncomment if require: \path (0,446); %set diagram left start at 0, and has height of 446

%Shape: Circle [id:dp39324764126664746] 
\draw  [fill={rgb, 255:red, 184; green, 233; blue, 134 }  ,fill opacity=1 ] (101.2,85) .. controls (101.2,75) and (109.31,66.9) .. (119.3,66.9) .. controls (129.3,66.9) and (137.4,75) .. (137.4,85) .. controls (137.4,95) and (129.3,103.1) .. (119.3,103.1) .. controls (109.31,103.1) and (101.2,95) .. (101.2,85) -- cycle ;
%Straight Lines [id:da279972863636345] 
\draw    (379.6,70.2) -- (248.8,70.2) -- (189.6,70.2) ;
%Straight Lines [id:da8063028764122899] 
\draw    (379.2,99.6) -- (189.4,99.6) ;
%Straight Lines [id:da6194955219154272] 
\draw    (210,70.4) -- (210,99.6) ;
%Straight Lines [id:da9291793988668682] 
\draw    (239.6,70.8) -- (239.6,100) ;
%Straight Lines [id:da1753997056206802] 
\draw    (269.2,70.4) -- (269.2,99.6) ;
%Straight Lines [id:da7326138590955908] 
\draw    (299.2,70.4) -- (299.2,99.6) ;
%Straight Lines [id:da8827634409952425] 
\draw    (329.2,70.4) -- (329.2,99.6) ;
%Straight Lines [id:da08240500809926843] 
\draw    (359.2,70.8) -- (359.2,100) ;
%Shape: Rectangle [id:dp8505881042891044] 
\draw  [line width=2.25]  (359.2,99.6) -- (329.6,99.6) -- (329.6,70.4) -- (359.2,70.4) -- cycle ;
%Straight Lines [id:da2295782410185434] 
\draw    (344.8,117.75) -- (344.8,103.55) ;
\draw [shift={(344.8,100.55)}, rotate = 90] [fill={rgb, 255:red, 0; green, 0; blue, 0 }  ][line width=0.08]  [draw opacity=0] (8.93,-4.29) -- (0,0) -- (8.93,4.29) -- cycle    ;
%Straight Lines [id:da2131860495885145] 
\draw [line width=1.5]    (344.8,123.4) -- (344.8,106.15) ;


% Text Node
\draw (107.4,79.2) node [anchor=north west][inner sep=0.75pt]  [font=\small] [align=left] {\begin{minipage}[lt]{18.04pt}\setlength\topsep{0pt}
\begin{center}
$q_\text{F}$
\end{center}

\end{minipage}};
% Text Node
\draw (224.8,85.4) node   [align=left] {\begin{minipage}[lt]{20.13pt}\setlength\topsep{0pt}
\begin{center}
\Vtextvisiblespace
\end{center}

\end{minipage}};
% Text Node
\draw (254.4,85) node   [align=left] {\begin{minipage}[lt]{20.13pt}\setlength\topsep{0pt}
\begin{center}
X
\end{center}

\end{minipage}};
% Text Node
\draw (284.4,85) node   [align=left] {\begin{minipage}[lt]{20.13pt}\setlength\topsep{0pt}
\begin{center}
X
\end{center}

\end{minipage}};
% Text Node
\draw (314,85) node   [align=left] {\begin{minipage}[lt]{20.13pt}\setlength\topsep{0pt}
\begin{center}
X
\end{center}

\end{minipage}};
% Text Node
\draw (344.4,85.4) node   [align=left] {\begin{minipage}[lt]{20.13pt}\setlength\topsep{0pt}
\begin{center}
\Vtextvisiblespace
\end{center}

\end{minipage}};
% Text Node
\draw (374,85.4) node   [align=left] {\begin{minipage}[lt]{20.13pt}\setlength\topsep{0pt}
\begin{center}
...
\end{center}

\end{minipage}};
% Text Node
\draw (194.2,86) node   [align=left] {\begin{minipage}[lt]{20.13pt}\setlength\topsep{0pt}
\begin{center}
...
\end{center}

\end{minipage}};
%Shape: Triangle [id:dp9549836358247636] 
\draw  [fill={rgb, 255:red, 0; green, 0; blue, 0 }  ,fill opacity=1 ] (134.4,85) -- (141.6,82.2) -- (141.6,87.8) -- cycle ;


\end{tikzpicture}

\end{adjustbox}

\end{frame}


\begin{frame}[fragile]{Programas en Python}

\begin{lstlisting}[language=Python]
import utilidades

def mas_a_que_b(palabra):
    # comprobamos si la codificación es correcta, es decir (si es del alfabeto {'a', 'b'}), si no es correcta devolvemos 'no'
    if not utilidades.en_alfabeto(palabra, {'a', 'b'}):
        return 'no'

    if palabra.count('a') > palabra.count('b'):
        return 'sí'

    return 'no'
\end{lstlisting}
\vspace*{-4mm}
\begin{center}
{\small Programa 4.1.: \texttt{mas\_a\_que\_b.py}}
\end{center}

\end{frame}


\begin{frame}<presentation:0>{Programas en Python}

\begin{table}[H]
\centering
\begin{tabular}{@{}lc@{}}
\toprule
Comando  & Salida \\ \midrule
\texttt{mas\_a\_que\_b(\palabra{abaab})} & \palabra{sí} \\
\texttt{mas\_a\_que\_b(\palabra{abbab})} & \palabra{no} \\
\texttt{mas\_a\_que\_b(\palabra{})} & \palabra{no} \\
\texttt{mas\_a\_que\_b(\palabra{abaabc})} & \palabra{no} \\ \bottomrule
\end{tabular}
\end{table}
\begin{center}
{\small Tabla 4.1.: Ejemplos de salidas de \texttt{mas\_a\_que\_b.py}}
\end{center}

\end{frame}


\begin{frame}[fragile]{Equivalencia}


\begin{itemize}
    \item Tratamos con programas en Python que tienen como entrada y como salida \emph{strings} (\hle{programas SISO}).
    \item Sería idóneo poder usar programas en Python como la representación matemática de un algoritmo.
\end{itemize}

\vspace{5mm}
\pause

\textbf{Teorema 3.1} (Equivalencia entre programas en Python y máquinas de Turing)\textbf{.}

\emph{Dado un programa en Python $P$, existe una máquina de Turing (de una sola cinta) $M$ que \hle{computa la misma función} que $P$.}

\emph{Dada una máquina de Turing (de una sola cinta) $M$, existe un programa en Python $P$ que \hle{computa la misma función} que $M$.}

\end{frame}

\begin{frame}<presentation:0>{Equivalencia}

\vspace{5mm}
\begin{adjustbox}{max width={.9\textwidth},center}

\tikzset{every picture/.style={line width=0.75pt}} %set default line width to 0.75pt        

\begin{tikzpicture}[x=0.75pt,y=0.75pt,yscale=-1,xscale=1]
%uncomment if require: \path (0,489); %set diagram left start at 0, and has height of 489

%Rounded Rect [id:dp38098198089187774] 
\draw   (241,29) .. controls (241,24.58) and (244.58,21) .. (249,21) -- (501.5,21) .. controls (505.92,21) and (509.5,24.58) .. (509.5,29) -- (509.5,53) .. controls (509.5,57.42) and (505.92,61) .. (501.5,61) -- (249,61) .. controls (244.58,61) and (241,57.42) .. (241,53) -- cycle ;
%Rounded Rect [id:dp7442940701895104] 
\draw   (410.5,104) .. controls (410.5,99.58) and (414.08,96) .. (418.5,96) -- (621.5,96) .. controls (625.92,96) and (629.5,99.58) .. (629.5,104) -- (629.5,128) .. controls (629.5,132.42) and (625.92,136) .. (621.5,136) -- (418.5,136) .. controls (414.08,136) and (410.5,132.42) .. (410.5,128) -- cycle ;
%Rounded Rect [id:dp05734792720775661] 
\draw   (410.5,180) .. controls (410.5,175.58) and (414.08,172) .. (418.5,172) -- (621.5,172) .. controls (625.92,172) and (629.5,175.58) .. (629.5,180) -- (629.5,204) .. controls (629.5,208.42) and (625.92,212) .. (621.5,212) -- (418.5,212) .. controls (414.08,212) and (410.5,208.42) .. (410.5,204) -- cycle ;
%Rounded Rect [id:dp6203792922315816] 
\draw   (121.5,105) .. controls (121.5,100.58) and (125.08,97) .. (129.5,97) -- (332.5,97) .. controls (336.92,97) and (340.5,100.58) .. (340.5,105) -- (340.5,129) .. controls (340.5,133.42) and (336.92,137) .. (332.5,137) -- (129.5,137) .. controls (125.08,137) and (121.5,133.42) .. (121.5,129) -- cycle ;
%Rounded Rect [id:dp5638157358715545] 
\draw   (121.5,181) .. controls (121.5,176.58) and (125.08,173) .. (129.5,173) -- (332.5,173) .. controls (336.92,173) and (340.5,176.58) .. (340.5,181) -- (340.5,205) .. controls (340.5,209.42) and (336.92,213) .. (332.5,213) -- (129.5,213) .. controls (125.08,213) and (121.5,209.42) .. (121.5,205) -- cycle ;
%Straight Lines [id:da13357301440091862] 
\draw    (479.2,61.14) .. controls (481.55,60.97) and (482.81,62.06) .. (482.98,64.41) .. controls (483.14,66.76) and (484.4,67.86) .. (486.75,67.69) .. controls (489.1,67.52) and (490.36,68.62) .. (490.53,70.97) .. controls (490.7,73.32) and (491.96,74.41) .. (494.31,74.24) .. controls (496.66,74.07) and (497.92,75.17) .. (498.08,77.52) .. controls (498.25,79.87) and (499.51,80.97) .. (501.86,80.8) .. controls (504.21,80.63) and (505.47,81.72) .. (505.64,84.07) .. controls (505.8,86.42) and (507.06,87.52) .. (509.41,87.35) -- (511.45,89.12) -- (517.49,94.36) ;
\draw [shift={(519,95.67)}, rotate = 220.95] [color={rgb, 255:red, 0; green, 0; blue, 0 }  ][line width=0.75]    (10.93,-3.29) .. controls (6.95,-1.4) and (3.31,-0.3) .. (0,0) .. controls (3.31,0.3) and (6.95,1.4) .. (10.93,3.29)   ;
%Straight Lines [id:da7755341344632458] 
\draw    (519.67,136.05) .. controls (521.34,137.72) and (521.34,139.38) .. (519.67,141.05) .. controls (518,142.72) and (518,144.38) .. (519.67,146.05) .. controls (521.34,147.72) and (521.34,149.38) .. (519.67,151.05) .. controls (518,152.72) and (518,154.38) .. (519.67,156.05) .. controls (521.34,157.72) and (521.34,159.38) .. (519.67,161.05) -- (519.67,161.67) -- (519.67,169.67) ;
\draw [shift={(519.67,171.67)}, rotate = 270] [color={rgb, 255:red, 0; green, 0; blue, 0 }  ][line width=0.75]    (10.93,-3.29) .. controls (6.95,-1.4) and (3.31,-0.3) .. (0,0) .. controls (3.31,0.3) and (6.95,1.4) .. (10.93,3.29)   ;
%Straight Lines [id:da5719619240034488] 
\draw    (229.83,172.67) .. controls (228.16,171) and (228.16,169.34) .. (229.83,167.67) .. controls (231.5,166) and (231.5,164.34) .. (229.83,162.67) .. controls (228.16,161) and (228.16,159.34) .. (229.83,157.67) .. controls (231.5,156) and (231.5,154.34) .. (229.83,152.67) .. controls (228.16,151) and (228.16,149.34) .. (229.83,147.67) -- (229.83,147.19) -- (229.83,139.19) ;
\draw [shift={(229.83,137.19)}, rotate = 90] [color={rgb, 255:red, 0; green, 0; blue, 0 }  ][line width=0.75]    (10.93,-3.29) .. controls (6.95,-1.4) and (3.31,-0.3) .. (0,0) .. controls (3.31,0.3) and (6.95,1.4) .. (10.93,3.29)   ;
%Straight Lines [id:da9776591080439316] 
\draw    (410.6,192.34) .. controls (408.93,194.01) and (407.27,194.01) .. (405.6,192.34) .. controls (403.93,190.67) and (402.27,190.67) .. (400.6,192.34) .. controls (398.93,194.01) and (397.27,194.01) .. (395.6,192.34) .. controls (393.93,190.67) and (392.27,190.67) .. (390.6,192.34) .. controls (388.93,194.01) and (387.27,194.01) .. (385.6,192.34) .. controls (383.93,190.67) and (382.27,190.67) .. (380.6,192.34) .. controls (378.93,194.01) and (377.27,194.01) .. (375.6,192.34) .. controls (373.93,190.67) and (372.27,190.67) .. (370.6,192.34) .. controls (368.93,194.01) and (367.27,194.01) .. (365.6,192.34) .. controls (363.93,190.67) and (362.27,190.67) .. (360.6,192.34) .. controls (358.93,194.01) and (357.27,194.01) .. (355.6,192.34) -- (351,192.34) -- (343,192.34) ;
\draw [shift={(341,192.34)}, rotate = 360] [color={rgb, 255:red, 0; green, 0; blue, 0 }  ][line width=0.75]    (10.93,-3.29) .. controls (6.95,-1.4) and (3.31,-0.3) .. (0,0) .. controls (3.31,0.3) and (6.95,1.4) .. (10.93,3.29)   ;
%Straight Lines [id:da4761638716364047] 
\draw    (230.33,96.67) .. controls (230.51,94.32) and (231.77,93.23) .. (234.12,93.41) .. controls (236.47,93.58) and (237.74,92.49) .. (237.91,90.14) .. controls (238.08,87.79) and (239.35,86.7) .. (241.7,86.88) .. controls (244.05,87.06) and (245.32,85.97) .. (245.49,83.62) .. controls (245.66,81.27) and (246.93,80.18) .. (249.28,80.36) .. controls (251.63,80.53) and (252.9,79.44) .. (253.07,77.09) .. controls (253.24,74.74) and (254.51,73.65) .. (256.86,73.83) .. controls (259.21,74.01) and (260.48,72.92) .. (260.65,70.57) -- (264.02,67.66) -- (270.08,62.44) ;
\draw [shift={(271.6,61.14)}, rotate = 139.27] [color={rgb, 255:red, 0; green, 0; blue, 0 }  ][line width=0.75]    (10.93,-3.29) .. controls (6.95,-1.4) and (3.31,-0.3) .. (0,0) .. controls (3.31,0.3) and (6.95,1.4) .. (10.93,3.29)   ;

% Text Node
\draw (375.25,41) node  [font=\normalsize] [align=left] {\begin{minipage}[lt]{171.7pt}\setlength\topsep{0pt}
\begin{center}
máquina de Turing de una sola cinta
\end{center}

\end{minipage}};
% Text Node
\draw (520,116) node  [font=\normalsize] [align=left] {\begin{minipage}[lt]{138.04pt}\setlength\topsep{0pt}
\begin{center}
máquina de Turing multicinta
\end{center}

\end{minipage}};
% Text Node
\draw (520,192) node  [font=\normalsize] [align=left] {\begin{minipage}[lt]{138.04pt}\setlength\topsep{0pt}
\begin{center}
máquina de acceso aleatorio
\end{center}

\end{minipage}};
% Text Node
\draw (231,117) node  [font=\normalsize] [align=left] {\begin{minipage}[lt]{138.04pt}\setlength\topsep{0pt}
\begin{center}
programa en Python
\end{center}

\end{minipage}};
% Text Node
\draw (231,193) node  [font=\normalsize] [align=left] {\begin{minipage}[lt]{138.04pt}\setlength\topsep{0pt}
\begin{center}
ordenador moderno
\end{center}

\end{minipage}};
% Text Node
\draw (551.17,77.84) node  [font=\small] [align=left] {\begin{minipage}[lt]{43.29pt}\setlength\topsep{0pt}
\end{minipage}};
% Text Node
\draw (561.17,153.84) node  [font=\small] [align=left] {\begin{minipage}[lt]{43.29pt}\setlength\topsep{0pt}
\end{minipage}};
% Text Node
\draw (375.8,177.5) node  [font=\small] [align=left] {\begin{minipage}[lt]{47.33pt}\setlength\topsep{0pt}
\begin{center}
\end{center}

\end{minipage}};
% Text Node
\draw (189.67,154.84) node  [font=\small] [align=left] {\begin{minipage}[lt]{43.29pt}\setlength\topsep{0pt}
\begin{flushright}
\end{flushright}

\end{minipage}};
% Text Node
\draw (199.67,78.84) node  [font=\small] [align=left] {\begin{minipage}[lt]{43.29pt}\setlength\topsep{0pt}
\begin{flushright}
\end{flushright}

\end{minipage}};


\end{tikzpicture}
\end{adjustbox}

\begin{center}
{\small Figura 3.3.: Cadena de simulaciones}
\end{center}

\end{frame}


\begin{frame}{Equivalencia}

\textbf{Consecuencias.}

\begin{itemize}
    \item Las máquinas de Turing son \hl{muy potentes} computacionalmente.

    \pause
    \item \hl{\textbf{Tesis de Church-Turing.}} \emph{Toda función calculable (algoritmo) puede calcularse mediante una máquina de Turing}.
    
\end{itemize}

\end{frame}


\section{El problema de la parada}

\begin{frame}{Problemas decidibles}

\textbf{Definición 4.1} (Problema computacional)\textbf{.}
$${ \Large
\begin{matrix}
    F & : & \;\;\; X  \;\;\; & \longrightarrow & \;\;\; Y \;\;\; \\
\end{matrix}
}$$
\begin{center}
\vspace*{-2mm}
    \hspace{11mm} entradas \hspace{9mm} soluciones
\end{center}
\vspace{5mm}
\textbf{Definición 4.2} (Problema de decisión)\textbf{.}
$${ \Large
\begin{matrix}
    F & : & \;\;\; X  \;\;\; & \longrightarrow & \;\;\; \{\text{\emph{`sí'}}, \text{\emph{`no'}}\} \;\;\; \\
\end{matrix}
}$$
\vspace{2mm}

\pause

\textbf{Proposición B.2.} \emph{Para cada problema} $F:X\longrightarrow Y$\emph{, existe un problema equivalente} $\widehat{F}:S\longrightarrow S$.

\end{frame}

\begin{frame}{Problemas decidibles}

\begin{framed}
$$\text{\large \textsc{MásAQueB}}$$

\begin{itemize}
    \item \textbf{Entrada:} una cadena de caracteres del alfabeto $\{\palabra{a}, \palabra{b}\}$.
    \item \textbf{Solución:} si la cadena de caracteres tiene más \palabra{a} que \palabra{b}.
\end{itemize}
\end{framed}
\begin{center}
{\small Problema 4.4.: \textsc{MásAQueB}}
\end{center}

\end{frame}


\begin{frame}{Problemas decidibles}

\textbf{Definición 4.4} (Resolver y decidir)\textbf{.}

Dados problema $F:S\longrightarrow S$, programa $P$, $P$ \hle{resuelve} $F$ si para cada entrada $I\in S$:

\begin{itemize}
    \item $P$ \hl{para}.
    \item $P$ devuelve un \emph{string} tal que \hl{$P(I)=F(I)$}.
\end{itemize}

Si $F$ es de decisión, $P$ \hle{decide} $F$.

\vspace{5mm}

\pause

\textbf{Definición 4.6} (Problema computable y problema decidible)\textbf{.}

\begin{itemize}
    \item $F$ es \hle{computable} si existe $P$ que lo resuelve.
    \item $D$ es \hle{decidible} si existe $P$ que lo decide.
\end{itemize}


\end{frame}

\begin{frame}<1-3>[label=masaqueb-decidible]{Problemas decidibles}

\textbf{Proposición 4.1.} \emph{El problema} \textsc{MásAQueB} \emph{es decidible.}

\vspace{5mm}

\pause

\textbf{\emph{Demostración.}} El programa \texttt{mas\_a\_que\_b.py} decide \textsc{MásAQueB}, puesto que para para todas las entradas $I\in S$ y:

\begin{itemize}
    \pause
    \item Si $I \notin \{\palabra{a}, \palabra{b}\}$, es $P(I)=\textsc{MásAQueB}(I)=\palabra{no}$ (entrada con codif. incorrecta).

    \pause
    \item Si $I \in \{\palabra{a}, \palabra{b}\}$:
    
    \begin{itemize}
        \item Si $I$ tiene más \palabra{a} que \palabra{b}, es $P(I)=\textsc{MásAQueB}(I)=\palabra{sí}$.
        \pause
        \item Si $I$ no tiene más \palabra{a} que \palabra{b}, es $P(I)=\textsc{MásAQueB}(I)=\palabra{no}$.\vspace{-8mm}
    \end{itemize}
\end{itemize}
\pause[5]{\hfill\square}

\end{frame}


\begin{frame}[fragile]{Problemas decidibles}

\begin{lstlisting}[language=Python]
import utilidades

def mas_a_que_b(palabra):
    # comprobamos si la codificación es correcta, es decir (si es del alfabeto {'a', 'b'}), si no es correcta devolvemos 'no'
    (*@\hl{if not utilidades.en\_alfabeto(palabra, \{\textquotesingle{}a\textquotesingle{}, \textquotesingle{}b\textquotesingle{}\}):}@*)
        (*@\hl{return 'no'}@*)

    if palabra.count('a') > palabra.count('b'):
        return 'sí'

    return 'no'
\end{lstlisting}
\vspace*{-4mm}
\begin{center}
{\small Programa 4.1.: \texttt{mas\_a\_que\_b.py}}
\end{center}

\end{frame}


\againframe<4>{masaqueb-decidible}

\begin{frame}[fragile]{Problemas decidibles}

\begin{lstlisting}[language=Python]
import utilidades

def mas_a_que_b(palabra):
    # comprobamos si la codificación es correcta, es decir (si es del alfabeto {'a', 'b'}), si no es correcta devolvemos 'no'
    if not utilidades.en_alfabeto(palabra, {'a', 'b'}):
        return 'no'

    (*@\hl{if palabra.count(\textquotesingle{}a\textquotesingle{}) > palabra.count(\textquotesingle{}b\textquotesingle{}): }@*)
        (*@\hl{return 'sí'}@*)

    return 'no'
\end{lstlisting}
\vspace*{-4mm}
\begin{center}
{\small Programa 4.1.: \texttt{mas\_a\_que\_b.py}}
\end{center}

\end{frame}

\againframe<5>{masaqueb-decidible}


\begin{frame}[fragile]{Problemas decidibles}

\begin{lstlisting}[language=Python]
import utilidades

def mas_a_que_b(palabra):
    # comprobamos si la codificación es correcta, es decir (si es del alfabeto {'a', 'b'}), si no es correcta devolvemos 'no'
    if not utilidades.en_alfabeto(palabra, {'a', 'b'}):
        return 'no'

    if palabra.count('a') > palabra.count('b'):
        return 'sí'

    (*@\hl{return 'no'}  @*)
\end{lstlisting}
\vspace*{-4mm}
\begin{center}
{\small Programa 4.1.: \texttt{mas\_a\_que\_b.py}}
\end{center}

\end{frame}

\againframe<5>{masaqueb-decidible}


\begin{frame}{Problemas semidecidibles}

\textbf{Definición 4.4} (Semidecidir)\textbf{.}

Dados problema de decisión $D:S\longrightarrow S$, programa $P$, $P$ \hle{semidecide} $F$ si para cada entrada $I\in S$:

\begin{itemize}
    \item si $D(I)=\palabra{sí}$, $P$ \hl{para} y $P(I)=\palabra{sí}$.
    \item si $D(I)=\palabra{no}$, $P(I)$ es \hl{indefinido} o $P(I)=\palabra{no}$.
\end{itemize}

\vspace{5mm}

\textbf{Definición 4.6} (Problema semidecidible)\textbf{.}

$D$ es \hle{semidecidible} si existe $P$ que lo semidecide.

\end{frame}


\begin{frame}{Máquina universal}

\textbf{Existencia de máquina universal.}

Existe una \hle{máquina universal} en Python, un programa \texttt{maquina\_universal.py} con:

\begin{itemize}
    \item \textbf{Entrada:} programa \texttt{programa}, entrada \texttt{entrada}.
    \item \textbf{Salida:} resultado de ejecutar programa \texttt{programa} con entrada \texttt{entrada}.
\end{itemize}

\end{frame}


\begin{frame}{Máquina universal}


\begin{table}[H]
\centering
\begin{tabular}{@{}lc@{}}
\toprule
Comando  & Salida \\ \midrule
\texttt{maquina\_universal(} & \texttt{\textquotesingle}\texttt{error}$...$\texttt{\textquotesingle} \\
\hspace{16pt}\texttt{\palabra{no es un programa},}\\
\hspace{16pt}\texttt{\palabra{una entrada cualquiera})} \\[3mm]
\texttt{maquina\_universal(} & \palabra{sí} \\
\hspace{16pt}\texttt{leer(\palabra{./mas\_a\_que\_b.py}), \palabra{abaab})}\\[3mm]
\texttt{maquina\_universal(} & \palabra{no} \\
\hspace{16pt}\texttt{leer(\palabra{./mas\_a\_que\_b.py}),}\\
\hspace{16pt}\texttt{leer(\palabra{./mas\_a\_que\_b.py}))} \\ \bottomrule
\end{tabular}
\end{table}
\begin{center}
{\small Tabla 4.1.: Ejemplos de salidas de \texttt{maquina\_universal.py}}
\end{center}

\end{frame}


\begin{frame}<presentation:0>{Problemas no decidibles}

\begin{framed}
$$\text{\large \textsc{C-Diagonal}}$$

\begin{itemize}
    \item \textbf{Entrada:} un programa $P$.
    \item \textbf{Solución:} \palabra{no} si $P(P)$ está bien definido y $P(P)=\text{\palabra{sí}}$, \palabra{sí} en caso contrario.
\end{itemize}
\end{framed}
\begin{center}
{\small Problema 4.8.: \textsc{C-Diagonal}}
\end{center}

\end{frame}


\begin{frame}<presentation:0>{Problemas no decidibles}

\begin{table}[H]
\centering
\begin{tabular}{@{}lc@{}}
\toprule
Comando  & Salida \\ \midrule
\texttt{c\_diagonal(\palabra{no es un programa})} & \palabra{sí} \\
\texttt{c\_diagonal(\texttt{leer}(\palabra{./mas\_a\_que\_b.py}))} & \palabra{no} \\
\texttt{c\_diagonal(\texttt{leer}(\palabra{./si.py}))} & \palabra{no} \\
\texttt{c\_diagonal(\texttt{leer}(\palabra{./c\_diagonal.py}))} & \textbf{¿}\palabra{sí} o \palabra{no}\textbf{?}\\ \bottomrule
\end{tabular}
\end{table}
\begin{center}
{\small Tabla 4.5.: Ejemplos de salidas de \texttt{c\_diagonal.py}}
\end{center}

\end{frame}



\begin{frame}<presentation:0>{Problemas no decidibles}

\textbf{Proposición 4.2.} \emph{El problema} \textsc{C-Diagonal} \emph{es no decidible.}

\textbf{Demostración.} Si el problema es decidible, existe un programa \texttt{c\_diagonal.py} que lo decide, pero:

\begin{center}
    \texttt{c\_diagonal(\texttt{leer}(\palabra{./c\_diagonal.py}))}

    $\updownarrow$

    ¿Es cierto que \texttt{c\_diagonal.py} \emph{no} devuelve \palabra{sí} cuando se ejecuta a sí mismo?
\end{center}

Pero la respuesta a esta pregunta no puede ser ni afirmativa ni negativa. \vspace{-6.5mm}

\hfill\square

\end{frame}


\begin{frame}{Problemas no decidibles}

\begin{framed}
$$\text{\large \textsc{Universal}}$$

\begin{itemize}
    \item \textbf{Entrada:} un programa $P$ y una entrada $I$.
    \item \textbf{Solución:} \palabra{sí} si $P$ está bien definido y $P(I)=\text{\palabra{sí}}$, \palabra{no} en caso contrario.
\end{itemize}
\end{framed}
\begin{center}
{\small Problema 4.6.: \textsc{Universal}}
\end{center}

\end{frame}


\begin{frame}<presentation:0>{Problemas no decidibles}

\begin{table}[H]
\centering
\begin{tabular}{@{}lc@{}}
\toprule
Comando  & Salida \\ \midrule
\texttt{universal(} & \palabra{no} \\
\hspace{16pt}\texttt{\palabra{no es un programa},}\\
\hspace{16pt}\texttt{\palabra{una entrada cualquiera})} \\[3mm]
\texttt{universal(} & \palabra{sí} \\
\hspace{16pt}\texttt{leer(\palabra{./mas\_a\_que\_b.py}), \palabra{aaba})} \\[3mm]
\pause[2]{\texttt{universal(}} & \pause[3]{\textbf{¿}\palabra{sí} o \palabra{no}\textbf{?}} \\
\pause[2]{\hspace{16pt}\texttt{leer(\palabra{./universal.py}), leer(\palabra{./universal.py}))}} \\ \bottomrule
\end{tabular}
\end{table}
\begin{center}
{\small Tabla 4.3. (modificada): Ejemplos de salidas de \texttt{universal.py}}
\end{center}

\end{frame}


\begin{frame}[fragile]{Problemas no decidibles}

\textbf{Proposición 4.3.} \emph{El problema} \textsc{Universal} \emph{es no decidible.}

\end{frame}


\begin{frame}<presentation:0>[fragile]{Problemas no decidibles}

\textbf{Proposición 4.3.} \emph{El problema} \textsc{Universal} \emph{es no decidible.}

\textbf{Demostración.} Si \textsc{Universal} fuese decidible, podríamos escribir un programa \texttt{c\_diagonal.py} que decidiese \textsc{C-Diagonal}:

\begin{lstlisting}[language=Python]
from universal import universal

def c_diagonal(programa):
    salida = universal(programa, programa)

    if salida == 'sí':
        return 'no'
    else:
        return 'sí'
\end{lstlisting}
 \vspace{-8.5mm}

\hfill\square

\end{frame}

\begin{frame}{Reducciones}

\textbf{Definición 4.7} (Reducción de Turing)\textbf{.}

Sean problemas $F:S\longrightarrow S$ y $G:S\longrightarrow S$.
{\Large $$F \;\;\; \leq_T \;\;\; G$$}
\begin{center}
\vspace*{-2mm}
    $F$ se \hle{reduce} a $G$
\end{center}
Si asumiendo $\exists P_G$ resuelve $G$ $\rightarrow$ $\exists P_F$ resuelve $F$.

\vspace{5mm}

\pause

\textbf{Proposición 4.5.} {\itshape Sean $F$ y $G$ problemas, y sea $F\leq_T G$. Entonces:
\begin{itemize}
    \item[a)] $G$ decidible $\Rightarrow$ $F$ decidible
    \item[b)] $F$ no decidible $\Rightarrow$ $G$ no decidible
\end{itemize}}

\end{frame}

\begin{frame}{El problema de la parada}

\begin{framed}
$$\text{\large \textsc{Parada}}$$

\begin{itemize}
    \item \textbf{Entrada:} un programa $P$ y una entrada $I$.
    \item \textbf{Solución:} si el programa $P$ para con entrada $I$.
\end{itemize}
\end{framed}
\begin{center}
{\small Problema 4.9.: \textsc{Parada}}
\end{center}

\end{frame}

\begin{frame}{El problema de la parada}

\textbf{Proposición 4.6.} \emph{El problema} \textsc{Parada} \emph{es no decidible.}

\vspace{5mm}

\pause

\textbf{Demostración.} Hacemos la reducción $\textsc{Universal}\leq_T\textsc{Parada}$.

\pause

\emph{Idea.} Transformar un programa $P$ en otro $P'$ de modo que:
$$\begin{matrix}
    $P$\text{ devuelve \palabra{sí}} & \rightarrow & $P'$\text{ para}\\
    $P$\text{ no devuelve \palabra{sí}} & \rightarrow & $P'$\text{ no para}\\
\end{matrix}$$

\end{frame}



\begin{frame}<presentation:0>{El problema de la parada}

\begin{center}
    \large{$\textsc{Universal}\leq_T\textsc{Parada}$}
\end{center}

$$\begin{matrix}
    $P$\text{ devuelve \palabra{sí}} & \rightarrow & $P'$\text{ para}\\
    $P$\text{ no devuelve \palabra{sí}} & \rightarrow & $P'$\text{ no para}\\
\end{matrix}$$

\begin{columns}
\begin{column}{0.5\textwidth}
\begin{framed}
$$\text{\large \textsc{Universal}}$$

\begin{itemize}
    \item \textbf{Entrada:} programa $P$ y entrada $I$.
    \item \textbf{Solución:} \palabra{sí} si $P$ está bien definido y $P(I)=\text{\palabra{sí}}$, \palabra{no} en otro caso.
\end{itemize}
\end{framed}
\end{column}
\begin{column}{0.5\textwidth}  %%<--- here
\begin{framed}
$$\text{\large \textsc{Parada}}$$

\begin{itemize}
    \item \textbf{Entrada:} programa $P$ y entrada $I$.
    \item \textbf{Solución:} \palabra{si} el programa $P$ para con entrada $I$, \palabra{no} en otro caso.
\end{itemize}
\end{framed}
\end{column}
\end{columns}

\end{frame}


\begin{frame}[fragile]{El problema de la parada}

\begin{lstlisting}[language=Python]
import utilidades
from maquina_universal import maquina_universal

def maquina_universal_parada(entrada_codificada):
    (programa, entrada) = entrada_codificada.split(SEP)
    salida = maquina_universal(programa, entrada)

    if salida == 'sí':
        return 'para'
    else:
        utilidades.ciclar()
\end{lstlisting}
\vspace*{-4mm}
\begin{center}
{\small Programa 4.9.: \texttt{maquina\_universal\_parada.py}}
\end{center}

\end{frame}

\begin{frame}[label=maquina-universal-parada]{El problema de la parada}

$$\begin{matrix}
\begin{matrix}P = \texttt{programa}\text{ }\text{\underline{devuelve \palabra{sí}}}\\\text{con entrada }\texttt{entrada}\end{matrix}
& \stackrel{(1)}{\longrightarrow} &
\begin{matrix}P' = \texttt{maquina\_universal\_parada}\\\text{\underline{para} con entrada}\\ \texttt{programa + SEP + entrada}\end{matrix} \\
\\
\\
\begin{matrix}P = \texttt{programa}\text{ }\text{\underline{no devuelve \palabra{sí}}}\\\text{con entrada }\texttt{entrada}\end{matrix}
& \stackrel{(2)}{\longrightarrow} &
\begin{matrix}P' = \texttt{maquina\_universal\_parada}\\\text{\underline{no para} con entrada}\\ \texttt{programa + SEP + entrada}\end{matrix}
\end{matrix}$$

\end{frame}



\begin{frame}[fragile]{El problema de la parada}

\begin{lstlisting}[language=Python]
import utilidades
from maquina_universal import maquina_universal

def maquina_universal_parada(entrada_codificada):
    (programa, entrada) = entrada_codificada.split(SEP)
    salida = maquina_universal(programa, entrada)

    (*@\hl{if salida == \textquotesingle{}sí\textquotesingle{}:}  @*)
        (*@\hl{return \textquotesingle{}para\textquotesingle{}}  @*)
    else:
        utilidades.ciclar()
\end{lstlisting}
\vspace*{-4mm}
\begin{center}
{\small Programa 4.9.: \texttt{maquina\_universal\_parada.py}}
\end{center}

\end{frame}

\againframe{maquina-universal-parada}

\begin{frame}[fragile]{El problema de la parada}

\begin{lstlisting}[language=Python]
import utilidades
from maquina_universal import maquina_universal

def maquina_universal_parada(entrada_codificada):
    (programa, entrada) = entrada_codificada.split(SEP)
    salida = maquina_universal(programa, entrada)

    if salida == 'sí':
        return 'para'
    (*@\hl{else:}  @*)
        (*@\hl{utilidades.ciclar()}  @*)
\end{lstlisting}
\vspace*{-4mm}
\begin{center}
{\small Programa 4.9.: \texttt{maquina\_universal\_parada.py}}
\end{center}

\end{frame}

\againframe{maquina-universal-parada}


\begin{frame}{El problema de la parada}

\begin{center}
    \large{$\textsc{Universal}(\texttt{programa}, \texttt{entrada})$}
    
    \large{$|\:|$}
    
    \large{$\textsc{Parada(\texttt{maquina\_universal\_parada}, \texttt{programa + SEP + entrada})}$}
\end{center}

\begin{columns}
\begin{column}{0.5\textwidth}
\begin{framed}
$$\text{\large \textsc{Universal}}$$

\begin{itemize}
    \item \textbf{Entrada:} programa $P$ y entrada $I$.
    \item \textbf{Solución:} \palabra{sí} si $P$ está bien definido y $P(I)=\text{\palabra{sí}}$, \palabra{no} en otro caso.
\end{itemize}
\end{framed}
\end{column}
\begin{column}{0.5\textwidth}  %%<--- here
\begin{framed}
$$\text{\large \textsc{Parada}}$$

\begin{itemize}
    \item \textbf{Entrada:} programa $P$ y entrada $I$.
    \item \textbf{Solución:} \palabra{si} el programa $P$ para con entrada $I$, \palabra{no} en otro caso.
\end{itemize}
\end{framed}
\end{column}
\end{columns}

\end{frame}



\begin{frame}[fragile]{El problema de la parada}

La reducción $\textsc{Universal}\leq_T\textsc{Parada}$ es:

\begin{lstlisting}[language=Python]
import utilidades
from parada import parada # oráculo

def universal_a_parada(programa, entrada):
    entrada_codificada = programa + SEP + entrada
    maquina_universal_parada = utilidades.leer('maquina_universal_parada.py')

    return parada(maquina_universal_parada, entrada_codificada)
\end{lstlisting}
\vspace*{-4mm}
\begin{center}
{\small Programa 4.10.: \texttt{universal\_a\_parada.py}}
\end{center}

\vspace*{-8.5mm}

\hfill\square


\end{frame}


\begin{frame}<presentation:0>{El problema de la parada}

Ya tenemos la reducción: \texttt{universal\_a\_parada.py}, dados como entrada un programa $P=\texttt{programa}$ y una entrada $I=\texttt{entrada}$:

\pause

\begin{itemize}
    \item Si $P(I)$ está bien definido y es $P(I)=\palabra{sí}$, entonces, aplicando (1), devolverá \palabra{sí}.
    \item En caso contrario, si $P(I)$ no está bien definido o no es $P(I)=\palabra{sí}$, entonces, aplicando (2), devolverá \palabra{no}.\vspace{-8mm}
\end{itemize}

\hfill\square

\end{frame}

\begin{frame}<presentation:0>{El problema \textsc{AdivinaConsistente}}

\begin{framed}
$$\text{\large \textsc{AdivinaConsistente}}$$

\begin{itemize}
    \item \textbf{Entrada:} un programa $P$.
    \item \textbf{Solución:} $\palabra{sí}$ si $P$ devuelve \palabra{sí} con entrada vacía (acepta), \palabra{no} en caso contrario (rechaza o cicla).
\end{itemize}
\end{framed}
\begin{center}
{\small Problema 4.11.: \textsc{AdivinaConsistente}}
\end{center}

\vspace{5mm}

\textbf{Proposición 4.8.} \textsc{AdivinaConsistente} \emph{es no decidible.}

\end{frame}

\begin{frame}{Resumen}

\begin{table}[H]
\hspace*{-2.5mm}
\centering
\begin{tabular}{@{}llll@{}}
\cmidrule[0.7pt]{2-4}
 & Entrada & Salida & Decidible$\;\;$ \\ \midrule
\textsc{MásAQueB} & $u\in\{\palabra{a}, \palabra{b}\}^*$ & -- \palabra{sí}, si $\#(u,a)>\#(u,b)$ & sí \\
 && -- \palabra{no}, en caso contrario & \\[4pt]
\textsc{Universal} & $P$, $I$ & -- \palabra{sí}, si $P(I)=\palabra{sí}$ & no \\
 & & -- \palabra{no}, en caso contrario &  \\[4pt]
\textsc{Parada} & $P$, $I$ & -- \palabra{sí}, si $P$ para con entrada $I$ & no \\
& & -- \palabra{no}, en caso contrario& \\\bottomrule
\end{tabular}
\end{table}

\end{frame}

\section{Sistemas lógicos}

\begin{frame}{Sistemas formales}

\textbf{Definición 5.1} (Sistema formal)\textbf{.}
{\Large $$\mathcal{S} = (\mathcal{L}, \mathcal{A}, \mathcal{R})$$}
\vspace*{-5mm}
\begin{itemize}
    \item $\mathcal{L}$: \hle{lenguaje formal} sobre alfabeto $A$. Sus elementos son \hle{fórmulas} (\emph{bien formadas}).
    \item $\mathcal{A}$: \hle{axiomas}, fórmulas de $\mathcal{L}$.
    \item $\mathcal{R}$: \hle{reglas de inferencia},\footnote{Todas deben ser computables.} \emph{premisas} \hl{$\mapsto$} \emph{derivaciones}.
\end{itemize}

\end{frame}

\begin{frame}{Sistemas formales}

\textbf{Sistema formal 5.1} (\textbf{SumaBinaria})\textbf{.}

$\textbf{SumaBinaria}=(\mathcal{L}, \mathcal{A}, \mathcal{R})$, con:

\begin{itemize}
    \item Lenguaje $\mathcal{L}$: reconocido por expresión regular sobre $\{\texttt{0}, \texttt{1}, \texttt{+}, \texttt{=}\}^*$:
    $$
        N(\texttt{+}N)^*=N(\texttt{+}N)^*,\;\;\text{con}\;\;N\text{ regex. }\texttt{1}(\texttt{0}|\texttt{1})^*
    $$
    \item Conjunto de axiomas $\mathcal{A}$:
    $$
        \mathcal{A} = \{\texttt{1=1}\}
    $$
    \item Conjunto de reglas de inferencia $\mathcal{R}$:
    $$
    \begin{matrix}
    \text{(R1)} & A\texttt{=}B & \mapsto & \texttt{1+}A\texttt{=}\texttt{1+}B \\
    \text{(R2)} & A\underline{N}\texttt{+}\underline{N}B & \mapsto & AN\texttt{0}B \\
    \text{(R3)} & A\texttt{\underline{1}+}\underline{N\texttt{0}}B & \mapsto & AN\texttt{1}B
    \end{matrix}
    $$
\end{itemize}

\end{frame}

\begin{frame}{Sistemas formales}

\textbf{Definición 5.2} (Teorema y demostración)\textbf{.}
{\large $$\phi\text{ \hle{teorema}} \;\;\; \stackrel{def.}{\iff} \;\;\; \psi \xmapsto{*} \phi,\text{ con }\psi\text{ axioma}$$}
\vspace*{-5mm}
\begin{itemize}
    \item \hle{Demostración:} secuencia de derivaciones.
\end{itemize}
\end{frame}

\begin{frame}{Sistemas formales}

\textbf{Ejemplo 5.1.}

\begin{center}
    \texttt{1+1+1=11} es un teorema de \textbf{SumaBinaria}
\end{center}

\begin{table}[h]
\centering
\begin{tabular}{lrcl}
& (F1) & \texttt{1=1} & (axioma) \\
$\mapsto$ & (F2) & \texttt{1+1=1+1} & (aplicar R1 a F1) \\
$\mapsto$ & (F3) & \texttt{1+1=10} & (aplicar R2 a F2) \\
$\mapsto$ & (F4) & \texttt{1+1+1=1+10} & (aplicar R1 a F3) \\
$\mapsto$ & (F5) & \texttt{1+1+1=11} & (aplicar R3 a F4) \\
\end{tabular}
\end{table}

\end{frame}

\begin{frame}{Sistemas lógicos}


\textbf{Definición 5.3} (Sistema lógico)\textbf{.}
{\Large $$\mathcal{G} = (\mathcal{S}, \mathcal{I})$$}
\vspace*{-5mm}
\begin{itemize}
    \item $\mathcal{S}$: sistema formal.
    \item $\mathcal{I}$: \hle{semántica} o \emph{asignación de verdad},\footnote{No necesariamente computable, y no necesariamente definida para todas las fórmulas.} $\mathcal{I}:\mathcal{L}\longrightarrow\{\text{verdadero}, \text{falso}\}$.
\end{itemize}

\vspace{5mm}

\textbf{Definición 5.4} (Fórmula verdadera y fórmula falsa)\textbf{.}

\begin{itemize}
    \item $\mathcal{I}(\phi)=\text{verdadero}$ $\;\stackrel{def.}{\iff}\;$ $\phi$ es \hle{verdadera}.
    \item $\mathcal{I}(\phi)=\text{falso}$ $\;\stackrel{def.}{\iff}\;$ $\phi$ es \hle{falsa}.
\end{itemize}

\end{frame}


\begin{frame}{Sistemas lógicos}

\textbf{Sistema lógico 5.2} (\textbf{SumaBinariaLógico})\textbf{.}

Definimos \textbf{SumaBinariaLógico} con:

\begin{itemize}
    \item Sintaxis: la de \textbf{SumaBinaria}.
    \item Semántica: $\phi$ verdadero $\iff$ lo es visto como una fórmula de aritmética binaria.
\end{itemize}

\end{frame}


\begin{frame}{Sistemas lógicos}

\textbf{Ejemplo 5.2.}

\begin{center}
    \texttt{1+1+1=11} es una fórmula verdadera de \textbf{SumaBinariaLógico}
\end{center}

En efecto, la siguiente operación aritmética en binario es cierta:
$$1+1+1=11$$

\end{frame}

\begin{frame}{Sistemas lógicos}

\begin{columns}
\begin{column}{0.35\textwidth}
\begin{center}
    \Large
    \textbf{sistema formal}

    sintaxis
\end{center}
\end{column}
\begin{column}{0.35\textwidth}  %%<--- here
\begin{center}
    \Large
    \textbf{sistema lógico}

    sintaxis + semántica
\end{center}
\end{column}
\end{columns}

\end{frame}

\begin{frame}{Solidez, completitud y decidibilidad}

\textbf{Definición 5.5} (Solidez)\textbf{.}
{\large $$\mathcal{G}\text{ \hle{sólido}} \;\;\; \stackrel{def.}{\iff} \;\;\; \underset{\forall \phi \text{ fórmula de }\mathcal{G}}{\phi \text{ teorema} \Rightarrow \phi\text{ verdadero}}$$}

\textbf{Definición 5.6} (Completitud)\textbf{.}
{\large $$\mathcal{G}\text{ \hle{(sintácticamente) completo}} \;\;\; \stackrel{def.}{\iff} \;\;\; \underset{\forall \phi \text{ fórmula de }\mathcal{G}}{\phi \text{ verdadero} \Rightarrow \phi\text{ teorema}}$$}

\end{frame}

\begin{frame}{Solidez, completitud y decidibilidad}

\textbf{Proposición 4.8.} \emph{El sistema lógico} \textbf{SumaBinariaLógico} \emph{es sólido y completo.}
\pause
{\large $$\text{verdadero}\iff\text{teorema}$$}

\end{frame}

\begin{frame}{Solidez, completitud y decidibilidad}

\begin{columns}
\begin{column}{0.5\textwidth}
\textbf{Sistema lógico 5.3}

(\textbf{SumaBinariaLógicoRoto})\textbf{.}\\[3mm]

Añadir a \textbf{SumaBinarioLógico} la regla (R1b).

$$
    \begin{matrix}
    \text{(R1)} & A\texttt{=}B & \mapsto & \texttt{1+}A\texttt{=}\texttt{1+}B \\
    \text{(R1b)} & A\texttt{=}B & \mapsto & \texttt{1+}A\texttt{=}B \\
    \text{(R2)} & A\underline{N}\texttt{+}\underline{N}B & \mapsto & AN\texttt{0}B \\
    \text{(R3)} & A\texttt{\underline{1}+}\underline{N\texttt{0}}B & \mapsto & AN\texttt{1}B
    \end{matrix}
$$

\vspace{3mm}

\textbf{Proposición 5.2.} \emph{El sistema lógico} \textbf{SumaBinariaLógico} \emph{es completo y no sólido.}
{\large $$\text{verdadero}\Rightarrow\text{teorema}$$}

\end{column}
\begin{column}{0.5\textwidth}  %%<--- here
\textbf{Sistema lógico 5.4}

(\textbf{SumaBinariaLógicoRestringido})\textbf{.}\\[3mm]

Eliminar de \textbf{SumaBinarioLógico} la regla (R3).

$$
    \begin{matrix}
    \text{(R1)} & A\texttt{=}B & \mapsto & \texttt{1+}A\texttt{=}\texttt{1+}B \\
     \\
    \text{(R2)} & A\underline{N}\texttt{+}\underline{N}B & \mapsto & AN\texttt{0}B \\
    \\
    \end{matrix}
$$

\vspace{3mm}

\textbf{Proposición 5.3.} \emph{El sistema lógico} \textbf{Suma-} \textbf{BinariaRestringido} \emph{es incompleto y sólido.}
{\large $$\text{verdadero}\Leftarrow\text{teorema}$$}
\end{column}
\end{columns}

\end{frame}


\begin{frame}{Solidez, completitud y decidibilidad}

\begin{columns}
\begin{column}{0.5\textwidth}
\begin{center}
    $\mathcal{G} \text{ sistema lógico}$
\end{center}

\begin{framed}
$$\text{\textsc{EsVerdadero}}_\mathcal{G}$$

\begin{itemize}
    \item \textbf{Entrada:} una fórmula $\phi$ de $\mathcal{G}$.
    \item \textbf{Solución:} \palabra{sí} si $\mathcal{I}(\phi)$ está bien definido y $\mathcal{I}(\phi)=\text{verdadero}$, y \palabra{no} en caso contrario.
\end{itemize}
\end{framed}
\begin{center}
{\small Problema 5.1.: $\text{\textsc{EsVerdadero}}_\mathcal{G}$}
\end{center}
\end{column}
\begin{column}{0.5\textwidth}
\begin{center}
    $\mathcal{S} \text{ sistema formal}$
\end{center}

\begin{framed}
$$\text{\textsc{EsTeorema}}_\mathcal{S}$$

\begin{itemize}
    \item \textbf{Entrada:} una fórmula $\phi$ de $\mathcal{S}$.
    \item \textbf{Solución:} \palabra{sí} si $\phi$ es un teorema de $\mathcal{S}$, y \palabra{no} en caso contrario.
\end{itemize}
\end{framed}
\begin{center}
{\small Problema 5.2.: $\text{\textsc{EsTeorema}}_\mathcal{S}$}
\end{center}
\end{column}
\end{columns}

\vspace{5mm}
\textbf{Definición 5.7} (Decidibilidad de un sistema lógico)\textbf{.}

Decimos que un sistema lógico $\mathcal{G}$ es \hle{decidible} si $\text{\textsc{EsVerdadero}}_\mathcal{G}$ es decidible.
\vspace{10mm}
\end{frame}


\begin{frame}{El sueño de Hilbert}

\textbf{Segundo problema de Hilbert.}

¿Son completas las matemáticas? \\ Esto es, ¿puede probarse o no cada sentencia matemática?

\vspace{5mm}

\pause

Spoiler: \textbf{no} (gracias a Gödel).

\end{frame}


\begin{frame}{La aritmética de Peano}

\begin{table}[H]
\centering
\begin{tabular}{@{}ccl@{}}
\toprule
Símbolo & Codificación & Descripción \\ \midrule
$\neg(p)$ & \texttt{NO (p)} & Negación \\
$(p) \vee (q)$ & \texttt{(p) Ó (q)} & Disyunción \\
$\forall x : p$ & \texttt{PARA TODO (x) : (p)} & Cuantificador universal \\
$S(a)$ & \texttt{S(a)} & Sucesor \\
$\exists x : p$ & \texttt{EXISTE (x) : (p)} & Cuantificador existencial \\
$(p) \Rightarrow (q)$ & \texttt{(p)}\texttt{ }\texttt{=>}\texttt{ }\texttt{(q)} & Implicación \\
$(p) \iff (q)$ & \texttt{(p)}\texttt{ }\texttt{<=>}\texttt{ }\texttt{(q)} & Coimplicación \\
$(p)=(q)$ & \texttt{(p) = (q)} & Igualdad \\\bottomrule
\end{tabular}
\end{table}
\begin{center}
{\small Tabla 5.2.: Una forma de traducir los símbolos de la aritmética de Peano a \emph{strings}}
\end{center}

\end{frame}

\begin{frame}{La aritmética de Peano}

\textbf{Podemos traducir fórmulas de la aritmética de Peano a \emph{strings}.}

\begin{itemize}
    \item[(PF2)] Para todo entero $x$, es $x^2+1>0$.

    \texttt{PARA TODO x : x{\^{}}2 + 1 > 0}

    \item[(PF6)] Si $n>3$, no existen $x,y,z$ enteros tal que $x^n+y^n=z^n$ (el \emph{último teorema de Fermat}).

    \texttt{n > 3 Y NO EXISTE x, y, z : x{\^{}}n + y{\^{}}n = z{\^{}}n}
\end{itemize}

\end{frame}


\begin{frame}<2>{La aritmética de Peano}

\begin{columns}
\begin{column}{0.5\textwidth}
\pause[2]{
\begin{framed}
$$\text{\textsc{EsVerdaderoPeano}}$$

\begin{itemize}
    \item \textbf{Entrada:} una fórmula $\phi$ de \textbf{Peano}.
    \item \textbf{Solución:} \palabra{sí} si $\phi$ es una fórmula bien formada, cerrada y verdadera en \textbf{Peano}; \palabra{no} en caso contrario.
\end{itemize}
\end{framed}
\begin{center}
{\small Problema 5.4.: $\text{\textsc{EsVerdaderoPeano}}$}
\end{center}
}
\end{column}
\begin{column}{0.5\textwidth}
\begin{framed}
$$\text{\textsc{EsTeoremaPeano}}$$
\vspace{-5mm}
\begin{itemize}
    \pause[1]{
    \item \textbf{Entrada:} una fórmula $\phi$ de \textbf{Peano}.
    \item \textbf{Solución:} \palabra{sí} si $\phi$ es una fórmula bien formada y cerrada en \textbf{Peano}, y se trata de un teorema de \textbf{Peano}; \palabra{no} en caso contrario.
    }
\end{itemize}

\end{framed}
\begin{center}
{\small Problema 5.3.: $\text{\textsc{EsTeoremaPeano}}$}
\end{center}
\end{column}
\end{columns}

\end{frame}

\section{Incompletitud}


\begin{frame}{Incompletitud de la aritmética de Peano}

\textbf{Afirmación.} \textbf{Peano} \emph{es sólido.}

\pause

\begin{itemize}
    \item Es razonable pensar que todos los teoremas son verdaderos.
    \item Es algo demostrable (fuera del trabajo).
\end{itemize}

\end{frame}



\begin{frame}<1>{Incompletitud de la aritmética de Peano}

\begin{columns}
\begin{column}{0.5\textwidth}
\pause[2]{
\begin{framed}
$$\text{\textsc{EsVerdaderoPeano}}$$

\begin{itemize}
    \item \textbf{Entrada:} una fórmula $\phi$ de \textbf{Peano}.
    \item \textbf{Solución:} \palabra{sí} si $\phi$ es una fórmula bien formada, cerrada y verdadera en \textbf{Peano}; \palabra{no} en caso contrario.
\end{itemize}
\end{framed}
\begin{center}
{\small Problema 5.4.: $\text{\textsc{EsVerdaderoPeano}}$}
\end{center}
}
\end{column}
\begin{column}{0.5\textwidth}
\begin{framed}
$$\text{\textsc{EsTeoremaPeano}}$$
\vspace{-5mm}
\pause[1]{
\begin{itemize}
    \item \textbf{Entrada:} una fórmula $\phi$ de \textbf{Peano}.
    \item \textbf{Solución:} \palabra{sí} si $\phi$ es una fórmula bien formada y cerrada en \textbf{Peano}, y se trata de un teorema de \textbf{Peano}; \palabra{no} en caso contrario.
\end{itemize}
}
\end{framed}
\begin{center}
{\small Problema 5.3.: $\text{\textsc{EsTeoremaPeano}}$}
\end{center}
\end{column}
\end{columns}

\end{frame}



\begin{frame}{Incompletitud de la aritmética de Peano}

\begin{enumerate}
    \item Encontrar una forma de traducir \textsc{Parada} a \textbf{Peano}.
    \item Ver cómo demostrar automáticamente teoremas de \textbf{Peano}.
    \item Encontrar una fórmula verdadera de \textbf{Peano} no demostrable.
\end{enumerate}

\end{frame}



\begin{frame}{Incompletitud de la aritmética de Peano}

\begin{center}
    \textbf{Paso 1.} Traducir \textbf{Parada} a \textsc{Peano}

    \line(1,0){50}
    %\rule{20mm}
\end{center}

\vspace{5mm}

\textbf{Proposición 5.5.} {\itshape
Sea $P$ un programa de ordenador (o una máquina de Turing). La afirmación $H$ dada por
\begin{center}
    ``$P$ para con entrada vacía''
\end{center}
puede ser traducida en una fórmula $\phi_H$ equivalente en \normalfont{\textbf{Peano}}.} {\itshape (En este caso, con ``equivalente'' nos referimos a que $H$ es verdadera si y solo si $\phi_H$ es verdadera.)}

\vspace{5mm}

\pause

\textbf{\emph{Demostración.}}
$\phi_H=\;$``Existe un natural $n$ de modo que $Parado(Paso^n(c_0))=1$''\hfill\square

\end{frame}


\begin{frame}{Incompletitud de la aritmética de Peano}

\textbf{Definición 6.1} (Sistema de tipo Turing)\textbf{.}

Un sistema (lógico o formal) es de \hle{tipo Turing} si puede expresar, dado un programa $P$, que
\begin{center}
    ``$P$ para con entrada vacía''
\end{center}

\end{frame}


\begin{frame}<presentation:0>{Incompletitud de la aritmética de Peano}

\textbf{\emph{Demostración.}} Usamos máquinas de Turing.

\begin{itemize}
    \pause
    \item Podemos codificar una configuración dada mediante \emph{strings}:
    $$
    q_R\;:\;X\:\fbox{$X$}\:X\:a \;\;\;\; \longrightarrow \;\;\;\; \texttt{q\_R : X \^{}X X a}
    $$
    \vspace*{-5mm}
    \pause
    \item A su vez los codificamos mediante números binarios en \emph{strings} (palabras del alfabeto $\{\texttt{0}, \texttt{1}\}$).
\end{itemize}



\end{frame}

\begin{frame}<presentation:0>{Incompletitud de la aritmética de Peano}

\begin{itemize}
    \item Función de transición: $Paso$ (de enteros a enteros).
    \pause
    \item Configuración inicial: $c_0$ (entero).
    \pause
    \item Proceso de cálculo: $n$ pasos la configuración será $Paso^n(c_0)$
    $$
    c_0 \rightarrow Paso(c_0) \rightarrow Paso(Paso(c_0)) = Paso^2(c_0) \rightarrow ...
    $$
    \pause
    \vspace*{-5mm}
    \item Condición de parada: $Parado$
    $$
    Parado(m) =
    \begin{cases} 
      0, & \text{si }m\text{ no es configuración de parada} \\
      1, & \text{si }m\text{ es configuración de parada}
   \end{cases}
    $$
\end{itemize}

\end{frame}

\begin{frame}<presentation:0>{Incompletitud de la aritmética de Peano}

\begin{center}
    ``$P$ para con entrada vacía''

    \pause
    $\updownarrow$

    ``Existe un natural $n$ de modo que $Paso^n(c_0)$ es una configuración de parada''

    \pause
    $\updownarrow$

    ``Existe un natural $n$ de modo que $Parado(Paso^n(c_0))=1$''
\end{center}

\pause
La última afirmación puede traducirse a una fórmula de \textbf{Peano} $\phi_H$.\hfill\square

\end{frame}

\begin{frame}{Incompletitud de la aritmética de Peano}

\begin{center}
    \textbf{Paso 2.} Demostrar automáticamente teoremas de \textsc{Peano}

    \line(1,0){50}
    %\rule{20mm}
\end{center}

\vspace{5mm}

\textbf{Proposición 5.8.} \textsc{EsTeoremaPeano} \emph{es semidecidible.}

\pause
\begin{itemize}
    \item Si la entrada es un teorema, para y devuelve \palabra{sí}.
    \item Si no, cicla.
\end{itemize}

\end{frame}


\begin{frame}[fragile]{Incompletitud de la aritmética de Peano}

\textbf{Demostración.}

\begin{lstlisting}[language=Python]
import utilidades
from es_prueba_peano import es_prueba_peano # NO oráculo

def es_teorema_peano(fórmula):
    demostración = ''

    while True:
        if es_prueba_peano(demostración, fórmula) == 'sí':
            return 'sí'
        demostración = utilidades.siguiente_string(demostración)
\end{lstlisting}
\vspace*{-4mm}
\begin{center}
{\small Programa 5.2.: \texttt{es\_teorema\_peano.py}}
\end{center}

\vspace*{-9.5mm}

\hfill\square

\end{frame}


\begin{frame}{Incompletitud de la aritmética de Peano}

\textbf{Definición 5.8} (Recursivamente axiomatizable)\textbf{.}

Un sistema formal $\mathcal{S}$ es \hle{recursivamente axiomatizable} si $\textsc{EsTeorema}_\mathcal{S}$ es semidecidible.
\end{frame}

\begin{frame}[fragile]{Incompletitud de la aritmética de Peano}

\begin{center}
    \textbf{Paso 3.} Encontrar una fórmula verdadera y no demostrable de \textsc{Peano}

    \line(1,0){50}
    %\rule{20mm}
\end{center}

\end{frame}


\begin{frame}[fragile]{Incompletitud de la aritmética de Peano}

\begin{lstlisting}[language=Python]
import utilidades

from parada_a_peano import parada_a_peano     # paso 1
from es_teorema_peano import es_teorema_peano # paso 2

def godel_peano(entrada):
    programa_godel = utilidades.leer('godel_peano.py')
    para_en_peano = parada_a_peano(programa_godel)
    no_para_en_peano = 'NO (' + para_en_peano + ')'

    if es_teorema_peano(no_para_en_peano) == 'sí':
        return 'para'
    else:
        utilidades.ciclar()
\end{lstlisting}
\vspace*{-4mm}
\begin{center}
{\small Programa 6.1.: \texttt{godel\_peano.py}}
\end{center}

\end{frame}


\begin{frame}<1-3>[label=godel-peano]{Incompletitud de la aritmética de Peano}

\textbf{Proposición 6.1.} \texttt{godel\_peano.py} \emph{no para con entrada vacía.}

\pause
\vspace{5mm}

\textbf{Demostración.} Asumimos el contrario, y entramos en contradicción.

\begin{center}
\pause
\texttt{godel\_peano.py} para con entrada vacía
\vspace*{-2mm}
\pause
$$\downarrow$$

\vspace*{-2mm}
\texttt{es\_teorema\_peano(no\_para\_en\_peano) == \palabra{sí}}
\vspace*{-2mm}
\pause
$$\downarrow$$

\vspace*{-2mm}
``\texttt{godel\_peano.py} no para con entrada vacía'' es un teorema
\vspace*{-2mm}
\pause
$$\begin{matrix}
    \hspace{22mm} & \downarrow & \text{\small \textbf{Peano} es sólido}
\end{matrix}$$

\vspace*{-2mm}
\texttt{godel\_peano.py} no para con entrada vacía \hspace{2mm} {\large\lightning}
\end{center}

\vspace*{-9.5mm}

\pause[6]{\hfill\square}

\end{frame}

\begin{frame}[fragile]{Incompletitud de la aritmética de Peano}

\begin{lstlisting}[language=Python]
import utilidades

from parada_a_peano import parada_a_peano     # paso 1
from es_teorema_peano import es_teorema_peano # paso 2

def godel_peano(entrada):
    programa_godel = utilidades.leer('godel_peano.py')
    para_en_peano = parada_a_peano(programa_godel)
    no_para_en_peano = 'NO (' + para_en_peano + ')'

    (*@\hl{if es\_teorema\_peano(no\_para\_en\_peano) == \textquotesingle{}sí\textquotesingle{}:}  @*)
        (*@\hl{return \textquotesingle{}para\textquotesingle{}}  @*)
    else:
        utilidades.ciclar()
\end{lstlisting}
\vspace*{-4mm}
\begin{center}
{\small Programa 6.1.: \texttt{godel\_peano.py}}
\end{center}

\end{frame}

\againframe<4->{godel-peano}


\begin{frame}<1-5>[label=incompletitud-peano]{Incompletitud de la aritmética de Peano}

\textbf{Teorema 6.1.} \emph{La aritmética de Peano} (\textbf{Peano}) \emph{es incompleta.}

\vspace{5mm}

\pause

\textbf{Demostración.} \texttt{no\_para\_en\_peano} es verdadero y no es un teorema.

\begin{itemize}

    \pause
    \item Es verdadero: por la proposición anterior.

    \pause
    \item No es un teorema:

    \begin{center}
    \pause
\texttt{godel\_peano.py} no para con entrada vacía
\vspace*{-2mm}
\pause
$$\downarrow$$

\vspace*{-2mm}
no puede ser \texttt{es\_teorema\_peano(no\_para\_en\_peano) == \palabra{sí}}
\vspace*{-2mm}
\pause
$$\downarrow$$

\vspace*{-2mm}
\texttt{no\_para\_en\_peano} no es un teorema
\end{center}

\vspace*{-5.5mm}

\pause[7]{\hfill\square}

\end{itemize}

\end{frame}

\begin{frame}[fragile]{Incompletitud de la aritmética de Peano}

\begin{lstlisting}[language=Python]
import utilidades

from parada_a_peano import parada_a_peano     # paso 1
from es_teorema_peano import es_teorema_peano # paso 2

def godel_peano(entrada):
    programa_godel = utilidades.leer('godel_peano.py')
    para_en_peano = parada_a_peano(programa_godel)
    no_para_en_peano = 'NO (' + para_en_peano + ')'

    if es_teorema_peano(no_para_en_peano) == 'sí':
        return 'para'
    (*@\hl{else:}  @*)
        (*@\hl{utilidades.ciclar()}  @*)
\end{lstlisting}
\vspace*{-4mm}
\begin{center}
{\small Programa 6.1.: \texttt{godel\_peano.py}}
\end{center}

\end{frame}


\againframe<6->{incompletitud-peano}



\begin{frame}{Incompletitud de la aritmética de Peano}

{\Large
$$\begin{matrix}
    \begin{matrix}
        \textbf{Peano} \\
        \text{rec. axiom.} \\
        \text{tipo Turing} \\
        \text{\ul{sólido}} \\
        $\;$
    \end{matrix}
    & \;\;\;\;\;\;\Longrightarrow\;\;\;\;\;\;
    \begin{matrix}
        \textbf{Peano} \\
        \text{\ul{incompleto}} \\
        $\;$
    \end{matrix}
\end{matrix}$$
}

\end{frame}


\begin{frame}{Versión semántica del Primer Teorema de Incompletitud}

{\Large
$$\begin{matrix}
    \begin{matrix}
        \text{rec. axiom.} \\
        \text{tipo Turing} \\
        \text{\ul{sólido}} \\
    \end{matrix}
    & \;\;\;\;\;\;\Longrightarrow\;\;\;\;\;\;
    \begin{matrix}
        \text{\ul{incompleto}} \\
    \end{matrix}
\end{matrix}$$
}

\end{frame}



\begin{frame}{Versión semántica del Primer Teorema de Incompletitud}

\textbf{Teorema 6.2.} {\itshape
Sea $\mathcal{G}$ un sistema lógico que verifique las siguientes condiciones:
\begin{itemize}
    \item[(1)] $\mathcal{G}$ es sólido
    
    (todo teorema es verdadero).
    \item[(2)] $\mathcal{G}$ es recursivamente axiomatizable.
    \item[(3)] $\mathcal{G}$ es de tipo Turing.
\end{itemize}
Entonces, \hl{$\mathcal{G}$ es (semánticamente) incompleto}\\
(existen fórmulas verdaderas de $\mathcal{G}$ que no son teoremas).
}

\end{frame}


\begin{frame}<presentation:0>{Versión sintáctica del Primer Teorema de Incompletitud}

\textbf{Proposición 6.2.} {\itshape
Sea $\mathcal{S}$ un sistema formal recursivamente axiomatizable y sintácticamente completo. Entonces, $\textsc{EsTeorema}_\mathcal{S}$ es decidible.
}

\end{frame}


\begin{frame}{Versión sintáctica del Primer Teorema de Incompletitud}

\textbf{Teorema 6.3.} {\itshape
Sea $\mathcal{S}$ un sistema formal que verifique las siguientes condiciones:
\begin{itemize}
    \item[(1)] $\mathcal{S}$ es sintácticamente consistente.

    (no es posible que $\phi$ y $\neg\phi$ sean teoremas).
    \item[(2)] $\mathcal{S}$ es recursivamente axiomatizable.
    \item[(2)] $\mathcal{S}$ es de tipo Turing.
\end{itemize}
Entonces, \hl{$\mathcal{S}$ es sintácticamente incompleto} \\
(existen fórmulas $\phi$ de $\mathcal{S}$ tal que ni $\phi$ ni $\neg\phi$ son teoremas).
}

\end{frame}


\begin{frame}<presentation:0>{Versión sintáctica del Primer Teorema de Incompletitud}

\textbf{Demostración.} Si $\mathcal{S}$ fuese sintácticamente completo, \textsc{AdivinaConsistente} sería decidible.

\vspace{5mm}

\begin{framed}
$$\text{\large \textsc{AdivinaConsistente}}$$

\begin{itemize}
    \item \textbf{Entrada:} un programa $P$.
    \item \textbf{Solución:} $\palabra{sí}$ si $P$ devuelve \palabra{sí} con entrada vacía (acepta), \palabra{no} en caso contrario (rechaza o cicla).
\end{itemize}
\end{framed}

\end{frame}


\begin{frame}<presentation:0>[fragile]{Versión sintáctica del Primer Teorema de Incompletitud}

\begin{lstlisting}[language=Python]
import utilidades
from parada_a_S import parada_a_S # NO es un oráculo|\label{line:adivina-consistente-import-1}|
from es_teorema_S import es_teorema_S # NO es un oráculo|\label{line:adivina-consistente-import-2}|

def adivina_consistente(programa):|\label{line:adivina-consistente-main}|
    programa_para = parada_a_S(programa)|\label{line:adivina-consistente-programa-para}|
    es_teorema = es_teorema_S(programa_para)|\label{line:adivina-consistente-es-teorema}|
    if es_teorema == 'sí':|\label{line:adivina-consistente-si}|
        salida = maquina_universal(programa, '')|\label{line:adivina-consistente-universal}|
        return salida  # es 'sí' o 'no'
    else:|\label{line:adivina-consistente-no}|
        return 'no'
\end{lstlisting}
\vspace*{-4mm}
\begin{center}
{\small Programa 6.4.: \texttt{adivina\_consistente.py}}
\end{center}

\vspace*{-9.5mm}

\hfill\square

\end{frame}


\begin{frame}{Primer Teorema de Incompletitud}

$$
\begin{array}{ccc}
    {
        \begin{matrix}
            \text{\textbf{versión semántica}} \\
            \small{\;}
        \end{matrix}
    }
    & \;\;\;\;\;\; &
    {
        \begin{matrix}
            \text{\textbf{versión sintáctica}} \\
            \small{\;}
        \end{matrix}
    }
    \\
    {
        \begin{matrix}
            \mathcal{G}\text{ sistema lógico}\\
            \text{rec. axiom., de tipo Turing}\\
            \text{y sólido} \\
            \small{\phi \text{ teorema} \Rightarrow \phi\text{ verdadero}}\\
            \tiny{\;}
        \end{matrix}
    }
    &  & 
    {
        \begin{matrix}
            \mathcal{S}\text{ sistema formal}\\
            \text{rec. axiom., de tipo Turing}\\
            \text{y sintácticamente consistente} \\
            \small{\nexists \phi : \phi\text{ y }\neg\phi \text{ son teoremas}} \\
            \tiny{\;}
        \end{matrix}
    } \\
    \Huge{\Downarrow} & & \Huge{\Downarrow} \\
    {
        \begin{matrix}
            \tiny{\;}\\
            \mathcal{G}\text{ no (semánticamente) completo}\\
            \small{\exists\phi:\phi\text{ es verdadero y no es teorema}}\\[2mm]
        \end{matrix}
    }
    & &
    {
        \begin{matrix}
            \tiny{\;}\\
            \mathcal{S}\text{ no sintácticamente completo}\\
            \small{\exists\phi:\text{ni }\phi\text{ ni }\neg\phi\text{ son teoremas}}\\[2mm]
        \end{matrix}
    }
    \\
\end{array}
$$
\end{frame}


\begin{frame}{Primer Teorema de Incompletitud}

\textbf{Teorema 6.4} (Primer Teorema de Incompletitud de Gödel-Rosser)\textbf{.}

Sea $\mathcal{S}$ un sistema formal que verifique las siguientes condiciones:
\begin{itemize}
    \item[(1)] $\mathcal{S}$ es sintácticamente consistente.
    \item[(2)] $\mathcal{S}$ es recursivamente axiomatizable.
    \item[(3)] $\mathcal{S}$ contiene una cierta cantidad de aritmética elemental.
\end{itemize}

Entonces, $\mathcal{S}$ es \hl{sintácticamente incompleto}, es decir, hay fórmulas $\phi$ de $\mathcal{S}$ de modo que ni $\phi$ ni $\neg\phi$ son teoremas.
\end{frame}


\begin{frame}{Consecuencias}

\begin{itemize}
    \item Respuesta negativa al segundo problema de Hilbert.

    ¡Y más! Hilbert pensaba en aritmética...
    \vspace{5mm}
    \pause
    \item La mente humana no es algorítmica (Roger Penrose, John Lucas).
    \item Los seres humanos pensamos de lo general a lo concreto, al contrario de las matemáticas, que lo hacen de lo concreto a lo general (Douglas Hofstadter).
\end{itemize}

\end{frame}

\section{Conclusiones}

\begin{frame}{Conclusiones}

\begin{itemize}
    \item Logramos probar Primer Teorema de Incompletitud mediante perspectiva computacional.
    \item Usamos programas en Python como representación matemática de algoritmo.
    \item Más allá de Gödel: resultado semántico.
    \item Precisamos los conceptos de verdad y demostración.
\end{itemize}

\end{frame}

\maketitle


\end{document}
